\chapter{Forces and free-body diagrams}
\label{vol03:ch01}

\epigraph{To every action there is always opposed an equal reaction: or
the mutual actions of two bodies upon each other are always equal, and
directed to contrary parts.}{Isaac Newton, \emph{Philosophiae Naturalis
Principia Mathematica} (1687), Axioms, or the Laws of Motion, Law III,
in the Cohen and Whitman translation \cite[p.~417]{hist:newton1687}.}

\chapmeta{Half-life: durable. Archetypes: balance, failure (load-path
error). Exercise target: 35. Project track: physical observation plus
analysis (hybrid). Hazard class: none. Reader path default: standard,
with sections explicitly tagged. Named cases: Hyatt Regency walkway
collapse (1981).}

\noindent
We begin Volume III with forces because every mechanical fact in the
volume is, in the end, a force acting on a body. A free-body diagram is
the engineer's instrument for turning a real situation, with its
clutter of contacts and supports and weights, into a set of vectors
that can be summed. The diagram is an accounting discipline: every force on the body
appears on it, and only forces do.

\section{Force as a vector}\pathtag{core}

A \engterm{force} is a vector quantity: it has magnitude, direction,
and a point of application. The magnitude is measured in newtons,
$1\,\si{\newton} = 1\,\si{\kilogram\meter\per\second\squared}$. The
direction is given by a unit vector. The point of application is the
location at which the force acts on the body. Two forces with the same
magnitude and direction but different points of application produce the
same translational acceleration of the centre of mass but different
moments about an arbitrary axis. The point of application is part of
the data.

Two forces $\vec{F}_1$ and $\vec{F}_2$ acting at the same point on a
rigid body have a resultant
\[
\vec{R} = \vec{F}_1 + \vec{F}_2,
\]
constructed geometrically by the parallelogram rule or, equivalently,
by the polygon rule with the tail of $\vec{F}_2$ placed at the head of
$\vec{F}_1$. The parallelogram rule was given a clean experimental
demonstration by Simon Stevin in 1586 with his \emph{Wisconstighe
Ghedaechtenissen}; the rule's algebraic generalisation is the
component decomposition
\[
\vec{F} = F_x \hat{\imath} + F_y \hat{\jmath} + F_z \hat{k},
\]
in any chosen Cartesian basis. The choice of basis is the engineer's;
the physics is invariant.

In practice the engineer chooses a basis aligned with the geometry of
the problem: one axis along the gravitational vertical, one along a
known constraint surface, one perpendicular to both. The axes are
chosen so that the largest number of forces lie along axes rather than
at oblique angles. A poorly chosen basis is the most common source of
arithmetic error in an undergraduate FBD; the cure is not better
arithmetic but a better choice of axes.

How much of the point of application matters depends on what the body is, and
mechanics sorts vectors into three kinds by that dependence. A force on a
\emph{rigid} body is a \engterm{sliding vector}: it may be slid anywhere along
its own line of action without changing its external effect, because sliding
adds a displacement parallel to the force whose moment contribution is zero.
This is the principle of transmissibility, and it is why a rigid-body FBD needs
the line of action of each force but not the exact point on that line. A force
on a \emph{deformable} body is a \engterm{bound vector}: it is fixed to its
point of application, because sliding it along its line, though it leaves the
resultant and the net moment unchanged, changes the internal stress the body
carries, and a bar pushed at one end is not the same, inside, as a bar pulled
at the other. A couple, met later in this section, is a \engterm{free vector}:
it may be moved anywhere at all, off its line and to any point, without changing
its effect, because it has no resultant force to bind it. The three kinds are
not a taxonomy for its own sake. They tell the engineer exactly how much of the
point-of-application data a given problem needs: the line only, for
rigid-body equilibrium; the full point, for internal stress; and none of it,
for a couple. This chapter treats bodies as rigid, so its forces are sliding
vectors and its couples are free, and Volume V, taking up the deformable body,
promotes every force back to a bound vector the moment the stress inside the
member becomes the question.

% Figure: a single force vector resolved into components, with the
% parallelogram (polygon) construction of a resultant of two forces.
\begin{figure}[htbp]
\centering
\begin{tikzpicture}[
  force/.style={draw=engblue, -{Stealth}, very thick},
  comp/.style={draw=enggray, -{Stealth}, thin, dashed},
  res/.style={draw=engred, -{Stealth}, very thick},
  axis/.style={draw=engblack, -{Stealth}, thick},
  scale=1.0,
]
% === Left panel: one force resolved into components ===
\begin{scope}[xshift=0cm]
  \draw[axis] (0,0) -- (4.2,0) node[anchor=west, font=\small] {$x$};
  \draw[axis] (0,0) -- (0,3.4) node[anchor=south, font=\small] {$y$};
  % the force F at ~40 deg
  \draw[force] (0,0) -- (3.06,2.57) node[midway, above left, font=\small, engblue] {$\vec{F}$};
  % components
  \draw[comp] (0,0) -- (3.06,0);
  \draw[comp] (3.06,0) -- (3.06,2.57);
  \node[font=\scriptsize, below] at (1.53,0) {$F_x = F\cos\theta$};
  \node[font=\scriptsize, right] at (3.06,1.28) {$F_y = F\sin\theta$};
  % angle
  \draw[draw=enggray, thin] (0.9,0) arc (0:40:0.9);
  \node[font=\scriptsize] at (1.15,0.35) {$\theta$};
  \node[font=\small, anchor=north] at (1.8,-0.7) {(a) resolution};
\end{scope}
% === Right panel: parallelogram rule for the resultant ===
\begin{scope}[xshift=6.2cm]
  \coordinate (O) at (0,0);
  \coordinate (F1) at (3.2,0.6);
  \coordinate (F2) at (1.1,2.6);
  \coordinate (R) at (4.3,3.2);
  \draw[force] (O) -- (F1) node[midway, below, font=\small, engblue] {$\vec{F}_1$};
  \draw[force] (O) -- (F2) node[midway, left, font=\small, engblue] {$\vec{F}_2$};
  % parallelogram completion
  \draw[draw=enggray, thin, dashed] (F1) -- (R);
  \draw[draw=enggray, thin, dashed] (F2) -- (R);
  \draw[res] (O) -- (R) node[midway, above left, font=\small, engred] {$\vec{R}$};
  \node[font=\small, anchor=north] at (2.0,-0.7) {(b) parallelogram rule};
\end{scope}
\end{tikzpicture}
\caption[Force as a vector: resolution and the parallelogram rule]{A
force is a vector with magnitude, direction, and point of application.
(a) A force $\vec{F}$ at angle $\theta$ to the horizontal resolves into
Cartesian components $F_x=F\cos\theta$ and $F_y=F\sin\theta$ in the
chosen basis; the choice of basis is the engineer's and the physics is
unchanged by it. (b) Two forces sharing a point of application add by
the parallelogram rule: the resultant $\vec{R}=\vec{F}_1+\vec{F}_2$ is
the diagonal. The polygon rule (place the tail of $\vec{F}_2$ at the
head of $\vec{F}_1$) gives the same resultant and generalises to any
number of concurrent forces.}
\label{fig:vol03ch01:force-vector}
\end{figure}


A worked component decomposition fixes the notation. A cable pulls on a
bracket with a force of $400\,\text{N}$ directed $35^{\circ}$ above the
horizontal. In a horizontal-vertical basis the components are
$F_x = 400\cos 35^{\circ} \approx 328\,\text{N}$ and
$F_y = 400\sin 35^{\circ} \approx 229\,\text{N}$. If the bracket face
is itself tilted, say $20^{\circ}$ from horizontal, a basis aligned
with the face is more convenient: the angle between the force and the
face is $35^{\circ} - 20^{\circ} = 15^{\circ}$, so the in-face
component is $400\cos 15^{\circ} \approx 386\,\text{N}$ and the
out-of-face component is $400\sin 15^{\circ} \approx 104\,\text{N}$.
Same force, two bases, two pairs of numbers, one physical pull. The
script \texttt{code/resolve\_forces.py} carries this arithmetic for an
arbitrary list of forces and checks that they sum to zero, the test
every equilibrium problem reduces to.

A force applied to a rigid body has two effects: a translational effect
(through Newton's second law applied to the centre of mass) and a
rotational effect (through the force's moment about an axis). The
\engterm{moment} of $\vec{F}$ about a point $O$ is
\[
\vec{M}_O = \vec{r} \times \vec{F},
\]
where $\vec{r}$ is the position vector from $O$ to the force's point
of application. The moment is itself a vector, perpendicular to the
plane containing $\vec{r}$ and $\vec{F}$, with magnitude
$|\vec{r}||\vec{F}|\sin\theta$. In two-dimensional problems we treat
the moment as a signed scalar with a chosen sign convention
(counterclockwise positive, by default in this work).

Forces and moments together form a six-component \engterm{wrench} that
describes the mechanical state at a point. Statics chapters in this
book treat the wrench's components as the unknowns; dynamics chapters
let them drive the motion. Volume III's first habit is to write the
wrench down before doing anything else.

The moment of a single force depends on the reference point, but the
\emph{net} moment of a balanced force system does not, and that fact is
the working engineer's freedom in choosing where to sum moments. Moving
the reference point from $O$ to $O'$ changes each moment by $\vec{M}_{O'}
= \vec{M}_O + (\vec{r}_{O \to O'}) \times \vec{R}$, where $\vec{R}$ is
the resultant force. When $\vec{R} = \vec{0}$, as it is in equilibrium,
the correction term vanishes and the net moment is the same about every
point. The engineer exploits this by summing moments about a point
where one or more unknown forces pass through, so those unknowns drop
out of the moment equation. The ladder example below sums about the
foot precisely to kill the two reactions acting there.

\begin{mastery}
\textbf{The moment-transfer rule, derived.} Let a system of forces
$\vec{F}_i$ act at points $\vec{r}_i$. The net moment about $O$ is
$\vec{M}_O = \sum_i (\vec{r}_i - \vec{r}_O) \times \vec{F}_i$. About a
second point $O'$,
\[
\vec{M}_{O'} = \sum_i (\vec{r}_i - \vec{r}_{O'}) \times \vec{F}_i
= \sum_i (\vec{r}_i - \vec{r}_O) \times \vec{F}_i
+ (\vec{r}_O - \vec{r}_{O'}) \times \sum_i \vec{F}_i,
\]
which is $\vec{M}_{O'} = \vec{M}_O + (\vec{r}_O - \vec{r}_{O'}) \times
\vec{R}$. Two consequences. First, in equilibrium $\vec{R} = \vec{0}$,
so $\sum M = 0$ about one point implies $\sum M = 0$ about every point;
the three planar equilibrium equations may therefore be taken as
$\sum F_x = 0$, $\sum M_A = 0$, $\sum M_B = 0$ about two distinct points
$A, B$, a form often more convenient than the canonical one. Second,
when $\vec{R} \ne \vec{0}$, the locus of points about which the moment
is minimal defines the resultant's line of action; the wrench reduces
to a single force along that line plus a moment parallel to it. Exercise
1.9 asks for this proof and its three-dimensional analogue.
\end{mastery}

\subsection{Resolution in three dimensions and direction cosines}

The planar decomposition $\vec{F} = F_x\hat{\imath} + F_y\hat{\jmath}$
generalises to space by adding a third component, and the bookkeeping that
keeps a spatial force honest is the unit vector that carries its direction.
A force of magnitude $F$ pointing along a direction is written
$\vec{F} = F\,\hat{u}$, where $\hat{u}$ is the unit vector along the line of
action. When the direction is given by the angles $\theta_x$, $\theta_y$,
$\theta_z$ that the force makes with the three axes, the components are the
projections
\[
F_x = F\cos\theta_x, \quad F_y = F\cos\theta_y, \quad F_z = F\cos\theta_z,
\]
and the three cosines $\ell = \cos\theta_x$, $m = \cos\theta_y$,
$n = \cos\theta_z$ are the \engterm{direction cosines}. They are the
entries of $\hat{u}$, so they obey the identity that any unit vector obeys,
\[
\ell^2 + m^2 + n^2 = \cos^2\theta_x + \cos^2\theta_y + \cos^2\theta_z = 1.
\]
The identity is the spatial check that the planar problem did not need:
two direction cosines fix the third up to a sign, and a stated triple that
fails to square-sum to one is a stated direction that does not exist.

% Figure: a force in three dimensions, resolved into Cartesian components
% by its direction cosines. The angles the force makes with each axis fix
% the unit vector; the cosines square-sum to one.
\begin{figure}[htbp]
\centering
\begin{tikzpicture}[
  axis/.style={draw=engblack, -{Stealth}, thick},
  force/.style={draw=engblue, -{Stealth}, very thick},
  comp/.style={draw=enggray, thin, dashed},
  scale=1.0,
]
% an oblique 3D frame: x toward lower-right, y toward right, z up
\coordinate (O) at (0,0);
\coordinate (Xax) at (-2.3,-1.5);   % x axis (out of page, down-left)
\coordinate (Yax) at (4.0,0);       % y axis (to the right)
\coordinate (Zax) at (0,3.6);       % z axis (up)
\draw[axis] (O) -- (Xax) node[anchor=north east, font=\small] {$x$};
\draw[axis] (O) -- (Yax) node[anchor=west, font=\small] {$y$};
\draw[axis] (O) -- (Zax) node[anchor=south, font=\small] {$z$};
% the force tip P, an oblique point in the box
\coordinate (P) at (1.9,2.5);
% projection of P onto the x-y "floor": move down from P by the z part
\coordinate (Q) at (1.9,0.55);      % foot of vertical from P onto floor
% the force vector
\draw[force] (O) -- (P) node[midway, above left, font=\small, engblue] {$\vec{F}$};
% vertical drop to the floor (z component)
\draw[comp] (P) -- (Q);
% floor projection from O to Q
\draw[comp] (O) -- (Q);
% small marks for the three axis angles
\draw[draw=engred, thin] (0.85,0) arc (0:53:0.85);
\node[font=\scriptsize, engred] at (1.15,0.45) {$\theta_z$};
% label the three direction angles schematically near O
\node[font=\scriptsize, anchor=west] at (1.6,1.45)
  {$\hat{u}=(\cos\theta_x,\cos\theta_y,\cos\theta_z)$};
% mark the foot
\fill[engblack] (Q) circle (0.04);
\node[font=\scriptsize, anchor=north west] at (Q) {$Q$};
\fill[engblue] (P) circle (0.04);
\node[font=\scriptsize, anchor=south west] at (P) {$P$};
\end{tikzpicture}
\caption[A force in three dimensions and its direction cosines]{A force
$\vec{F}$ in space is fixed by its magnitude and a unit vector
$\hat{u}=(\cos\theta_x,\cos\theta_y,\cos\theta_z)$, whose entries are the
cosines of the angles the force makes with the three axes. The point $P$
is the head of the force; $Q$ is its projection onto the $xy$ floor, and
the dashed drop $PQ$ is the vertical ($z$) component. The three cosines
are not independent: because $\hat{u}$ is a unit vector, they satisfy
$\cos^2\theta_x+\cos^2\theta_y+\cos^2\theta_z=1$, so any two of the angles
fix the third up to a sign. Resolving a spatial force is the same act as
the planar resolution of figure~\ref{fig:vol03ch01:force-vector}, carried
out three times rather than once.}
\label{fig:vol03ch01:direction-cosines}
\end{figure}


In practice the direction of a force is rarely handed over as three angles.
It is read off the geometry, as the line from one known point to another.
A force of magnitude $F$ directed from a point $A$ toward a point $B$ has
unit vector
\[
\hat{u} = \frac{\vec{r}_{AB}}{|\vec{r}_{AB}|}
= \frac{(x_B - x_A)\hat{\imath} + (y_B - y_A)\hat{\jmath}
+ (z_B - z_A)\hat{k}}{\sqrt{(x_B - x_A)^2 + (y_B - y_A)^2 + (z_B - z_A)^2}},
\]
and the force is $\vec{F} = F\,\hat{u}$. The direction cosines fall out as
the entries of $\hat{u}$ without any angle ever being measured. This is the
working route for every cable, strut, and leg in a three-dimensional FBD:
write the position vector along the member, normalise it, scale by the
magnitude. The tripod example of section 1.6.7 is exactly this route run
three times, and the script \texttt{code/resolve\_forces.py} builds the
unit vector from two endpoints precisely so that the spatial arithmetic is
never done by hand.

\subsection{The moment as a cross product}

The two-dimensional ``force times perpendicular distance'' is a shadow of
the vector definition $\vec{M}_O = \vec{r} \times \vec{F}$. Written out in
the Cartesian basis with $\vec{r} = (r_x, r_y, r_z)$ and
$\vec{F} = (F_x, F_y, F_z)$, the cross product is the determinant
\[
\vec{M}_O = \vec{r} \times \vec{F}
= \begin{vmatrix}
\hat{\imath} & \hat{\jmath} & \hat{k} \\
r_x & r_y & r_z \\
F_x & F_y & F_z
\end{vmatrix}
= (r_y F_z - r_z F_y)\,\hat{\imath}
+ (r_z F_x - r_x F_z)\,\hat{\jmath}
+ (r_x F_y - r_y F_x)\,\hat{k}.
\]
Three scalar moments, one about each axis, fall out of the three
components. In a planar problem the forces and position vectors lie in the
$xy$ plane, so $r_z = F_z = 0$ and the first two components vanish; only
the $\hat{k}$ component survives,
\[
M_z = r_x F_y - r_y F_x,
\]
the signed scalar moment of the plane, positive counterclockwise. The
planar convention is not a separate rule; it is the one nonzero entry of
the determinant when the geometry is flat.

The magnitude of the moment is $|\vec{M}_O| = |\vec{r}|\,|\vec{F}|\sin\theta
= |\vec{F}|\,d$, where $d = |\vec{r}|\sin\theta$ is the perpendicular
distance from $O$ to the force's line of action, the \engterm{moment arm}.
The vector definition and the moment-arm picture agree, and the engineer
uses whichever is faster: the determinant when the components are at hand,
the moment arm when the perpendicular distance is obvious from a sketch.
The position vector $\vec{r}$ may run from $O$ to \emph{any} point on the
line of action, because sliding the application point along the line adds a
vector parallel to $\vec{F}$, whose cross product with $\vec{F}$ is zero.
This freedom, the \engterm{principle of transmissibility}, is what lets the
engineer take the moment arm to the most convenient point on the line.

\begin{principle}[Varignon's theorem]
The moment of a force about a point equals the sum of the moments of its
components about that point. The proof is the distributive law of the cross
product: if $\vec{F} = \vec{F}_1 + \vec{F}_2$, then
$\vec{r}\times\vec{F} = \vec{r}\times\vec{F}_1 + \vec{r}\times\vec{F}_2$.
The working use is the reverse: resolve an awkward force into components
whose moment arms are obvious (one horizontal, one vertical, each with a
clean perpendicular distance), take the moment of each, and add. The single
oblique moment arm, which a sketch gets wrong as often as right, is never
computed.
\end{principle}

\subsection{Couples, the free vector, and the force-couple resultant}

A \engterm{couple} is a pair of equal and opposite forces whose lines of
action do not coincide. Its resultant force is zero, $\vec{F} +
(-\vec{F}) = \vec{0}$, so it produces no net push, yet it carries a moment.
Take the two forces at points $\vec{r}_1$ and $\vec{r}_2$; the couple's
moment about any point $O$ is
\[
\vec{M} = \vec{r}_1 \times \vec{F} + \vec{r}_2 \times (-\vec{F})
= (\vec{r}_1 - \vec{r}_2) \times \vec{F},
\]
which contains the position vectors only through their difference. The
reference point $O$ has dropped out. A couple has the same moment about
every point in space, of magnitude $Fd$ with $d$ the perpendicular
separation of the two lines, and that point-independence makes the couple
a \engterm{free vector}: it may be moved to any location and turned in its
own plane without changing its effect on the body. A force is bound to its
line of action; a couple is free of any line.

% Figure: (a) a couple as two equal-and-opposite forces, whose moment is
% the same about every point; (b) reduction of a general force system to a
% single force at a chosen point plus a couple.
\begin{figure}[htbp]
\centering
\begin{tikzpicture}[
  force/.style={draw=engblue, -{Stealth}, very thick},
  redforce/.style={draw=engred, -{Stealth}, very thick},
  body/.style={draw=engblack, line width=1.0pt, fill=engblack!4},
  arc/.style={draw=engorange, -{Stealth}, thick},
  scale=1.0,
]
% === Left panel: a couple ===
\begin{scope}[xshift=0cm]
  \draw[body] (0,0) rectangle (2.6,1.6);
  % two equal-and-opposite horizontal forces, offset vertically by d
  \draw[force] (0.2,0.3) -- ++(1.6,0) node[anchor=west, font=\scriptsize, engblue] {$\vec{F}$};
  \draw[force] (2.4,1.3) -- ++(-1.6,0) node[anchor=east, font=\scriptsize, engblue] {$-\vec{F}$};
  % separation d
  \draw[draw=enggray, thin, {Stealth}-{Stealth}] (2.85,0.3) -- (2.85,1.3)
    node[midway, anchor=west, font=\scriptsize] {$d$};
  \node[font=\small, anchor=north] at (1.3,-0.45) {(a) a couple, $M=Fd$};
\end{scope}
% === Right panel: reduction to force-couple at O ===
\begin{scope}[xshift=6.4cm]
  \draw[body] (0,0) rectangle (2.6,1.6);
  \coordinate (Oc) at (0.3,0.3);
  \fill[engblack] (Oc) circle (0.05);
  \node[font=\scriptsize, anchor=north east] at (Oc) {$O$};
  % the single resultant force at O
  \draw[redforce] (Oc) -- ++(1.7,0.9) node[anchor=south west, font=\scriptsize, engred] {$\vec{R}$};
  % the residual couple, drawn as a curved arrow
  \draw[arc] (1.7,1.15) arc (-30:210:0.35)
    node[anchor=south, font=\scriptsize, engorange] {$\vec{M}_O$};
  \node[font=\small, anchor=north] at (1.3,-0.45) {(b) force-couple at $O$};
\end{scope}
\end{tikzpicture}
\caption[A couple and the force-couple reduction of a force system]{(a) A
couple is a pair of equal and opposite forces separated by a distance
$d$. Its resultant force is zero, so it exerts no net push, yet it carries
a moment $M=Fd$ that is the same about every point, which makes the couple
a free vector: it may be slid and turned anywhere on the body without
changing its effect. (b) Any system of forces and couples on a rigid body
reduces to a single resultant force $\vec{R}=\sum\vec{F}_i$ placed at a
chosen point $O$, together with a single resultant couple
$\vec{M}_O=\sum\vec{r}_i\times\vec{F}_i$ taken about $O$. Moving the point
$O$ leaves $\vec{R}$ unchanged and shifts $\vec{M}_O$ by the transfer
rule. Equilibrium is the statement that both the resultant force and the
resultant couple vanish.}
\label{fig:vol03ch01:couple-reduction}
\end{figure}


The couple is the device that lets a single force be moved off its line.
To shift a force $\vec{F}$ from its line of action to a parallel line
through a chosen point $O$, add at $O$ both $\vec{F}$ and $-\vec{F}$; the
added pair has zero net effect. Now $\vec{F}$ at $O$ stands alone, while
the original $\vec{F}$ on its old line and the new $-\vec{F}$ at $O$ form a
couple of moment $\vec{M}_O = \vec{r} \times \vec{F}$, with $\vec{r}$ from
$O$ to the original line. A force moved to a new point is the same force
plus a couple equal to its moment about the new point. This single move,
applied to every force in a system, collapses the whole system to one place.

The general reduction follows. Any system of forces $\vec{F}_i$ acting at
points $\vec{r}_i$, together with any applied couples, reduces about a
chosen point $O$ to a single \engterm{force-couple resultant}: a resultant
force
\[
\vec{R} = \sum_i \vec{F}_i
\]
acting at $O$, plus a resultant couple
\[
\vec{M}_O = \sum_i \vec{r}_i \times \vec{F}_i + \sum_j \vec{C}_j,
\]
the sum of every force's moment about $O$ and every applied couple. The
resultant force $\vec{R}$ is independent of $O$; the resultant couple
$\vec{M}_O$ depends on $O$ through the transfer rule derived above. The
force-couple pair carries the complete mechanical effect of the original
system. When $\vec{R}$ and $\vec{M}_O$ are not perpendicular, a further
reduction slides $\vec{R}$ to the line about which the couple is parallel
to the force, the \engterm{wrench}, a force and a collinear couple, the
irreducible form of a general spatial force system.

\subsection{The scalar equilibrium equations as vector components}

Equilibrium is the single pair of vector statements $\vec{R} = \vec{0}$ and
$\vec{M}_O = \vec{0}$. The scalar equations the rest of the chapter writes
are the components of that pair in the chosen basis. In three dimensions
the two vector equations have three components each,
\[
\sum F_x = 0, \quad \sum F_y = 0, \quad \sum F_z = 0, \qquad
\sum M_x = 0, \quad \sum M_y = 0, \quad \sum M_z = 0,
\]
six scalar equations, the count of the six rigid-body degrees of freedom
(three translations, three rotations) the equations pin down. In a planar
problem the forces lie in the $xy$ plane, so $\sum F_z = 0$ is empty and
the two in-plane moment equations $\sum M_x = 0$ and $\sum M_y = 0$ are
satisfied identically; only $\sum F_x = 0$, $\sum F_y = 0$, and
$\sum M_z = 0$ carry information, the familiar three. The planar
equilibrium equations are not a special model. They are the three nonzero
components of the spatial conditions when the geometry is flat, and the
three empty ones are the price-free record that a flat problem cannot
resist an out-of-plane force or an in-plane tilting moment without leaving
the plane.

\begin{archetype}[Balance]
A rigid body is in static equilibrium when the resultant force and the
resultant moment (about any point) are both zero. In three dimensions
this is six scalar equations; in two dimensions, three. Every static
problem in this book ultimately reduces to writing those equations and
solving them. The archetype is the same whether the body is a beam, a
truss, a parked car, a hand holding a coffee cup, or the inner casing
of a turbocharger: identify the body, list the forces, write the
balance, solve. The pleasure of mechanics is that the same six
equations resolve cases separated by ten orders of magnitude in scale.
\end{archetype}

\section{Newton's laws as observations}\pathtag{core}

Newton's three laws of motion, stated by Newton in the
\emph{Principia} of 1687 and given here in their working engineering
form, are the empirical content the rest of the book operates on
\cite[Axioms, pp.~416--417]{hist:newton1687}.

\begin{principle}[Newton's first law]
A body subject to no net force moves with constant velocity. Constant
velocity includes zero velocity. The frame in which this holds is
called an inertial frame; for the purposes of nearly all terrestrial
engineering, a frame fixed to the local ground is inertial to within
the precision the engineer needs.
\end{principle}

\begin{principle}[Newton's second law]
The net force on a body equals the rate of change of its momentum,
$\vec{F}_{\text{net}} = d\vec{p}/dt$. For a body of constant mass $m$,
this reduces to $\vec{F}_{\text{net}} = m\vec{a}$. The two forms are
not interchangeable when mass changes: rockets, conveyors, and any
control volume with mass flow across its boundary require the momentum
form.
\end{principle}

\begin{principle}[Newton's third law]
Every force exerted by body $A$ on body $B$ is accompanied by an equal
and opposite force exerted by $B$ on $A$. The two forces act on
different bodies. A free-body diagram of $A$ shows only the force from
$B$ on $A$; the reaction lives in the FBD of $B$.
\end{principle}

The third law is the law a student most often gets wrong. Two students
arguing about whether ``the bridge pulls back on the truck'' have
usually conflated the FBD of the bridge with the FBD of the truck. The
discipline of the next section eliminates the confusion by forcing the
engineer to choose one body and stay with it.

The first law hides an assumption the working engineer should be able to
state: the frame in which it holds. A frame is \engterm{inertial} when a body
free of net force moves in a straight line at constant speed in it, and the
second law $\vec{F} = m\vec{a}$ is written for accelerations measured in such a
frame. A frame that accelerates is not inertial, and in it the second law
acquires extra terms, the inertial forces the D'Alembert box names, that have
no material source. The Earth's surface is not truly inertial: it turns once a
day and orbits once a year, so a frame fixed to the ground accelerates, and a
body moving in it feels the Coriolis and centrifugal terms that turn a
long-range shell, spin a weather system, and swing a Foucault pendulum. The
question for the engineer is whether those terms matter at the scale in hand.
For a bracket, a beam, or a parked car, the ground-fixed frame is inertial to
far better than the precision any FBD needs: the centrifugal acceleration from
the Earth's spin is about $0.03\,\si{\meter\per\second\squared}$ at the
equator, three parts in a thousand of $g$, and the Coriolis term on a body
moving at workshop speeds is smaller still. For a rifle shot across a valley, a
guided missile, or a turbine blade whose reference frame spins thousands of
times a minute, the terms are the whole problem, and the FBD is drawn in an
inertial frame or corrected with the inertial forces that stand in for one. The
half-life rule applies to the frame as much as to the model: name it, check its
acceleration against the precision the problem asks, and record that the check
was made.

Newton stated the laws as axioms, in the geometric style of his time.
For our purposes they are best read as observations: they are the
patterns that survived every controlled experiment in mechanical
behaviour for two and a half centuries before relativity and quantum
mechanics revealed the limits of their domain. Inside that domain,
which contains every artifact this volume designs, they are exact to
the precision the engineer can measure.

The limits of the domain matter. Newton's laws fail when the speed
approaches the speed of light (special relativity replaces them with
covariant equations), when gravitational fields are strong enough to
curve spacetime appreciably (general relativity replaces gravity with
geodesic motion), and when the action approaches Planck's constant
(quantum mechanics replaces deterministic trajectories with
probability amplitudes). For the artifacts in this book, none of
those regimes applies. We will note one boundary case in chapter 13
of this volume, where compressibility effects in air make a Mach
correction unavoidable at speeds above about $0.3 M$. The point of
naming the limits is the half-life rule of the book: a model that is
honest about where it stops is a model the reader can use.

\begin{mastery}
\textbf{D'Alembert's principle: a moving body drawn as a static one.} The
free-body diagram is a statics instrument, yet most artifacts move. D'Alembert's
principle, published in his \emph{Trait\'e de dynamique} of 1743, extends the
diagram to accelerating bodies by a change of viewpoint. Newton's second law
$\vec{F}_{\text{net}} = m\vec{a}$ rearranges to $\vec{F}_{\text{net}} - m\vec{a}
= \vec{0}$, which reads as a statement of equilibrium if the term $-m\vec{a}$ is
treated as a force. That term is the \engterm{inertial force} (or D'Alembert
force): a fictitious force of magnitude $m|\vec{a}|$ directed opposite the
acceleration, added to the diagram at the centre of mass. With it drawn, the
accelerating body is in a state of \engterm{dynamic equilibrium}, and every
technique the rest of this chapter builds for statics, moment-summing to kill a
reaction, the three planar equations, the choice of axes, applies unchanged.
The rotational analogue adds an inertial couple $-I_G\dot\omega$ opposing the
angular acceleration. A car braking at $0.5g$ is drawn as a static body with a
forward inertial force $0.5mg$ at its centre of mass, and the load transfer
that pitches weight onto the front axle falls straight out of a moment balance,
no differential equation in sight. The caution is that the inertial force is a
bookkeeping device, not a physical push: it exists only in the accelerating
frame, has no third-law partner, and disappears the moment the analysis returns
to an inertial frame. Used with that caution, D'Alembert's principle is the
bridge that lets the FBD discipline carry from the parked car of section 1.6.4
into every dynamics chapter that follows, and it is why the diagram is worth
mastering as statics before a single body is allowed to move.
\end{mastery}

\section{Free-body diagrams: the discipline}\pathtag{core}

A free-body diagram is a representation of one chosen body, isolated
from its surroundings, with every external force and moment acting on
it shown as a labelled vector. The diagram is not a picture of the
artifact. The diagram is a picture of the forces.

The discipline has six steps. We state them as a procedure because the
procedure works, and beginners who skip steps produce diagrams that
hide loads they will later be punished for missing.

\begin{enumerate}
  \item \engterm{Choose one body.} Decide which body the diagram is
    for. A single beam, a single block, a single joint, the rope only,
    the rope and the weight as one body. Once chosen, do not let the
    diagram drift to include or exclude parts.
  \item \engterm{Isolate.} Mentally cut the body away from everything
    it touches. The body now floats in space.
  \item \engterm{Replace each contact with its forces.} At every
    boundary where the body touched something else, draw the force
    that the something else exerted on the body. A pin replaces a
    rotational constraint with two perpendicular reactions. A roller
    replaces a single normal constraint with a single normal
    reaction. A weld or built-in connection replaces a rotational and
    two translational constraints with two perpendicular reactions
    and a moment. A surface contact replaces the surface with a
    normal force and, if friction is present, a tangential friction
    force.
  \item \engterm{Add the body's weight} at its centre of mass, as a
    vector $-mg\hat{j}$ (with $\hat{j}$ vertical up). On Earth, $g
    \approx 9.81\,\si{\meter\per\second\squared}$; on Mars, about
    $3.71\,\si{\meter\per\second\squared}$; on the lunar surface, about
    $1.62\,\si{\meter\per\second\squared}$. The body's weight is a
    real external force; it is not a contact force, and it must not be
    omitted because the body looks ``at rest.''
  \item \engterm{Label every force.} Magnitude, direction, point of
    application. Use one symbol per unknown. Reactions get their own
    names ($A_x$, $A_y$, $B_y$, etc.). Friction is named
    separately from its normal companion. Tensions are named per
    rope segment; a rope passing over a pulley has different
    tensions on either side unless the pulley is explicitly
    frictionless.
  \item \engterm{State the equations.} For static equilibrium:
    \[
    \sum F_x = 0, \quad \sum F_y = 0, \quad \sum M_O = 0.
    \]
    For dynamics: substitute $m\vec{a}$ for the force sum and
    $I_O \dot\omega$ (or its full inertia-tensor version in three
    dimensions) for the moment sum.
\end{enumerate}

A diagram that omits a step omits a force. A diagram that omits a
force gives a wrong answer with a confident face. The next section
walks through three worked diagrams; the failure section at the end of
the chapter studies a real building in which a missing step killed
114 people.

% Figure: a block on a rough inclined plane, real scene (left) and
% free-body diagram (right). The heart of the chapter: turning a scene
% into vectors.
\begin{figure}[htbp]
\centering
\begin{tikzpicture}[
  weight/.style={draw=engblue, -{Stealth}, very thick},
  normal/.style={draw=engred, -{Stealth}, very thick},
  fric/.style={draw=engamber, -{Stealth}, very thick},
  axis/.style={draw=enggray, -{Stealth}, thin, dashed},
  scale=1.0,
]
% === Left: the scene ===
\begin{scope}[xshift=0cm]
  % incline triangle, angle ~ 25 deg
  \coordinate (A) at (0,0);
  \coordinate (B) at (5.2,0);
  \coordinate (C) at (5.2,2.42);
  \draw[draw=engblack, very thick] (A) -- (B) -- (C) -- cycle;
  % hatching on ground
  \foreach \x in {0.3,0.8,...,4.8} {
    \draw[draw=enggray, thin] (\x,0) -- (\x-0.25,-0.25);
  }
  % angle
  \draw[draw=enggray, thin] (1.1,0) arc (0:25:1.1);
  \node[font=\scriptsize] at (1.4,0.22) {$\theta$};
  % block on the incline
  \begin{scope}[shift={(3.3,1.54)}, rotate=25]
    \draw[draw=engblack, fill=engblue!12, thick] (-0.55,0) rectangle (0.55,0.85);
    \node[font=\scriptsize] at (0,0.42) {$m$};
  \end{scope}
  \node[font=\small, anchor=north] at (2.6,-0.7) {(a) the scene};
\end{scope}
% === Right: the FBD of the block ===
\begin{scope}[xshift=8.3cm, yshift=0.6cm]
  % the block, tilted 25 deg, centered at origin
  \begin{scope}[rotate=25]
    \draw[draw=engblack, fill=engblue!8, thick] (-0.7,-0.45) rectangle (0.7,0.45);
  \end{scope}
  \coordinate (G) at (0,0);
  % weight straight down
  \draw[weight] (G) -- ++(0,-2.3) node[anchor=north, font=\small, engblue] {$mg$};
  % normal perpendicular to incline (incline at 25 deg, normal at 25+90=115 deg)
  \draw[normal] (G) -- ++(115:2.2) node[anchor=south, font=\small, engred] {$N$};
  % friction up the incline (incline direction 25 deg, up-slope = 25+180=205 deg)
  \draw[fric] (G) -- ++(205:2.0) node[anchor=east, font=\small, engamber] {$F_f$};
  % local axes along/normal to incline
  \draw[axis] (G) -- ++(25:1.4) node[anchor=west, font=\scriptsize, enggray] {$x$ (up-slope $+$)};
  \node[font=\small, anchor=north] at (0,-2.9) {(b) free-body diagram};
\end{scope}
\end{tikzpicture}
\caption[A block on a rough incline and its free-body diagram]{The
move at the centre of the chapter. (a) A block of mass $m$ rests on a
rough plane inclined at $\theta$. (b) Isolate the block and replace
every contact with the force it exerts: the weight $mg$ acts straight
down at the centre of mass, the normal force $N$ acts perpendicular to
the incline surface, and the friction force $F_f$ acts along the
surface, here drawn up-slope to oppose impending downhill slip. The
basis is chosen along and perpendicular to the incline so that only the
weight needs resolving: $\sum F_x: F_f - mg\sin\theta = 0$ and $\sum
F_y: N - mg\cos\theta = 0$. No part of the incline appears in the
diagram; only the forces it transmits.}
\label{fig:vol03ch01:fbd-block-incline}
\end{figure}


Figure~\ref{fig:vol03ch01:fbd-block-incline} runs the six steps on the
simplest non-trivial case, a block on a rough incline. The scene on the
left has a block, a wedge, a floor, and the planet. The diagram on the
right has three vectors and nothing else. The wedge is gone; in its
place is the normal force it transmits. The planet is gone; in its
place is the weight it pulls with. The choice of basis along and across
the incline means that two of the three forces already lie on axes, and
only the weight needs resolving. This is the payoff of step zero, the
choice of axes: the arithmetic that follows is one sine and one cosine
rather than three oblique projections.

\begin{principle}[The FBD is exhaustive]
Every external force acting on the chosen body is drawn. No internal
force (force from one part of the body on another part) is drawn. No
imagined force is drawn. The diagram and the equilibrium equations are
in one-to-one correspondence: one force on the diagram is one term in
the equation.
\end{principle}

A common beginner error is to draw the force the body exerts on its
support rather than the force the support exerts on the body. The
third law guarantees these are equal in magnitude and opposite in
direction, but they appear in different diagrams. A box sitting on a
table receives a normal force from the table, drawn upward; the box
exerts a force on the table, drawn downward, but that force belongs in
the FBD of the table, not the box. Crossing the bookkeeping is the
fastest path to a sign error that compounds through the rest of the
problem.

The standard reference text for FBD methodology in this book is
Meriam, Kraige, and Bolton's \emph{Engineering Mechanics}
\cite{text:meriam-kraige2018}, which we will cite again in chapters 2,
4, and 5 of this volume.

\section{Constraints: rolling, pinning, fixing}\pathtag{standard}

A constraint is a geometric restriction on a body's motion. Every
constraint translates into a reaction force or moment on the FBD. The
working engineer learns to read a connection diagram and write down the
reactions it implies without thinking. We list the canonical
two-dimensional constraints; the three-dimensional generalisations are
in chapter 2 of this volume.

\begin{itemize}[leftmargin=*]
  \item \engterm{Smooth surface contact (or roller).} Restricts motion
    along the surface normal only. Reaction: one force, along the
    surface normal, pointing into the body. No friction, no moment.
  \item \engterm{Rough surface contact.} Restricts motion along the
    normal and, up to the friction limit, along the tangent.
    Reaction: a normal force and a tangential friction force, the
    latter bounded by $|F_f| \le \mu_s N$. The inequality is the
    subject of the next section.
  \item \engterm{Pin (hinge).} Restricts translation in two
    directions; allows rotation. Reaction: two perpendicular force
    components ($A_x$, $A_y$), no moment.
  \item \engterm{Built-in (fixed) support.} Restricts translation in
    two directions and rotation. Reaction: two perpendicular force
    components and a moment ($A_x$, $A_y$, $M_A$).
  \item \engterm{Rope or cable.} A flexible member that supports
    tension only. Reaction along the rope, away from the body. A rope
    that goes slack carries no force; a rope in compression has
    buckled (no rope ever pushes).
  \item \engterm{Rigid link (two-force member).} A straight member
    pinned at both ends, with no transverse load and no weight,
    transmits force along its axis only. The direction of the
    reaction is therefore along the link. Trusses are assemblies of
    such members.
  \item \engterm{Frictionless pulley.} An ideal pulley redirects a
    rope without changing its tension. Real pulleys have bearing
    friction and a finite rope-stiffness term, both of which Coulomb
    measured in his 1781 memoir on simple machines
    \cite{hist:coulomb1781}.
\end{itemize}

% Figure: the canonical 2D supports and the reactions each one implies.
% Roller (1 reaction), pin (2 reactions), fixed support (2 + moment),
% cable (1 along the cable).
\begin{figure}[htbp]
\centering
\begin{tikzpicture}[
  rxn/.style={draw=engred, -{Stealth}, very thick},
  mom/.style={draw=engred, -{Stealth}, very thick},
  member/.style={draw=engblack, very thick},
  scale=1.0,
]
% ==== Roller ====
\begin{scope}[xshift=0cm]
  \draw[member] (0,1.4) -- (1.6,1.4);
  \draw[draw=engblack, fill=enggray!15] (0.8,1.1) circle (0.28);
  \foreach \x in {0.2,0.5,...,1.4} {\draw[draw=enggray, thin] (\x,0.7) -- (\x-0.2,0.45);}
  \draw[draw=enggray, thick] (0.1,0.7) -- (1.5,0.7);
  \draw[rxn] (0.8,1.1) -- (0.8,2.6) node[anchor=south, font=\small, engred] {$N$};
  \node[font=\small, anchor=north] at (0.8,0.2) {roller: 1};
\end{scope}
% ==== Pin ====
\begin{scope}[xshift=4.0cm]
  \draw[member] (0.8,1.4) -- (0.8,2.9);
  \draw[draw=engblack, fill=enggray!15] (0.8,1.1) -- (0.45,0.55) -- (1.15,0.55) -- cycle;
  \foreach \x in {0.3,0.6,...,1.3} {\draw[draw=enggray, thin] (\x,0.55) -- (\x-0.2,0.3);}
  \draw[draw=enggray, thick] (0.1,0.55) -- (1.5,0.55);
  \draw[draw=engblack, fill=white] (0.8,1.1) circle (0.1);
  \draw[rxn] (0.8,1.1) -- (2.2,1.1) node[anchor=west, font=\small, engred] {$A_x$};
  \draw[rxn] (0.8,1.1) -- (0.8,2.6) node[anchor=south, font=\small, engred] {$A_y$};
  \node[font=\small, anchor=north] at (0.8,0.05) {pin: 2};
\end{scope}
% ==== Fixed support ====
\begin{scope}[xshift=8.2cm]
  \draw[draw=engblack, thick] (0,0.4) -- (0,2.4);
  \foreach \y in {0.5,0.85,...,2.3} {\draw[draw=enggray, thin] (0,\y) -- (-0.3,\y-0.3);}
  \draw[member] (0,1.4) -- (1.7,1.4);
  \draw[rxn] (0,1.4) -- (1.4,1.4) node[anchor=south, font=\small, engred] {$A_x$};
  \draw[rxn] (0,1.4) -- (0,2.9) node[anchor=south, font=\small, engred] {$A_y$};
  \draw[mom] (0.55,1.4) arc (0:300:0.4);
  \node[font=\small, engred] at (1.1,2.1) {$M_A$};
  \node[font=\small, anchor=north] at (0.6,0.05) {fixed: 3};
\end{scope}
% ==== Cable ====
\begin{scope}[xshift=12.2cm]
  \draw[draw=enggray, very thick] (0.8,2.9) -- (0.8,1.2);
  \draw[draw=engblack, fill=engblue!12, thick] (0.45,0.5) rectangle (1.15,1.2);
  \node[font=\scriptsize] at (0.8,0.85) {$m$};
  \draw[rxn] (0.8,1.2) -- (0.8,2.6) node[anchor=west, font=\small, engred] {$T$};
  \node[font=\small, anchor=north] at (0.8,0.3) {cable: 1};
\end{scope}
\end{tikzpicture}
\caption[Canonical two-dimensional supports and their reactions]{Each
support type translates into a fixed set of reaction unknowns on the
free-body diagram. A roller (or smooth contact) restricts motion along
the surface normal only and contributes one reaction $N$. A pin (hinge)
restricts translation in two directions but allows rotation, and
contributes two perpendicular reactions $A_x, A_y$. A built-in (fixed)
support restricts translation in two directions and rotation, and
contributes two reactions plus a moment $M_A$. A cable supports tension
only, directed along its own axis away from the body. The count of
reaction unknowns drives the statical-determinacy bookkeeping: a planar
rigid body has three equilibrium equations, so a determinate support
set sums to exactly three reactions.}
\label{fig:vol03ch01:constraint-reactions}
\end{figure}


\engterm{Rolling without slipping} is a kinematic constraint: the
contact point of the wheel has zero instantaneous velocity. The
constraint introduces a relationship between the wheel's translational
and angular velocities, $v = r\omega$, but does not specify the
friction force; the friction force adjusts to whatever value enforces
the constraint, up to the static-friction limit. When the required
friction exceeds the limit, the wheel slips, and the friction force
saturates at the kinetic value.

The number of independent reactions a constraint introduces is its
\engterm{degree}. A two-dimensional pin contributes two; a fixed
support contributes three; a roller contributes one. The body's
\engterm{statical determinacy} is the count: a planar rigid body with
three equilibrium equations is statically determinate if the
reactions sum to three. Fewer reactions and the body is a mechanism
(it moves); more and the body is statically indeterminate (the
deflection of the body itself is needed to find the reactions, a
topic of chapter 3 of this volume).

A horizontal beam fixed at both ends shows the count at work. Each
built-in support contributes three reactions, so the support count is
six. The beam has three planar equilibrium equations. The excess,
$6 - 3 = 3$, is the \engterm{degree of static indeterminacy}: three
reaction quantities cannot be found from equilibrium alone, and the
beam's bending stiffness enters the bookkeeping. Replace one fixed end
with a pin (two reactions) and a roller elsewhere (one reaction), and
the count becomes $3 + 1 = 4$, still one redundant. Replace it instead
with a pin plus a single roller, $2 + 1 = 3$, and the beam is
determinate: every reaction follows from equilibrium. The determinacy
count is the first thing to compute on any new structure, because it
decides whether the methods of this chapter suffice or whether
chapter 3's deflection methods are required.

The count is necessary but not sufficient, and the exception is worth naming
because it fails silently. A body can have exactly the right number of
reactions and still be unstable if those reactions are arranged badly, a
condition called \engterm{improper constraint}. Two arrangements cause it.
First, if all three reaction lines of a planar body are parallel, the body can
still translate in the one direction none of them opposes: three vertical
rollers under a beam count three but leave the beam free to slide sideways, and
the equilibrium equations return no finite horizontal reaction because there is
none to be had. Second, if all three reaction lines are concurrent, meeting at
a single point, the body can still rotate about that point: no reaction has a
moment about the common point, so any applied moment about it goes unresisted,
and the moment equation about that point reduces to $0 = M_{\text{applied}}$, a
contradiction with no solution. The formal signature of both cases is the same:
the equilibrium matrix, whose columns are the reaction directions and their
moment arms, drops rank, so the system that looked square is singular. A count
of three reactions guarantees determinacy only when the three are independent,
neither all parallel nor all concurrent. The practical check is to look before
counting: three reactions that share a line of action or a common point are a
mechanism wearing the costume of a determinate structure, and a solver handed
such a system returns either a division by zero or a plausible-looking number
that a real load would fling the body straight past. The concurrence theorem
that solves a three-force body is the same geometry that condemns a
three-reaction body: when the lines meet at a point, rotation about that point
is free.

The determinacy count is worth running on a handful of stock connections until
it is automatic, because the whole decision of whether this chapter's methods
suffice rides on it. Take a single horizontal beam and vary only how its two
ends are held. Pin at one end (two reactions), roller at the other (one): three
reactions against three planar equations, determinate, and every reaction
follows from equilibrium. This is the simply supported beam the worked examples
lean on. Now fix one end fully (three reactions) and free the other: still
three, determinate, the cantilever. Fix one end and pin the other (three plus
two): five reactions, two redundant, the propped cantilever that needs
chapter 3's deflection methods. Fix both ends (three plus three): six, three
redundant, the built-in beam met earlier. Pin both ends (two plus two): four,
one redundant. The pattern is arithmetic, $r$ reactions against $e = 3$
equations, with $r - e$ the degree of indeterminacy, and it is worth tabulating
the four cases side by side once so the count becomes a reflex rather than a
derivation. The trap the count alone does not catch is the improper arrangement
above: three rollers all vertical count three but leave the beam free to slide,
so a determinate count is a necessary condition the engineer confirms by
checking that the reactions are not all parallel and not all concurrent before
trusting it.

A three-dimensional body raises the stakes because it has six equations, not
three, and a connection that looks fully fixed on a flat drawing can leave a
degree of freedom out of the plane. A door on two hinges is the everyday case:
each hinge, idealised, restrains translation in the two directions across its
pin but allows rotation about the pin axis and, unless it has a thrust collar,
slide along it. Two such hinges together restrain five of the six degrees of
freedom and leave the door free to swing, which is the one motion a door is for.
Remove the assumption that the lower hinge takes the vertical thrust and the
door is a mechanism in the vertical direction, held up only by whatever collar
or shim actually carries its weight, which is why a heavy gate whose bottom
hinge cannot take thrust drops until something else does. Counting to six, and
naming which degree each connection removes, is the spatial discipline that the
planar count of three only rehearses.

\section{Friction}\pathtag{standard}

Friction is the tangential force at a contact that resists relative
sliding. Coulomb's 1781 prize memoir on simple machines is the source
of what is now called Coulomb friction
\cite{hist:coulomb1781}.

The Coulomb model states:

\begin{itemize}[leftmargin=*]
  \item \engterm{Static friction} adjusts to whatever value is
    required to prevent sliding, up to a maximum
    \[
    |F_{f,\max}| = \mu_s N,
    \]
    where $N$ is the normal force at the contact and $\mu_s$ is the
    coefficient of static friction, a function of the two materials
    in contact and their surface conditions.
  \item \engterm{Kinetic friction} acts opposite to the direction of
    sliding, with magnitude
    \[
    |F_k| = \mu_k N,
    \]
    where $\mu_k$ is the coefficient of kinetic friction. Typically
    $\mu_k < \mu_s$.
  \item Within the Coulomb model the friction coefficients do not
    depend on the apparent contact area or, to leading order, on the
    sliding speed. Both independences are empirical approximations
    that break down under high contact pressure, lubrication, and
    high sliding speed; the corrections live in tribology and
    surface mechanics, with the canonical engineering reference being
    Shigley's \emph{Mechanical Engineering Design}
    \cite[ch.~7--8]{text:shigley2020}.
\end{itemize}

% Figure: friction force versus applied force, showing the static plateau
% (Ff = applied), the slip threshold mu_s N, and the drop to kinetic mu_k N.
\begin{figure}[htbp]
\centering
\begin{tikzpicture}
\begin{axis}[
  width=12cm, height=7.5cm,
  xlabel={applied tangential force $P$}, ylabel={friction force $F_f$},
  xmin=0, xmax=10, ymin=0, ymax=7,
  axis lines=left,
  xtick=\empty, ytick={4,6}, yticklabels={$\mu_k N$, $\mu_s N$},
  tick label style={font=\scriptsize},
  label style={font=\small},
  legend pos=south east,
  legend style={font=\scriptsize},
  legend cell align=left,
]
% static region: Ff = P, up to P = mu_s N = 6
\addplot[engblue, very thick, domain=0:6] {x};
\addlegendentry{static: $F_f = P \le \mu_s N$}
% slip: drop from 6 to 4 at P = 6
\addplot[engred, very thick, dashed] coordinates {(6,6) (6,4)};
% kinetic: constant at 4 once sliding
\addplot[engred, very thick, domain=6:10] {4};
\addlegendentry{kinetic: $F_f = \mu_k N$}
% the 45-degree guide showing P
\addplot[enggray, dotted, thin, domain=0:10] {x};
% markers
\node[engblack, font=\scriptsize, anchor=south west] at (axis cs:2.5,2.7) {$F_f = P$};
\node[engred, font=\scriptsize, anchor=west] at (axis cs:6.1,5.0) {slip};
\draw[draw=enggray, dashed, thin] (axis cs:6,0) -- (axis cs:6,6);
\node[font=\scriptsize, anchor=north] at (axis cs:6,-0.15) {$P=\mu_s N$};
\end{axis}
\end{tikzpicture}
\caption[Coulomb friction: static plateau and kinetic drop]{The
friction force is not a fixed quantity. While the body does not slide,
static friction takes whatever value enforces $\sum F_x = 0$ along the
tangent: $F_f = P$, the blue line of slope one. This holds only up to
the threshold $\mu_s N$. At $P = \mu_s N$ the contact reaches its
static limit; the body breaks free and the friction force drops to the
kinetic value $\mu_k N$ (red), typically lower because $\mu_k < \mu_s$.
The common error is to write $F_f = \mu_s N$ when the body is nowhere
near the threshold; the correct statement below threshold is the
inequality $|F_f| \le \mu_s N$, with the equality reserved for the slip
condition itself.}
\label{fig:vol03ch01:friction-cone}
\end{figure}


Representative dry-sliding coefficients, current as of 2025 per
\textcite{text:meriam-kraige2018}: steel on steel
$\mu_s \approx 0.6$, $\mu_k \approx 0.4$; rubber on dry concrete
$\mu_s \approx 1.0$, $\mu_k \approx 0.7$; PTFE on steel
$\mu_s \approx 0.04$, $\mu_k \approx 0.04$; ice on steel
$\mu_s \approx 0.03$. These values are working figures, not standards;
in any specific design the coefficient is measured for the actual
contact under the actual conditions.

A common error in FBDs is to write the friction force as $F_f = \mu_s
N$ when the body is not on the verge of sliding. The correct statement
is the inequality $|F_f| \le \mu_s N$; the equality holds only at the
slip threshold. Outside that threshold, $F_f$ is whatever value makes
$\sum F_x = 0$ along the contact tangent. Treating the inequality as
an equality assigns the friction a magnitude it does not have, and
produces wrong answers in problems where the static friction is
nowhere near its maximum. The data file
\texttt{data/friction\_coefficients.csv} collects the working
coefficients used throughout this chapter, with their conditions; the
script \texttt{code/block\_on\_incline.py} sweeps the incline angle and
reproduces the static plateau and the kink at the slip threshold drawn
in figure~\ref{fig:vol03ch01:friction-cone}.

The area-independence that Coulomb's model asserts is the surprising claim,
and its explanation is that the area that matters is not the one the eye sees.
Two surfaces touch only at the summits of their microscopic roughness, so the
\engterm{real} area of contact is a tiny fraction of the \engterm{apparent}
area, the fraction growing in proportion to the normal load as the summits
squash. Friction is set by the shearing of those load-bearing junctions, so it
scales with the real area, which scales with the load, which is why the
friction force tracks $N$ and not the footprint. Doubling the apparent area
halves the pressure and halves the real contact per unit apparent area, leaving
the total real contact, and the friction, unchanged. The model breaks where the
picture breaks: under very high pressure the junctions saturate and friction
climbs less than linearly with load; with a soft or tacky material (rubber,
adhesive tape) the real area approaches the apparent and a wider tyre does grip
harder; and a lubricant film that keeps the summits apart replaces solid
shearing with fluid shearing and drops the coefficient by an order of magnitude.
The engineer keeps $F_f = \mu N$ as the working law and remembers that each of
its independences is a junction picture that a real contact can violate.

Rolling is the other case the block-on-incline picture omits. A wheel that
rolls without slipping still meets a resistance, but it is not the sliding
friction of the Coulomb model; it is \engterm{rolling resistance}, and its
origin is the small deformation of wheel and surface at the contact. The
material ahead of the contact is being compressed and the material behind is
recovering, and because neither is perfectly elastic the pressure distribution
is not symmetric: it crowds forward, so the normal reaction's resultant sits a
small distance $a$ ahead of the wheel centreline. That offset gives the normal
force a moment $N a$ opposing rotation, which the tow force must overcome. The
resistance is written $F_r = (a/r)\,N = C_{rr}\,N$, with $r$ the wheel radius
and $C_{rr} = a/r$ the \engterm{coefficient of rolling resistance}, a number
one to two orders of magnitude below a sliding coefficient: a steel railway
wheel on rail runs near $C_{rr} \approx 0.001$, a car tyre on asphalt near
$0.01$, a soft tyre on sand near $0.3$. The tiny coefficient is why rolling
beats dragging and why rail beats road for freight: the wheel does not fight
$\mu_s$ at all, only the hysteresis offset $a$, and a harder, larger wheel on a
harder surface shrinks $a/r$ toward zero. Rolling resistance is a separate
mechanism from the friction that drives the wheel, and a free-body diagram of a
rolling vehicle carries both: static friction at the contact to propel or brake
(bounded by $\mu_s N$) and the small forward offset of the normal reaction that
the drive must continuously overcome to keep the wheel turning.

Friction sets the angle at which a block slides down an incline, but a tall
enough block topples before it slides, and the two thresholds are worth carrying
together because the friction analysis is only half the picture. A block of
width $b$ and height $H$ sits on a plane tilted to an angle $\theta$. Sliding
impends when the down-slope weight component reaches the friction limit, at the
angle $\tan\theta_{\text{slide}} = \mu_s$, the result the friction cone gives.
Tipping impends when the weight's line of action, which drops vertically from
the block's centre of mass, falls outside the lower edge of the base, at the
angle $\tan\theta_{\text{tip}} = b/H$: past that tilt the weight has no moment
arm inside the contact to resist the overturn and the block rotates about its
lower edge. The block does whichever comes first as $\theta$ is raised, so it
slides if $\mu_s < b/H$ and tips if $\mu_s > b/H$. A short, wide block on a
slick plane slides; a tall, narrow block on a grippy plane tips. The two are
equal at the critical aspect ratio $b/H = \mu_s$, and the same competition, read
with a horizontal push instead of a tilt, is the filing-cabinet problem of
section 1.6.10 and the tower-crane stability of the third case study, the tilt
and the push being two ways to write one moment-versus-force contest. A friction
analysis that reports the slide angle and stops has answered half the question
for any block tall enough to matter, and the missing half is a moment balance
the friction coefficient never enters.

\begin{mastery}
\textbf{The friction cone.} Combine the normal force $N$ and the
friction force $F_f$ at a contact into a single contact reaction
$\vec{R}_c$. Its direction makes an angle $\phi$ with the surface normal
where $\tan\phi = F_f / N$. Static equilibrium requires $F_f \le \mu_s
N$, so $\tan\phi \le \mu_s$, that is $\phi \le \phi_s = \arctan\mu_s$,
the \engterm{angle of friction}. The admissible contact reactions
therefore sweep out a cone of half-angle $\phi_s$ about the normal: the
\engterm{friction cone}. A body stays put if and only if the line of
action of the required contact reaction lies inside the cone. This
reframes every block-on-incline problem geometrically: a block on a
plane inclined at $\theta$ needs a contact reaction tilted at $\theta$
from the normal to balance gravity, so it slides exactly when $\theta >
\phi_s$, recovering the threshold angle $\theta^\ast = \arctan\mu_s$
without writing a single component equation. The cone generalises to
three dimensions as a circular cone and underlies the analysis of
self-locking threads, wedges, and clamps in chapter 4 of this volume.
\end{mastery}

\begin{estimation}
\textbf{How heavy a wooden crate can a person push across a concrete
floor at steady speed?}

\smallskip
\textit{Estimate, before reading on.}

\smallskip
A working estimate goes through the following moves. A horizontal
push against kinetic friction balances $F_k = \mu_k N = \mu_k m g$.
For wood on dry concrete, $\mu_k$ is roughly $0.4$. A determined adult
can push with perhaps $300\,\text{N}$ horizontally for short
distances. Solving for the crate mass:
\[
m = \frac{F}{\mu_k g} \approx \frac{300}{0.4 \times 9.81}
\approx 76\,\si{\kilogram}.
\]
Estimate: a crate of roughly $50$ to $100\,\si{\kilogram}$ is at the
edge of what a single person can shove along.

A heavier crate requires either greater normal-force reduction (slide
the crate on rollers, or lubricate the floor) or a different
mechanical advantage (a hand truck, a lever, a winch). The
order-of-magnitude check explains why moving companies own equipment
rather than relying on muscle: muscle saturates at $\sim 80\,\si{
\kilogram}$ under steady push on a concrete floor, and a domestic
refrigerator weighs more.

The postmortem on the estimate: the $300\,\text{N}$ push is the
riskiest move. A trained mover can briefly produce more, perhaps
$500\,\text{N}$; over a sustained shove, less. The friction
coefficient is the second-riskiest move; a smoother floor can drop
$\mu_k$ to $0.3$, raising the answer to $\sim 100\,\si{\kilogram}$.
The structural answer, that a single human's pushing capacity is in
the range $50$--$100\,\si{\kilogram}$ on dry concrete, survives
either correction.
\end{estimation}

\begin{mastery}
\textbf{The resultant of a distributed load.} A real contact rarely acts
at a point. Weight is spread through a body, a fluid presses over an area,
a beam carries a load along its length. A distributed load is described by
an intensity $w(x)$, a force per unit length (or per unit area), and the
free-body diagram needs its single equivalent force: a resultant of the
right magnitude placed on the right line so that it produces the same force
sum and the same moment as the distribution. The magnitude is the integral
of the intensity,
\[
R = \int w(x)\,dx,
\]
the area under the load curve. The line of action passes through the
distribution's centroid, the point $\bar{x}$ that reproduces the moment,
\[
\bar{x} = \frac{\int x\,w(x)\,dx}{\int w(x)\,dx},
\]
the first moment of the load divided by its total. A uniform load $w_0$
over a length $L$ has resultant $w_0 L$ acting at the midpoint; a load that
ramps linearly from zero to $w_0$ over $L$ has resultant $\tfrac{1}{2}w_0 L$
acting two-thirds of the way along, at the centroid of the triangle. The
equivalence is exact for the external effect on a rigid body: the resultant
gives the same reactions at the supports as the full distribution. It is
not exact for the internal effect, the bending and shear along the member,
which depend on where the load actually sits; a beam loaded at its centre
and the same beam loaded uniformly share a support reaction but not a
deflection, the distinction Volume V's beam theory turns on. The resultant
is the tool for finding reactions; the distribution itself is kept for
anything that asks what happens inside the body.
\end{mastery}

\section{Worked examples}\pathtag{standard}

The worked diagrams below carry the FBD discipline through a range of
planar and spatial cases. A reader who follows the bookkeeping on them
should be able to construct an FBD for any rigid-body problem in this
volume. The first three are the canonical trio, a ladder, a truss joint,
and a hanging mass; the rest add a contact, a constraint, or a dimension
one at a time.

\subsection{A ladder against a wall}

% Figure: the ladder against a smooth wall, rough floor, with its FBD.
\begin{figure}[htbp]
\centering
\begin{tikzpicture}[
  force/.style={draw=engblue, -{Stealth}, very thick},
  rxn/.style={draw=engred, -{Stealth}, very thick},
  fric/.style={draw=engamber, -{Stealth}, very thick},
  scale=1.0,
]
% wall (vertical) and floor (horizontal)
\draw[draw=engblack, very thick] (0,0) -- (0,4.6);
\draw[draw=engblack, very thick] (0,0) -- (5.6,0);
% hatching on wall and floor
\foreach \y in {0.3,0.7,...,4.4} {\draw[draw=enggray, thin] (0,\y) -- (-0.28,\y-0.2);}
\foreach \x in {0.3,0.7,...,5.4} {\draw[draw=enggray, thin] (\x,0) -- (\x-0.2,-0.28);}
% the ladder: foot at (4.4,0), top at (0,3.8). angle theta with floor.
\coordinate (foot) at (4.4,0);
\coordinate (top) at (0,3.8);
\draw[draw=engblack, line width=2pt] (foot) -- (top);
% midpoint = centre of mass
\coordinate (mid) at (2.2,1.9);
\fill[engblack] (mid) circle (0.06);
% angle theta at foot
\draw[draw=enggray, thin] (3.7,0) arc (180:139:0.7);
\node[font=\scriptsize] at (3.5,0.3) {$\theta$};
% weight at midpoint
\draw[force] (mid) -- ++(0,-1.7) node[anchor=north, font=\small, engblue] {$mg$};
% wall normal: horizontal, away from wall (to the right) at top
\draw[rxn] (top) -- ++(1.5,0) node[anchor=west, font=\small, engred] {$N_W$};
% floor normal: up at foot
\draw[rxn] (foot) -- ++(0,1.6) node[anchor=south, font=\small, engred] {$N_F$};
% friction at foot toward wall (to the left)
\draw[fric] (foot) -- ++(-1.5,0) node[anchor=east, font=\small, engamber] {$F_f$};
% length label
\node[font=\scriptsize, anchor=south east, rotate=-41] at (1.1,2.85) {length $L$};
\end{tikzpicture}
\caption[Ladder against a smooth wall on a rough floor]{The chosen body
is the ladder. The smooth wall exerts a horizontal normal force $N_W$
only (no vertical friction). The rough floor exerts a vertical normal
$N_F$ and a horizontal friction $F_f$ directed toward the wall, opposing
the impending slip of the foot away from the wall. The weight $mg$ acts
at the midpoint. Taking moments about the foot eliminates $N_F$ and
$F_f$ at a stroke and isolates $N_W$, which is why the foot is the
moment centre of choice. The slip condition $F_f \le \mu_s N_F$ then
yields the minimum floor coefficient $\mu_s \ge 1/(2\tan\theta)$: steeper
ladders need less friction, and a shallow ladder can demand more
friction than wood on concrete provides.}
\label{fig:vol03ch01:ladder}
\end{figure}


A uniform ladder of length $L$ and mass $m$ leans against a smooth
vertical wall, with its base on a rough horizontal floor. The angle
between the ladder and the floor is $\theta$. Find the minimum
coefficient of static friction at the floor that keeps the ladder from
slipping.

\textbf{FBD.} The chosen body is the ladder. The forces:

\begin{itemize}[leftmargin=*]
  \item Weight $mg$, acting downward at the ladder's midpoint.
  \item Normal force $N_W$ from the wall, acting horizontally, away
    from the wall (perpendicular to the wall surface).
  \item Normal force $N_F$ from the floor, acting vertically upward.
  \item Friction force $F_f$ from the floor, acting horizontally,
    toward the wall (opposing the impending slip of the ladder's foot
    away from the wall).
\end{itemize}

The wall is smooth, so the wall exerts no tangential (vertical)
friction. The floor is rough, so the floor exerts both normal and
friction. Four unknowns ($N_W$, $N_F$, $F_f$, and the friction limit)
in three equations, with one inequality.

\textbf{Equilibrium.} Sum forces and moments. Taking moments about the
ladder's foot eliminates $N_F$ and $F_f$:
\begin{align*}
\sum F_x &= 0: & F_f - N_W &= 0, \\
\sum F_y &= 0: & N_F - mg &= 0, \\
\sum M_{\text{foot}} &= 0: & N_W L \sin\theta - mg \frac{L}{2}
\cos\theta &= 0.
\end{align*}
From the third equation, $N_W = \frac{mg}{2 \tan\theta}$. Then
$F_f = N_W = \frac{mg}{2 \tan\theta}$ and $N_F = mg$. The slip
condition is $F_f \le \mu_s N_F$, so
\[
\mu_s \ge \frac{1}{2 \tan\theta}.
\]
For $\theta = 60^{\circ}$, $\mu_s \ge 0.29$; for $\theta =
45^{\circ}$, $\mu_s \ge 0.5$; for $\theta = 30^{\circ}$, $\mu_s \ge
0.87$, which exceeds typical wood-on-concrete coefficients and means
the ladder will slip. A working painter knows the rule of thumb that
the base of a leaning ladder should be a quarter of its length from
the wall; that ratio corresponds to $\theta \approx 76^{\circ}$ and
$\mu_s \ge 0.12$, comfortably below any real surface coefficient.

A numerical instance closes the ladder example. Take $m = 12\,\si{
\kilogram}$, $L = 4\,\si{\meter}$, $\theta = 60^{\circ}$. Then $N_F =
mg = 12 \times 9.81 = 117.7\,\text{N}$, and $N_W = F_f = mg / (2\tan
60^{\circ}) = 117.7 / 3.464 = 34.0\,\text{N}$, so the required floor
coefficient is $\mu_s \ge F_f / N_F = 34.0/117.7 = 0.289$. The script
\texttt{code/resolve\_forces.py} assembles these four forces and
confirms that their vector sum and their net moment about the foot both
vanish, the numerical signature of a correct FBD. The table
\texttt{data/ladder\_slip\_angles.csv} tabulates the threshold across a
range of angles, the curve a ladder safety code is built on.

\subsection{A simple truss joint}

% Figure: method of joints, the FBD of a single pin loaded by an external
% force with three two-force members meeting at it.
\begin{figure}[htbp]
\centering
\begin{tikzpicture}[
  member/.style={draw=engblack, very thick},
  mforce/.style={draw=engred, -{Stealth}, very thick},
  load/.style={draw=engblue, -{Stealth}, very thick},
  scale=1.0,
]
% === Left: the joint in the truss ===
\begin{scope}[xshift=0cm]
  \coordinate (A) at (2.4,1.8);
  % three members radiating: AB up-left, AC down-right, AD straight down
  \draw[member] (A) -- ++(150:2.3) node[anchor=south east, font=\small] {$B$};
  \draw[member] (A) -- ++(-30:2.3) node[anchor=north west, font=\small] {$C$};
  \draw[member] (A) -- ++(-90:2.0) node[anchor=north, font=\small] {$D$};
  \draw[draw=engblack, fill=white, thick] (A) circle (0.12);
  \node[font=\small, anchor=south] at (2.4,2.0) {$A$};
  % external load P downward
  \draw[load] (A) -- ++(0,1.4) node[anchor=south, font=\small, engblue] {(joint $A$)};
  \node[font=\small, anchor=north] at (2.4,-1.0) {(a) joint in truss};
\end{scope}
% === Right: the FBD of the pin ===
\begin{scope}[xshift=7.4cm, yshift=0.4cm]
  \coordinate (A) at (0,0);
  \fill[engblack] (A) circle (0.07);
  % member forces along member axes, drawn as tension (away from joint)
  \draw[mforce] (A) -- ++(150:2.0) node[anchor=south east, font=\small, engred] {$F_{AB}$};
  \draw[mforce] (A) -- ++(-30:2.0) node[anchor=north west, font=\small, engred] {$F_{AC}$};
  \draw[mforce] (A) -- ++(-90:1.7) node[anchor=north, font=\small, engred] {$F_{AD}$};
  % external load P downward
  \draw[load] (A) -- ++(0,-1.7) node[anchor=north east, font=\small, engblue] {$P$};
  % angle annotations
  \draw[draw=enggray, thin] (0.7,0) arc (0:150:0.7);
  \node[font=\scriptsize] at (0.45,0.55) {$30^\circ$};
  \node[font=\small, anchor=north] at (0,-2.4) {(b) FBD of pin $A$};
\end{scope}
\end{tikzpicture}
\caption[Method of joints: the free-body diagram of a pin]{In a pin-
jointed truss every member is a two-force member, so its force lies
along its own axis. (a) A loaded joint $A$ where three members meet.
(b) Its free-body diagram treats the pin as a particle: the only forces
are the member forces along each member axis (drawn here as tension,
pulling away from the joint) and the external load $P$. A particle in
a plane gives two scalar equations $\sum F_x = 0$, $\sum F_y = 0$, so a
joint with two unknown member forces is solvable on its own; a joint
with three is closed only by stepping to a neighbouring joint. That
stepping, joint to joint, is the method of joints developed in chapter
2 of this volume. A member force that comes out negative was assumed in
the wrong sense: the member is in compression, not tension.}
\label{fig:vol03ch01:truss-joint}
\end{figure}


A pin joint in a planar truss connects three two-force members. The
joint is loaded by an external force $\vec{P}$ acting downward.
Members $AB$, $AC$, and $AD$ meet at the joint at angles $30^{\circ}$
above horizontal (member $AB$, going up and to the left),
$30^{\circ}$ below horizontal (member $AC$, going down and to the
right), and vertical (member $AD$, going straight down). Find the
forces in each member.

\textbf{FBD.} The chosen body is the joint itself, treated as a point.
The forces are the external load $\vec{P}$ downward and the three
member forces $F_{AB}$, $F_{AC}$, $F_{AD}$ along the lines of their
respective members. We assume tension positive (force pulling away
from the joint along the member's axis).

\textbf{Equilibrium.} Two scalar equations:
\begin{align*}
\sum F_x &= 0: & -F_{AB} \cos 30^{\circ} + F_{AC} \cos 30^{\circ} &=
0, \\
\sum F_y &= 0: & F_{AB} \sin 30^{\circ} - F_{AC} \sin 30^{\circ} -
F_{AD} - P &= 0.
\end{align*}
The first equation gives $F_{AB} = F_{AC}$. Substituting into the
second:
\[
2 F_{AB} \sin 30^{\circ} = F_{AD} + P, \quad F_{AB} = F_{AC} =
\frac{F_{AD} + P}{2 \sin 30^{\circ}} = F_{AD} + P.
\]
The joint equations alone leave one of $F_{AB}$, $F_{AC}$, $F_{AD}$
indeterminate; a second joint or a global truss equation is needed
to close the system. This is the structure of the method of joints:
each joint provides two equations in (typically) three unknowns, and
the truss is solved by stepping joint to joint. Chapter 2 develops
the full method.

\subsection{A mass hanging from two ropes}

% Figure: a mass on two ropes, the FBD as a concurrent force triangle,
% and the tension-divergence as the ropes flatten.
\begin{figure}[htbp]
\centering
\begin{tikzpicture}[
  rope/.style={draw=enggray, very thick},
  ten/.style={draw=engred, -{Stealth}, very thick},
  wt/.style={draw=engblue, -{Stealth}, very thick},
  scale=1.0,
]
% === Left: the scene ===
\begin{scope}[xshift=0cm]
  \draw[draw=engblack, very thick] (-0.3,3.6) -- (4.3,3.6);
  \foreach \x in {-0.1,0.3,...,4.1} {\draw[draw=enggray, thin] (\x,3.6) -- (\x-0.2,3.85);}
  \coordinate (P1) at (0.4,3.6);
  \coordinate (P2) at (3.6,3.6);
  \coordinate (M) at (2.0,1.4);
  \draw[rope] (P1) -- (M);
  \draw[rope] (P2) -- (M);
  \draw[draw=engblack, fill=engblue!12, thick] (1.6,0.7) rectangle (2.4,1.4);
  \node[font=\scriptsize] at (2.0,1.05) {$m$};
  % angles
  \node[font=\scriptsize] at (1.0,2.9) {$\alpha$};
  \node[font=\scriptsize] at (3.0,2.9) {$\beta$};
  \node[font=\small, anchor=north] at (2.0,0.3) {(a) scene};
\end{scope}
% === Right: concurrent force triangle ===
\begin{scope}[xshift=7.0cm, yshift=0.4cm]
  \coordinate (O) at (0,0);
  % weight straight down
  \draw[wt] (O) -- (0,-2.4) node[anchor=north, font=\small, engblue] {$mg$};
  % T1 up-left, T2 up-right (away from mass)
  \draw[ten] (O) -- (-1.7,2.2) node[anchor=south east, font=\small, engred] {$T_1$};
  \draw[ten] (O) -- (1.1,2.6) node[anchor=south west, font=\small, engred] {$T_2$};
  \node[font=\small, anchor=north] at (-0.3,-3.0) {(b) FBD of the mass};
\end{scope}
\end{tikzpicture}
\caption[A mass on two ropes and the divergence of tension]{(a) A mass
hangs from two ropes at angles $\alpha$ and $\beta$ to the vertical.
(b) Its free-body diagram is three concurrent forces: weight $mg$ down,
and tensions $T_1, T_2$ directed along each rope away from the mass.
Resolving gives $T_1 = mg\sin\beta / \sin(\alpha+\beta)$ and $T_2 =
mg\sin\alpha / \sin(\alpha+\beta)$. As the ropes flatten toward
horizontal, $\alpha+\beta \to \pi$ and $\sin(\alpha+\beta) \to 0$, so
both tensions diverge: two nearly horizontal ropes must carry tensions
far larger than the weight they support. The clothesline that barely
sags under a wet towel, and the rigging sling rated for a shallow
lifting angle, are the same fact seen twice.}
\label{fig:vol03ch01:hanging-mass}
\end{figure}


A point mass $m$ hangs in static equilibrium from two ropes attached
to the ceiling. The ropes make angles $\alpha$ and $\beta$ with the
vertical, with $\alpha + \beta < \pi$ (the mass hangs below both
attachment points). Find the tension in each rope.

\textbf{FBD.} The chosen body is the mass. Three forces: weight $mg$
downward, tension $T_1$ along the first rope (away from the mass,
upward and to one side), tension $T_2$ along the second rope (away
from the mass, upward and to the other side).

\textbf{Equilibrium.} Two equations:
\begin{align*}
\sum F_x &= 0: & T_1 \sin\alpha - T_2 \sin\beta &= 0, \\
\sum F_y &= 0: & T_1 \cos\alpha + T_2 \cos\beta - mg &= 0.
\end{align*}
Solving,
\[
T_1 = \frac{mg \sin\beta}{\sin(\alpha+\beta)}, \quad
T_2 = \frac{mg \sin\alpha}{\sin(\alpha+\beta)}.
\]
Two consequences worth noting. First, when $\alpha = \beta = 45^{\
circ}$, each rope carries $T = mg / \sqrt{2} \approx 0.71 mg$,
appreciably less than the full weight. Second, when the ropes
approach horizontal ($\alpha + \beta \to \pi$), the tension diverges:
two nearly horizontal ropes supporting a weight at their midpoint
must carry tensions much larger than the weight itself. A
clothesline supporting a wet towel demonstrates the divergence at
small angles to horizontal; a rigger lifting a load with a sling at
shallow angle demonstrates the same divergence with serious
consequences for the sling rating. The ratios for a range of angle
pairs are tabulated in \texttt{data/two\_rope\_tensions.csv}; at
$\alpha = \beta = 85^{\circ}$ each rope already carries more than five
times the weight.

\subsection{A parked car on a slope}

% Figure: a car parked on a slope, FBD with axle normals, weight, and
% the friction holding it. Used in the worked example and exercise 1.6.
\begin{figure}[htbp]
\centering
\begin{tikzpicture}[
  wt/.style={draw=engblue, -{Stealth}, very thick},
  nrm/.style={draw=engred, -{Stealth}, very thick},
  fric/.style={draw=engamber, -{Stealth}, very thick},
  scale=1.0,
]
% incline at ~12 deg for visibility
\begin{scope}[rotate=12]
  % road surface
  \draw[draw=engblack, very thick] (-0.6,0) -- (7.4,0);
  \foreach \x in {-0.3,0.2,...,7.2} {\draw[draw=enggray, thin] (\x,0) -- (\x-0.2,-0.25);}
  % car body (simplified)
  \draw[draw=engblack, fill=engblue!10, thick]
    (1.3,0.55) -- (2.0,1.25) -- (4.6,1.25) -- (5.4,0.55) -- cycle;
  \draw[draw=engblack, fill=white, thick] (2.2,0.55) circle (0.45);
  \draw[draw=engblack, fill=white, thick] (4.5,0.55) circle (0.45);
  \fill[engblack] (2.2,0.55) circle (0.05);
  \fill[engblack] (4.5,0.55) circle (0.05);
  % centre of mass
  \coordinate (G) at (3.6,1.0);
  \fill[engblack] (G) circle (0.06);
  \node[font=\scriptsize, anchor=south] at (G) {$G$};
  % normals at each axle (perpendicular to road = vertical in this frame)
  \draw[nrm] (2.2,0.1) -- ++(0,1.6) node[anchor=south, font=\scriptsize, engred] {$N_r$};
  \draw[nrm] (4.5,0.1) -- ++(0,1.6) node[anchor=south, font=\scriptsize, engred] {$N_f$};
  % friction up the slope
  \draw[fric] (2.2,0.1) -- ++(-1.6,0) node[anchor=east, font=\scriptsize, engamber] {$F_f$};
  % weight straight down (true vertical = -90-12 deg in the rotated frame)
  \draw[wt] (3.6,1.0) -- ++(-102:2.0) node[anchor=north, font=\scriptsize, engblue] {$mg$};
\end{scope}
% slope angle marker at origin
\draw[draw=enggray, thin] (1.2,0) arc (0:12:1.2);
\node[font=\scriptsize] at (1.45,0.12) {$\theta$};
\end{tikzpicture}
\caption[Free-body diagram of a car parked on a slope]{A parked car on
a slope of angle $\theta$. The chosen body is the whole car. Each axle
contact contributes a normal force perpendicular to the road, $N_f$ at
the front and $N_r$ at the rear; the braked or chocked wheels supply a
friction force $F_f$ up the slope; the weight $mg$ acts straight down at
the centre of mass $G$. Resolving along and perpendicular to the road
gives $F_f = mg\sin\theta$ and $N_f + N_r = mg\cos\theta$, while moments
about one axle split the normal load between front and rear according to
where $G$ sits in the wheelbase. The steeper the slope, the more of the
weight the friction must hold and the more the normal load shifts toward
the downhill axle.}
\label{fig:vol03ch01:parked-car}
\end{figure}


A car of mass $m$ is parked, brakes set, on a slope of angle $\theta$.
The wheelbase is $\ell$, and the centre of mass sits a distance $a$
behind the front axle and a height $h$ above the road. Find the normal
force at each axle and the friction the brakes must supply.

\textbf{FBD.} The chosen body is the whole car. Forces: the weight $mg$
down at the centre of mass $G$; a normal force at the front axle $N_f$
and at the rear axle $N_r$, both perpendicular to the road; and a total
friction force $F_f$ up the slope at the contact patches (the brakes
lock the wheels, so friction is static). We resolve along the road
($x$, up-slope positive) and perpendicular to it ($y$).

\textbf{Equilibrium.}
\begin{align*}
\sum F_x &= 0: & F_f - mg\sin\theta &= 0, \\
\sum F_y &= 0: & N_f + N_r - mg\cos\theta &= 0, \\
\sum M_{\text{front axle}} &= 0: & N_r \ell - mg\cos\theta\, a -
mg\sin\theta\, h &= 0.
\end{align*}
The friction follows at once: $F_f = mg\sin\theta$. The moment equation
gives the rear normal, $N_r = mg(a\cos\theta + h\sin\theta)/\ell$, and
the force balance then gives $N_f = mg\cos\theta - N_r$. The
$h\sin\theta$ term is the load transfer: on a slope the weight's line of
action shifts toward the downhill axle, off-loading the uphill wheels.
For a $1{,}500\,\si{\kilogram}$ car with $\ell = 2.6\,\si{\meter}$,
$a = 1.2\,\si{\meter}$, $h = 0.5\,\si{\meter}$, parked nose-up on a
$10^{\circ}$ slope, $F_f = 1500 \times 9.81 \times \sin 10^{\circ}
\approx 2{,}555\,\text{N}$, $N_r \approx 7{,}180\,\text{N}$, and
$N_f \approx 7{,}312\,\text{N}$. The two normals are nearly equal
because $G$ sits near mid-wheelbase; the small height term $h\sin\theta$
nudges load to the rear. The brakes carry the whole $2.56\,\text{kN}$
of down-slope pull, which is why a parking pawl or a wheel chock is the
prudent backup on a steep grade.

\subsection{An ideal pulley and the load on its axle}

% Figure: an ideal pulley redirecting a rope, with the FBD of the pulley
% axle showing the resultant axle force.
\begin{figure}[htbp]
\centering
\begin{tikzpicture}[
  rope/.style={draw=enggray, very thick},
  ten/.style={draw=engred, -{Stealth}, very thick},
  res/.style={draw=engblue, -{Stealth}, very thick},
  scale=1.0,
]
% === Left: the scene, pulley redirecting horizontal rope to vertical ===
\begin{scope}[xshift=0cm]
  \coordinate (P) at (2.0,2.6);
  \draw[draw=engblack, fill=enggray!12, thick] (P) circle (0.5);
  \fill[engblack] (P) circle (0.06);
  % horizontal rope coming from the left, wrapping over the top
  \draw[rope] (-0.6,3.1) -- (2.0,3.1);
  % vertical rope down the right side to a mass
  \draw[rope] (2.5,2.6) -- (2.5,0.9);
  \draw[draw=engblack, fill=engblue!12, thick] (2.15,0.2) rectangle (2.85,0.9);
  \node[font=\scriptsize] at (2.5,0.55) {$m$};
  % support bracket to axle
  \draw[draw=engblack, thick] (2.0,2.6) -- (2.0,4.2);
  \draw[draw=engblack, very thick] (1.3,4.2) -- (2.7,4.2);
  \foreach \x in {1.4,1.7,...,2.6} {\draw[draw=enggray, thin] (\x,4.2) -- (\x-0.2,4.45);}
  \node[font=\small, anchor=north] at (1.4,-0.1) {(a) scene};
\end{scope}
% === Right: FBD of the pulley (axle) ===
\begin{scope}[xshift=7.2cm, yshift=1.4cm]
  \coordinate (O) at (0,0);
  \draw[draw=engblack, fill=enggray!8, thick] (O) circle (0.5);
  \fill[engblack] (O) circle (0.06);
  % tension pulling left (horizontal segment) and down (vertical segment)
  \draw[ten] (0,0.5) -- (-2.0,0.5) node[anchor=south, font=\small, engred] {$T$};
  \draw[ten] (0.5,0) -- (0.5,-2.0) node[anchor=west, font=\small, engred] {$T$};
  % resultant rope pull is down-left; axle reaction balances it (up-right)
  \draw[res] (O) -- (1.6,1.6) node[anchor=south west, font=\small, engblue] {$R$ (axle)};
  \node[font=\small, anchor=north] at (-0.2,-2.6) {(b) FBD of pulley};
\end{scope}
\end{tikzpicture}
\caption[An ideal pulley and the force on its axle]{An ideal (massless,
frictionless) pulley redirects a rope without changing its tension, so
the rope tension $T$ is the same on both sides. (a) A horizontal rope
is turned to vertical and supports a hanging mass, $T = mg$. (b) The
free-body diagram of the pulley shows the rope pulling with tension $T$
along each departing segment; the axle reaction $R$ balances their
vector sum. For a right-angle turn the two equal tensions are
perpendicular, so the axle carries $R = \sqrt{T^2 + T^2} = T\sqrt{2}
\approx 1.41\,T$, directed along the bisector of the rope angle. The
ideal pulley changes a rope's direction for free, but the axle still
pays the vector cost.}
\label{fig:vol03ch01:pulley}
\end{figure}


A frictionless, massless pulley does not change a rope's tension; it
changes the rope's direction. The two facts together fix the force on
the pulley's axle. A horizontal rope turns over the top of a pulley and
descends vertically to support a mass $m$, as in
figure~\ref{fig:vol03ch01:pulley}. The vertical segment carries
$T = mg$, and because the pulley is ideal the horizontal segment carries
the same $T$. The FBD of the pulley shows two equal tensions, one
pulling left and one pulling down, plus the axle reaction. Summing
forces, the axle must supply $\vec{R}$ equal and opposite to the rope
pull: $R_x = T$, $R_y = T$, so $R = \sqrt{T^2 + T^2} = T\sqrt{2} \approx
1.41\,mg$, directed up and to the right along the bisector of the
right-angle turn. The lesson recurs across the volume: redirecting a
rope is free in tension but not free in support. Every direction change
a rope makes is paid for at the device that makes it, and that device's
mounting must carry the vector resultant, not the rope tension.

\subsection{A three-force member and the concurrence theorem}

% Figure: the three-force-member theorem. A bent bracket loaded at one
% pin and supported at two others; the three reaction lines meet at one
% point, the concurrence that closes the problem geometrically.
\begin{figure}[htbp]
\centering
\begin{tikzpicture}[
  member/.style={draw=engblack, line width=1.6pt},
  line/.style={draw=enggray, thin, dashed},
  force/.style={draw=engred, -{Stealth}, very thick},
  scale=1.0,
]
% the bent member: pin at A (lower left), roller at B (lower right),
% applied load at the top corner D
\coordinate (A) at (0,0);
\coordinate (B) at (4.4,0);
\coordinate (D) at (1.6,3.0);
% body outline (a triangle-ish bracket)
\draw[member] (A) -- (B) -- (D) -- cycle;
% pins/rollers
\draw[draw=engblack, fill=white] (A) circle (0.13);
\fill[engblack] (A) circle (0.06);
\draw[draw=engblack, fill=white] (B) circle (0.13);
\fill[engblack] (B) circle (0.06);
\node[font=\scriptsize, anchor=north east] at (A) {$A$};
\node[font=\scriptsize, anchor=north west] at (B) {$B$};
\node[font=\scriptsize, anchor=south] at (D) {$D$};
% roller symbol under B
\draw[draw=engblack] (4.15,-0.13) -- (4.65,-0.13);
\draw[draw=engblack, fill=white] (4.22,-0.22) circle (0.07);
\draw[draw=engblack, fill=white] (4.58,-0.22) circle (0.07);
% applied load P at D, downward and to the right
\draw[force] (D) -- ++(1.2,-1.2) node[anchor=west, font=\small, engred] {$P$};
% the roller reaction at B is vertical (normal to the ground)
% extend its line of action upward to the concurrence point
\coordinate (X) at (4.4,2.05); % concurrence point (on B's vertical line)
\draw[line] (B) -- (X);
% the load line of action through D, extended to X
\draw[line] (D) -- (X);
% the pin reaction at A must pass through X too: draw that line
\draw[line] (A) -- (X);
% mark concurrence
\fill[engblue] (X) circle (0.07);
\node[font=\scriptsize, anchor=west, engblue] at (X) {concurrence};
% reaction arrows (schematic) along the deduced lines
\draw[draw=engblue, -{Stealth}, very thick] (B) -- ++(0,1.3)
  node[anchor=west, font=\scriptsize, engblue] {$R_B$};
\draw[draw=engblue, -{Stealth}, very thick] (A) -- ($(A)!0.45!(X)$)
  node[anchor=south east, font=\scriptsize, engblue] {$R_A$};
\end{tikzpicture}
\caption[The three-force member: lines of action meet at a point]{A rigid
body acted on by exactly three non-parallel forces is in equilibrium only
if their three lines of action are concurrent. Here a bracket carries an
applied load $P$ at $D$, a roller reaction $R_B$ at $B$ (whose direction
is fixed vertical by the roller), and a pin reaction $R_A$ at $A$ (whose
direction is unknown). The load line and the roller line meet at one
point; moment balance about that point forces the pin reaction to pass
through it as well, which fixes the otherwise-unknown direction of $R_A$.
The magnitudes then follow from a single force triangle. Concurrence turns
a three-equation problem into one drawing.}
\label{fig:vol03ch01:three-force}
\end{figure}


A rigid body acted on by exactly three non-parallel forces obeys a
geometric constraint that often solves the problem faster than the three
scalar equations. The constraint is concurrence: the three lines of
action meet at a single point. The proof is one moment equation. Pick the
intersection point $X$ of any two of the forces' lines of action. Both of
those forces have zero moment about $X$, because each line passes through
$X$. Equilibrium demands $\sum M_X = 0$, so the third force must also have
zero moment about $X$, which forces its line of action through $X$ as
well. Three lines, one point.

The bracket of figure~\ref{fig:vol03ch01:three-force} shows the payoff. A
bent bracket is pinned at $A$, supported by a roller at $B$, and loaded by
a known force $P$ at $D$. The roller fixes the direction of its reaction
$R_B$ (vertical, normal to the rolling surface) but not its magnitude. The
pin at $A$ fixes neither the direction nor the magnitude of $R_A$. Three
unknowns, as it should be for a determinate planar body: $R_B$, and the
two components of $R_A$. Rather than write three equations, extend the
known line of $P$ and the known line of $R_B$ until they cross, at the
point $X$. Concurrence then says $R_A$ acts along the line $AX$, which
fixes the one remaining direction. The three magnitudes follow from a
single closed force triangle: $P$, $R_A$, and $R_B$ drawn tip to tail
must return to the start, because their vector sum is zero. Take $P =
500\,\text{N}$ acting at $45^{\circ}$ below the horizontal at $D$, with
the geometry of the figure ($A$ at the origin, $B$ at $(4.4, 0)\,\si{
\meter}$, $D$ at $(1.6, 3.0)\,\si{\meter}$). The line of $P$ through $D$
meets the vertical through $B$ at $X = (4.4, 0.2)\,\si{\meter}$, so $R_A$
lies along $AX$, at $\arctan(0.2/4.4) = 2.6^{\circ}$ above the horizontal.
Resolving the force triangle, $R_A \approx 354\,\text{N}$ and $R_B \approx
354\,\text{N}$ for this near-symmetric placement. The lesson carries past
this one bracket: whenever a body sees three forces, look for the
concurrence before reaching for components, because a crossing of two
lines is faster to draw than a system of equations is to solve, and the
drawing exposes a sign error that the algebra hides.

\subsection{A wedge at limiting equilibrium}

% Figure: a friction wedge driven under a load, with the FBD of the wedge
% showing the two contact reactions and the driving force at the verge of
% slipping. Supports the limiting-equilibrium worked example.
\begin{figure}[htbp]
\centering
\begin{tikzpicture}[
  force/.style={draw=engblue, -{Stealth}, very thick},
  nrm/.style={draw=engred, -{Stealth}, very thick},
  fric/.style={draw=engamber, -{Stealth}, very thick},
  scale=1.0,
]
% ===== Left: the scene =====
\begin{scope}[xshift=0cm]
  % ground
  \draw[draw=engblack, very thick] (-0.4,0) -- (5.0,0);
  \foreach \x in {-0.2,0.2,...,4.8} {\draw[draw=enggray, thin] (\x,0) -- (\x-0.2,-0.25);}
  % the block being lifted (rests on top of the wedge incline)
  \draw[draw=engblack, fill=engblue!10, thick] (2.4,0.9) rectangle (4.6,2.1);
  \node[font=\scriptsize] at (3.5,1.5) {load $W$};
  % the wedge: thin triangle, apex angle alpha, driven to the right under
  % the block. base on the ground, top face under the block.
  \draw[draw=engblack, fill=engamber!15, thick]
    (0.2,0) -- (4.6,0) -- (4.6,0.9) -- cycle;
  \node[font=\scriptsize] at (3.0,0.32) {wedge};
  % apex angle alpha at the thin (left) end
  \draw[draw=enggray, thin] (0.9,0) arc (0:11:0.7);
  \node[font=\scriptsize, anchor=west] at (0.95,0.08) {$\alpha$};
  % driving force P on the wedge, horizontal to the right
  \draw[force] (0.2,0.18) -- ++(-1.2,0)
    node[anchor=east, font=\small, engblue] {$P$};
  \node[font=\small, anchor=north] at (2.2,-0.45) {(a) the scene};
\end{scope}
% ===== Right: FBD of the wedge alone =====
\begin{scope}[xshift=7.0cm]
  % wedge outline repeated
  \draw[draw=engblack, fill=engamber!15, thick]
    (0.2,0) -- (4.6,0) -- (4.6,0.9) -- cycle;
  % driving force
  \draw[force] (0.2,0.18) -- ++(-1.3,0)
    node[anchor=east, font=\scriptsize, engblue] {$P$};
  % bottom contact (with ground): normal up, friction left (opposing the
  % wedge's rightward motion)
  \draw[nrm] (2.6,0) -- ++(0,1.3) node[anchor=south, font=\scriptsize, engred] {$N_1$};
  \draw[fric] (2.6,0) -- ++(-1.3,0) node[anchor=south, font=\scriptsize, engamber] {$F_1$};
  % top contact (with the load): normal perpendicular to the inclined top
  % face, pointing down-and-left into the wedge; friction along that face
  % opposing the wedge sliding under the block (down-the-incline on wedge)
  \draw[nrm] (3.6,0.74) -- ++(-0.5,-1.25)
    node[anchor=north, font=\scriptsize, engred] {$N_2$};
  \draw[fric] (3.6,0.74) -- ++(-1.25,-0.25)
    node[anchor=east, font=\scriptsize, engamber] {$F_2$};
  \node[font=\small, anchor=north] at (2.2,-0.45) {(b) FBD of the wedge};
\end{scope}
\end{tikzpicture}
\caption[A friction wedge at limiting equilibrium]{A wedge of apex angle
$\alpha$ driven horizontally under a load $W$. (a) The scene: the driving
force $P$ pushes the wedge to the right, lifting the load along the
wedge's top face. (b) The free-body diagram of the wedge has two
contacts, each contributing a normal force and a friction force. At the
ground the normal $N_1$ is vertical and the friction $F_1$ opposes the
rightward motion. At the top face the normal $N_2$ is perpendicular to the
incline and the friction $F_2$ acts along it, both opposing the wedge
sliding under the block. At the verge of motion each friction takes its
limiting value $\mu N$, and the three equilibrium equations close on the
driving force $P$ needed to advance the wedge.}
\label{fig:vol03ch01:wedge}
\end{figure}


The friction examples so far balanced a single contact. A wedge balances
two, and the two contacts couple through the wedge's geometry, which is
what makes the wedge a force multiplier and a self-locking device at once.
A wedge of apex angle $\alpha$ is driven horizontally under a load $W$, as
in figure~\ref{fig:vol03ch01:wedge}. Find the horizontal force $P$ that
just advances the wedge, taking the same friction coefficient $\mu$ at
both contacts.

\textbf{FBD.} The chosen body is the wedge alone. Two contacts, so four
contact forces. At the ground: a vertical normal $N_1$ and a horizontal
friction $F_1$ opposing the rightward motion. At the top face (inclined at
$\alpha$): a normal $N_2$ perpendicular to the face and a friction $F_2$
along the face, both opposing the wedge sliding under the block. The
driving force $P$ is horizontal. At the verge of motion both frictions
reach their limit, $F_1 = \mu N_1$ and $F_2 = \mu N_2$.

\textbf{The load contact first.} The block sits on the wedge's inclined
top face and presses down with $W$ (the block's own equilibrium, with
smooth vertical guides assumed so the block moves only vertically). The
face reaction on the block is $N_2$ normal plus $\mu N_2$ along the face.
Vertical balance of the block gives the relation between $N_2$ and $W$;
for a small apex angle the leading result is $N_2 \approx W / (\cos\alpha
- \mu\sin\alpha)$.

\textbf{Equilibrium of the wedge.} Resolve horizontally and vertically. The
top-face normal $N_2$ tilts at $\alpha$ from vertical, so it contributes a
horizontal component $N_2\sin\alpha$ resisting $P$ and the top-face
friction $\mu N_2$ contributes a horizontal component $\mu N_2\cos\alpha$
also resisting $P$. The ground friction adds $\mu N_1$. Horizontal
balance:
\[
P = N_2\sin\alpha + \mu N_2\cos\alpha + \mu N_1.
\]
Vertical balance of the wedge gives $N_1 = N_2\cos\alpha - \mu N_2
\sin\alpha$ (the wedge's own weight neglected). Substituting and
collecting,
\[
P = N_2\big[\sin\alpha + \mu\cos\alpha + \mu(\cos\alpha -
\mu\sin\alpha)\big]
= N_2\big[\sin\alpha(1 - \mu^2) + 2\mu\cos\alpha\big].
\]
A numerical instance closes it. Take $W = 5{,}000\,\text{N}$, $\alpha =
8^{\circ}$, $\mu = 0.2$. Then $N_2 = 5000/(\cos 8^{\circ} - 0.2\sin
8^{\circ}) = 5000/(0.990 - 0.028) = 5{,}198\,\text{N}$, and
\[
P = 5198\big[\sin 8^{\circ}(1 - 0.04) + 0.4\cos 8^{\circ}\big]
= 5198\,[0.1337 + 0.3961] = 2{,}754\,\text{N}.
\]
The wedge lifts a $5{,}000\,\text{N}$ load with a $2{,}754\,\text{N}$
push, a mechanical advantage of $1.8$, and the advantage climbs sharply as
$\alpha$ shrinks. The self-locking question is the mirror image: remove
$P$ and ask whether the load drives the wedge back out. The wedge holds
without the push exactly when the friction available exceeds the back-drive
component, which for equal coefficients reduces to $\tan\alpha < 2\mu/(1 -
\mu^2)$, satisfied here ($\tan 8^{\circ} = 0.14$ against $0.42$) so the
wedge stays put when released. The wedge that lifts and then locks is the
same device a door wedge, a splitting maul, and a machine-tool clamp all
rely on, and the locking condition is a friction inequality, not a
strength one.

\subsection{Concurrent forces in three dimensions: a tripod load}

The planar examples used two force equations and one moment equation. A
concurrent three-dimensional system uses three force equations and no
moment equation, because all forces pass through one point and so have no
moment about it. A camera tripod holds a head at the apex $P =
(0, 0, 1.5)\,\si{\meter}$ above a level floor; three legs run from $P$ to
feet at $A = (0.8, 0, 0)$, $B = (-0.4, 0.69, 0)$, and $C = (-0.4, -0.69,
0)\,\si{\meter}$, spaced $120^{\circ}$ apart on a circle of radius
$0.8\,\si{\meter}$. A vertical load $W = 60\,\text{N}$ (camera plus head)
presses down at $P$. Find the axial force in each leg, taking the legs as
two-force members.

\textbf{FBD.} The chosen body is the apex joint $P$, a point. Four forces
meet there: the load $W = (0, 0, -60)\,\text{N}$ and the three leg forces,
each directed from $P$ toward its foot (compression, since the legs push
up on the joint) along the unit vectors
\[
\hat{u}_A = \frac{A - P}{|A - P|}, \quad
\hat{u}_B = \frac{B - P}{|B - P|}, \quad
\hat{u}_C = \frac{C - P}{|C - P|}.
\]
Each leg has length $|A - P| = \sqrt{0.8^2 + 0 + 1.5^2} = \sqrt{0.64 +
2.25} = \sqrt{2.89} = 1.70\,\si{\meter}$, the same for all three by
symmetry. So $\hat{u}_A = (0.8, 0, -1.5)/1.70 = (0.471, 0, -0.883)$, and
likewise for the other two with their own foot coordinates.

\textbf{Equilibrium.} With $F_A, F_B, F_C$ the (signed) axial magnitudes,
force balance is the single vector equation $F_A\hat{u}_A + F_B\hat{u}_B +
F_C\hat{u}_C + W = \vec{0}$, three scalar equations. Symmetry collapses
them: the three legs are identical and the load is vertical and central, so
$F_A = F_B = F_C = F$, and the horizontal components cancel by the
$120^{\circ}$ spacing. The vertical equation is all that remains:
\[
3F\,(-0.883) + (-60) = 0
\quad\Longrightarrow\quad
F = \frac{-60}{-2.648} = 22.7\,\text{N} \;\text{per leg}.
\]
The sign convention makes $F$ the magnitude of the force the leg exerts on
the joint along $\hat{u}_A$ (which points down-and-out), so the joint feels
each leg pushing up: the legs are in compression at $22.7\,\text{N}$ each.
The check is immediate: three legs each supplying $22.7 \times 0.883 =
20.0\,\text{N}$ vertical sum to $60\,\text{N}$, balancing $W$. Spreading
the legs wider (larger foot radius) lengthens them, lowers the vertical
fraction $0.883$ of each leg force, and raises the compression each leg
carries, the three-dimensional echo of the shallow-rope divergence of
section 1.6.3. A tripod set up too splayed loads its legs harder, not
less, the opposite of a beginner's intuition.

\subsection{A cable with a central point load}

The hanging-mass example fixed the two ropes at given angles. A cable
strung between two level supports finds its own angle once a weight is hung
on it, and the angle is set by how far the cable is allowed to sag. A light
cable of horizontal span $2\ell$ runs between two supports at the same
height and carries a single weight $W$ at its midpoint. The cable sags a
depth $h$ below the supports at that point. Find the tension in the cable
and the force the supports must resist.

\textbf{FBD.} The chosen body is the midpoint node where the weight hangs.
Three forces meet there: the weight $W$ straight down, and the two cable
tensions $T$, each directed up along its half of the cable toward a support.
By symmetry the two tensions are equal in magnitude. Each makes an angle
$\phi$ below the horizontal at the node, with $\tan\phi = h/\ell$, the sag
over the half-span.

\textbf{Equilibrium.} The horizontal components of the two tensions cancel
by symmetry, leaving the vertical balance:
\[
2T\sin\phi = W \quad\Longrightarrow\quad T = \frac{W}{2\sin\phi}.
\]
The support must resist both a vertical pull $T\sin\phi = W/2$ (half the
weight, as it must be) and a horizontal pull $H = T\cos\phi$, the horizontal
component that the two cable halves transmit outward to the supports. That
horizontal pull is the part a beginner forgets. Take $W = 1{,}000\,\text{N}$,
half-span $\ell = 5\,\si{\meter}$, sag $h = 0.5\,\si{\meter}$. Then $\tan\phi
= 0.5/5 = 0.1$, so $\phi = 5.71^{\circ}$, $\sin\phi = 0.0995$, and
\[
T = \frac{1000}{2 \times 0.0995} = 5{,}025\,\text{N},
\]
five times the hung weight. The horizontal support pull is $H = T\cos\phi =
5025 \times 0.995 = 5{,}000\,\text{N}$ at each end, drawn together by the
fact that $H = W\ell/(2h)$ exactly: the horizontal thrust scales inversely
with the sag. Allow only half the sag, $h = 0.25\,\si{\meter}$, and the
tension and the horizontal thrust both double. The lesson recurs wherever a
load hangs from a near-level line: a flat cable buys its flatness with
tension and with a horizontal force that the anchorages, not just the cable,
must carry. A guy wire, a power line, a tightrope, and a suspension-bridge
hanger are the same equation read at different scales.

\subsection{A body resting in a symmetric V-groove}

A cylinder dropped into a V-groove touches two surfaces at once, and the
two contacts share the weight in a ratio the groove angle sets. A uniform
cylinder of weight $W$ rests in a smooth V-groove whose two faces each make
an angle $\psi$ with the horizontal, symmetric about the vertical. Find the
normal force on each face.

\textbf{FBD.} The chosen body is the cylinder. The contacts are smooth, so
each face supplies a normal force only, directed along its own surface
normal toward the cylinder's centre. Call them $N_1$ and $N_2$. By symmetry
$N_1 = N_2 = N$. The third force is the weight $W$ down through the centre.
Each contact normal is perpendicular to its face. With the face at $\psi$ to
the horizontal, the inward normal rises at $\psi$ from the horizontal toward
the centre, so its vertical component is $N\sin\psi$ and its horizontal
component is $N\cos\psi$.

\textbf{Equilibrium.} The horizontal components of the two normals cancel by
symmetry. The vertical balance is
\[
2N\sin\psi = W \quad\Longrightarrow\quad N = \frac{W}{2\sin\psi}.
\]
A wide groove ($\psi$ near $90^{\circ}$, faces nearly vertical, a deep
narrow vee) presses the cylinder lightly against each face, $N \to W/2$. A
shallow groove ($\psi$ small, faces nearly flat) drives the normals up
without bound, $N \to \infty$, because two nearly horizontal faces must
generate a large normal to find any vertical component at all. For a
$300\,\text{N}$ shaft in a $60^{\circ}$ vee (each face $60^{\circ}$ from
horizontal, the common machinist's groove), $N = 300/(2\sin 60^{\circ}) =
300/1.732 = 173\,\text{N}$ on each face, more than half the weight apiece,
the price the groove pays for locating the shaft against sideways motion as
well as supporting it. The V-block on a surface plate and the prism that
cradles a rotating shaft both live on this single equation.

\subsection{Three concurrent cables in a plane}

A ring loaded by three cables in a plane is the smallest problem that
forces the engineer to balance two non-trivial force equations at once,
and it is the model for any fitting where several lines pull on one point.
A small steel ring is held in equilibrium by three cables lying in a
vertical plane. Cable $1$ pulls at $T_1 = 600\,\text{N}$ along a direction
$20^{\circ}$ above the horizontal to the right. Cable $2$ pulls at an
unknown tension $T_2$ straight up. Cable $3$ pulls at an unknown tension
$T_3$ along a direction $50^{\circ}$ below the horizontal to the left.
Find $T_2$ and $T_3$.

\textbf{FBD.} The chosen body is the ring, treated as a point. Three
forces meet there, all tensions pulling away from the ring along their
own cables: $T_1$ up-and-right, $T_2$ straight up, and $T_3$ down-and-left.
A point in equilibrium under three forces has no moment equation to write,
because every line of action passes through the point; the two force
equations carry the whole problem.

\textbf{Equilibrium.} Resolve into horizontal and vertical components,
counting rightward and upward as positive. Cable $3$ points down and to
the left, so both of its components are negative:
\begin{align*}
\sum F_x &= 0: & T_1\cos 20^{\circ} - T_3\cos 50^{\circ} &= 0, \\
\sum F_y &= 0: & T_1\sin 20^{\circ} + T_2 - T_3\sin 50^{\circ} &= 0.
\end{align*}
The horizontal equation has only one unknown and gives $T_3$ at once:
\[
T_3 = \frac{T_1\cos 20^{\circ}}{\cos 50^{\circ}}
= \frac{600 \times 0.940}{0.643} = 877\,\text{N}.
\]
Substituting into the vertical equation returns $T_2$:
\[
T_2 = T_3\sin 50^{\circ} - T_1\sin 20^{\circ}
= 877 \times 0.766 - 600 \times 0.342 = 672 - 205 = 467\,\text{N}.
\]
The check is the closure of the force polygon: the three tension vectors,
laid tip to tail, must return to the start, and they do to within rounding.
The lesson is the order of solution. A two-equation system is solved
fastest by finding an equation with a single unknown and clearing it first,
which here the horizontal balance supplies because cable $2$ is vertical
and contributes nothing to it. Choosing axes so that one unknown drops out
of one equation is the planar twin of summing moments about a point a force
passes through, and it is the habit that keeps a two-cable or three-cable
fitting from turning into a simultaneous-equation grind.

\subsection{A boom held by a cable, with a wall pin}

% Figure: a boom hinged to a wall at its foot and held by an inclined
% cable to a higher wall point, carrying a load at its tip. Three forces
% act on the boom: the cable tension, the wall-pin reaction, and the load.
% The figure shows the assembly and the boom FBD with the pin reaction
% direction left as an unknown to be found by concurrence.
\begin{figure}[htbp]
\centering
\begin{tikzpicture}[
  member/.style={draw=engblack, line width=1.6pt},
  cable/.style={draw=engblue, line width=1.3pt},
  load/.style={draw=engred, -{Stealth}, very thick},
  react/.style={draw=engamber, -{Stealth}, very thick},
  scale=1.0,
]
% ===== Left: the assembly =====
\begin{scope}[xshift=0cm]
  % wall
  \draw[draw=engblack, very thick] (0,-0.2) -- (0,3.6);
  \foreach \y in {0.1,0.5,...,3.4} {\draw[draw=enggray, thin] (0,\y) -- (-0.28,\y-0.2);}
  \coordinate (O) at (0,0.6);        % pin at foot
  \coordinate (Tip) at (3.4,1.4);    % boom tip, boom inclined up
  \coordinate (Cw) at (0,3.0);       % cable anchor high on wall
  % boom
  \draw[member] (O) -- (Tip);
  % cable from high wall point to tip
  \draw[cable] (Cw) -- (Tip);
  % pin marker at O
  \fill[engblack] (O) circle (0.07); \draw[draw=engblack, fill=white] (O) circle (0.12);
  % cable anchor marker
  \fill[engblack] (Cw) circle (0.07); \draw[draw=engblack, fill=white] (Cw) circle (0.12);
  % tip pin
  \fill[engblack] (Tip) circle (0.06);
  \node[font=\scriptsize, anchor=east] at (O) {$A$};
  \node[font=\scriptsize, anchor=east] at (Cw) {$C$};
  \node[font=\scriptsize, anchor=south west] at (Tip) {$B$};
  % load at tip
  \draw[load] (Tip) -- ++(0,-1.4) node[anchor=north, font=\small, engred] {$W$};
  \node[font=\small, anchor=north] at (1.8,-0.55) {(a) the assembly};
\end{scope}
% ===== Right: FBD of the boom with the three forces =====
\begin{scope}[xshift=6.4cm]
  \coordinate (O2) at (0,0.6);
  \coordinate (Tip2) at (3.4,1.4);
  \draw[member] (O2) -- (Tip2);
  \fill[engblack] (O2) circle (0.06);
  \fill[engblack] (Tip2) circle (0.06);
  \node[font=\scriptsize, anchor=north east] at (O2) {$A$};
  \node[font=\scriptsize, anchor=south west] at (Tip2) {$B$};
  % cable tension at tip: up and toward the wall along BC
  \draw[cable, -{Stealth}, very thick] (Tip2) -- ++(-1.7,0.8)
    node[anchor=south, font=\scriptsize, engblue] {$T$};
  % load down at tip
  \draw[load] (Tip2) -- ++(0,-1.4) node[anchor=north, font=\small, engred] {$W$};
  % pin reaction at A, direction unknown (drawn provisionally)
  \draw[react] (O2) -- ++(1.3,0.7)
    node[anchor=south west, font=\scriptsize, engamber] {$R_A$};
  \node[font=\scriptsize, anchor=north, engamber] at (0.7,0.55) {(direction by concurrence)};
  \node[font=\small, anchor=north] at (1.8,-0.55) {(b) FBD of the boom};
\end{scope}
\end{tikzpicture}
\caption[A boom held by a cable with a wall pin]{A boom hinged to a wall
at its foot $A$, held up by a cable from a higher wall point $C$ to the
boom tip $B$, with a load $W$ hanging at $B$. (a) The assembly. (b) The
free-body diagram of the boom. Exactly three forces act on it: the cable
tension $T$ along $BC$, the hanging load $W$ down at $B$, and the pin
reaction $R_A$ at the foot. Because the cable tension and the load both
pass through $B$, their lines of action cross there, so concurrence forces
the pin reaction to pass through $B$ as well: $R_A$ acts along $AB$,
fixing the one direction the pin would otherwise leave open. The three
magnitudes then close a single force triangle.}
\label{fig:vol03ch01:boom-wall-pin}
\end{figure}


The three-force concurrence theorem of section 1.6.6 earns its keep on a
boom. A uniform boom $AB$ of negligible weight is hinged to a wall at its
foot $A$ and held up by a cable running from a higher wall point $C$ to the
boom tip $B$, where a load $W = 2{,}000\,\text{N}$ hangs, as in
figure~\ref{fig:vol03ch01:boom-wall-pin}. The boom rises at $25^{\circ}$
above the horizontal, and the cable $BC$ makes an angle $\gamma =
40^{\circ}$ with the boom. Find the cable tension and the pin reaction at
$A$.

\textbf{FBD.} The chosen body is the boom. Exactly three forces act on it:
the cable tension $T$ at $B$ directed along $BC$, the load $W$ down at $B$,
and the pin reaction $R_A$ at the foot. The cable tension and the load both
pass through the tip $B$, so by concurrence the third force, the pin
reaction, must also pass through $B$: $R_A$ acts along the line $AB$, that
is along the boom at $25^{\circ}$ above horizontal. The pin reaction's
direction, which the hinge alone would leave open, is fixed before any
arithmetic.

\textbf{Equilibrium.} The cleanest route sums moments about $A$, which
kills the pin reaction (it passes through $A$) and leaves one equation in
$T$. The load $W$ acts at the tip a distance $L$ from $A$ along the boom,
with a horizontal lever arm $L\cos 25^{\circ}$ about $A$; the cable
tension's lever arm about $A$ is $L\sin\gamma$, since the perpendicular
distance from $A$ to the cable's line through $B$ is the boom length times
the sine of the angle between boom and cable. Moment balance:
\[
\sum M_A = 0: \quad T\,(L\sin\gamma) - W\,(L\cos 25^{\circ}) = 0
\quad\Longrightarrow\quad
T = \frac{W\cos 25^{\circ}}{\sin\gamma}.
\]
The boom length $L$ cancels, as it must for a weightless member. Numerically,
\[
T = \frac{2000 \times \cos 25^{\circ}}{\sin 40^{\circ}}
= \frac{2000 \times 0.906}{0.643} = 2{,}819\,\text{N}.
\]
The pin reaction follows from force balance, and concurrence has already
told us it lies along the boom. Resolving the cable tension into horizontal
and vertical components and adding the load, the pin must supply the balance.
The cable rises at $25^{\circ} + 40^{\circ} = 65^{\circ}$ above horizontal
(boom angle plus the angle from boom to cable), so its components are
$T\cos 65^{\circ} = 1{,}191\,\text{N}$ toward the wall and $T\sin 65^{\circ}
= 2{,}555\,\text{N}$ up. Horizontal balance gives the pin's outward push
$R_{Ax} = 1{,}191\,\text{N}$ and vertical balance gives $R_{Ay} = W -
2{,}555 = 2000 - 2555 = -555\,\text{N}$, a downward pull of $555\,\text{N}$.
The resultant is $R_A = \sqrt{1191^2 + 555^2} = 1{,}314\,\text{N}$ at
$\arctan(555/1191) = 25^{\circ}$ below the horizontal, along the boom toward
$A$, exactly the line concurrence predicted. The two methods agree, and the
concurrence direction is the cross-check that catches a sign slip in the
component work before it propagates.

\subsection{A stacked block: tipping versus sliding}

% Figure: a block pushed by a horizontal force near its top, illustrating
% the competition between sliding (friction limit) and tipping (moment
% about the leading bottom edge). Two reactions are drawn: the friction at
% the base and the leading-edge pivot the block tips about.
\begin{figure}[htbp]
\centering
\begin{tikzpicture}[
  push/.style={draw=engred, -{Stealth}, very thick},
  wt/.style={draw=engblue, -{Stealth}, very thick},
  nrm/.style={draw=engamber, -{Stealth}, very thick},
  fric/.style={draw=engorange, -{Stealth}, very thick},
  dim/.style={draw=enggray, thin, {Stealth}-{Stealth}},
  scale=1.0,
]
% ground
\draw[draw=engblack, very thick] (-1.0,0) -- (5.6,0);
\foreach \x in {-0.7,-0.2,...,5.3} {\draw[draw=enggray, thin] (\x,0) -- (\x-0.2,-0.25);}
% the block: width b, height H
\draw[draw=engblack, fill=engblue!8, thick] (1.2,0) rectangle (2.6,2.6);
% centre of mass
\coordinate (G) at (1.9,1.3);
\fill[engblack] (G) circle (0.06);
\node[font=\scriptsize, anchor=south east] at (G) {$G$};
% weight at G
\draw[wt] (G) -- ++(0,-1.0) node[anchor=north east, font=\scriptsize, engblue] {$mg$};
% applied push near the top, at height h
\coordinate (Ppt) at (1.2,2.2);
\draw[push] ($(Ppt)+(-1.1,0)$) -- (Ppt) node[midway, anchor=south, font=\scriptsize, engred] {$P$};
\node[font=\scriptsize, anchor=east, engred] at ($(Ppt)+(-1.1,0)$) {};
% normal reaction (resultant) shifted toward the leading edge
\draw[nrm] (2.3,0) -- ++(0,1.3) node[anchor=south, font=\scriptsize, engamber] {$N$};
% friction at the base, opposing the push
\draw[fric] (1.9,0.0) -- ++(0.9,0) node[anchor=west, font=\scriptsize, engorange] {$F_f$};
% leading bottom edge: the tipping pivot
\fill[engred] (2.6,0) circle (0.07);
\node[font=\scriptsize, anchor=north west, engred] at (2.6,0) {$O$};
% dimension hints: width b and push height h
\draw[dim] (1.2,-0.55) -- (2.6,-0.55);
\node[font=\scriptsize, anchor=north] at (1.9,-0.55) {$b$};
\draw[dim] (3.0,0) -- (3.0,2.2);
\node[font=\scriptsize, anchor=west] at (3.0,1.1) {$h$};
% a faint tipping arc about O
\draw[draw=engred, thin, dashed] (1.2,2.6) arc (143:175:3.2);
\node[font=\scriptsize, anchor=south west, engred] at (4.1,1.7) {tips about $O$};
\end{tikzpicture}
\caption[Tipping versus sliding of a pushed block]{A block of width $b$
and weight $mg$ pushed by a horizontal force $P$ applied at height $h$.
The chosen body is the block. The base supplies a distributed reaction
whose resultant is a normal force $N$ (which migrates toward the leading
edge as $P$ grows) and a friction force $F_f$ opposing the push. Two
failure modes compete. Sliding begins when $P$ reaches the friction
limit, $P = \mu_s mg$, independent of where $P$ is applied. Tipping
begins when $P$ generates enough moment about the leading bottom edge
$O$ to lift the trailing edge, $P h = mg\,(b/2)$, which depends on the
push height $h$. Whichever threshold is lower governs: a high push on a
high-friction floor tips, a low push on a slick floor slides. The two
lines cross at $h = b/(2\mu_s)$, the geometry at which the block is
equally close to both failures.}
\label{fig:vol03ch01:tipping-vs-sliding}
\end{figure}


A single contact can fail two ways at once, and the FBD must check both. A
filing cabinet of mass $m = 40\,\si{\kilogram}$, width $b = 0.5\,\si{\meter}$,
and height $H = 1.5\,\si{\meter}$ stands on a floor with $\mu_s = 0.5$. A
person pushes it horizontally at a height $h = 1.2\,\si{\meter}$, as in
figure~\ref{fig:vol03ch01:tipping-vs-sliding}. Find the push that starts it
sliding and the push that starts it tipping, and state which happens first.

\textbf{FBD.} The chosen body is the cabinet. The forces are the weight
$mg$ down at the centre of mass (height $H/2$, centred in width), the
applied push $P$ horizontal at height $h$, and the floor's distributed
reaction, whose resultant is a normal force $N$ and a friction force $F_f$.
As $P$ grows, the normal-force resultant migrates toward the leading bottom
edge $O$, because the floor cannot pull down and the pressure crowds to the
side the cabinet is rotating onto.

\textbf{The sliding threshold.} Sliding impends when the friction reaches
its limit. Vertical balance gives $N = mg = 40 \times 9.81 = 392\,\text{N}$,
so the friction ceiling is $\mu_s N = 0.5 \times 392 = 196\,\text{N}$.
Horizontal balance at impending slip is $P_{\text{slide}} = \mu_s mg =
196\,\text{N}$, and the push height does not enter: sliding is a force
balance only.

\textbf{The tipping threshold.} Tipping impends when the normal-force
resultant reaches the leading edge $O$ and the cabinet is about to rotate
about it. Take moments about $O$. The weight acts at the centre of width, a
horizontal distance $b/2$ from $O$, and resists tipping with moment $mg\,
(b/2)$. The push acts at height $h$ and drives tipping with moment $P h$. At
impending tip the two balance:
\[
P_{\text{tip}}\,h = mg\,\frac{b}{2}
\quad\Longrightarrow\quad
P_{\text{tip}} = \frac{mg\,b}{2h}
= \frac{392 \times 0.5}{2 \times 1.2} = 81.7\,\text{N}.
\]
\textbf{The verdict.} The tipping push, $81.7\,\text{N}$, is well below the
sliding push, $196\,\text{N}$, so the cabinet tips before it slides: the
push topples it onto the pusher rather than sliding it across the floor. The
two thresholds are equal at the critical height $h^\ast = b/(2\mu_s) = 0.5/
(2 \times 0.5) = 0.5\,\si{\meter}$; a push applied below $0.5\,\si{\meter}$
slides the cabinet, a push above it tips the cabinet. The lesson carries to
every tall load handled by pushing: a hand truck takes a load low precisely
to keep the push under the tipping height, and a fully loaded filing cabinet
with the heavy drawers high is a tipping hazard the moment it is shoved near
the top. The FBD that checks only sliding passes a load that the floor will
nonetheless let fall over.

\subsection{A particle on a frictional incline under an applied force}

The block-on-incline of figure~\ref{fig:vol03ch01:fbd-block-incline}
balanced gravity against friction alone. Add a force the operator controls
and the friction inequality becomes a band: there is a smallest push that
keeps the block from sliding down and a largest push before it is shoved up,
and the band between them is the range the friction holds. A crate of mass
$m = 30\,\si{\kilogram}$ rests on a ramp inclined at $\theta = 25^{\circ}$,
with $\mu_s = 0.4$ between crate and ramp. A force $P$ is applied parallel to
the ramp surface, directed up the slope. Find the range of $P$ for which the
crate stays put.

\textbf{FBD.} The chosen body is the crate. Four forces: the weight $mg$
down, the ramp normal $N$ perpendicular to the surface, the friction $F_f$
along the surface (direction not yet known, since it depends on which way the
crate tends to move), and the applied push $P$ up the slope. Resolve along
the ramp ($x$, up-slope positive) and across it ($y$).

\textbf{Across the ramp.} The normal balance is unaffected by $P$, which is
parallel to the surface: $N = mg\cos\theta = 30 \times 9.81 \times \cos
25^{\circ} = 266.7\,\text{N}$. The friction ceiling is $\mu_s N = 0.4 \times
266.7 = 106.7\,\text{N}$.

\textbf{The lower bound: about to slide down.} With $P$ too small, the crate
tends to slide down, so friction acts up the slope at its limit. Up-slope
balance: $P_{\min} + \mu_s N - mg\sin\theta = 0$, where the down-slope weight
component is $mg\sin\theta = 30 \times 9.81 \times \sin 25^{\circ} =
124.4\,\text{N}$. Then
\[
P_{\min} = mg\sin\theta - \mu_s N = 124.4 - 106.7 = 17.7\,\text{N}.
\]
\textbf{The upper bound: about to slide up.} With $P$ too large, the crate
tends to slide up, so friction reverses and acts down the slope at its limit.
Up-slope balance: $P_{\max} - \mu_s N - mg\sin\theta = 0$, so
\[
P_{\max} = mg\sin\theta + \mu_s N = 124.4 + 106.7 = 231.1\,\text{N}.
\]
\textbf{The band.} Any push between $17.7\,\text{N}$ and $231.1\,\text{N}$
holds the crate at rest, the friction quietly taking whatever value balances
the up-slope equation. The midpoint of the band, $124.4\,\text{N}$, is the
push that exactly cancels the weight component and leaves friction at zero;
on either side of it friction grows until it saturates at one of the two
bounds. The width of the band is $2\mu_s N = 213\,\text{N}$, twice the
friction ceiling, which is the operator's whole margin for error. The same
band underlies any held position on an incline: a car on a hill held by
partial throttle, a load nudged onto a ramp, a screw on its thread. The
friction force is not a number until the FBD says which way motion impends,
and on an incline under an applied force it can point either way.

\subsection{A rope over a fixed drum: distributed friction}

% Figure: a rope wrapped over a fixed drum (capstan). (a) the wrap geometry
% with the small-element FBD that yields the capstan equation; (b) the
% exponential tension ratio against wrap angle for two friction
% coefficients. Numbers match the capstan worked example.
\begin{figure}[htbp]
\centering
\begin{tikzpicture}[
  rope/.style={draw=engblue, line width=1.6pt},
  frc/.style={draw=engred, -{Stealth}, very thick},
  scale=1.0,
]
% ===== Left: wrap geometry and the element =====
\begin{scope}[xshift=0cm]
  % the drum
  \draw[draw=engblack, fill=enggray!12, very thick] (2.0,2.0) circle (1.2);
  \fill[engblack] (2.0,2.0) circle (0.05);
  % rope wraps over the top half (from left-down to right-down)
  \draw[rope] (0.8,0.6) -- (0.8,1.4)
    arc (180:0:1.2) -- (3.2,0.6);
  % tensions at the two ends
  \draw[frc] (0.8,1.0) -- ++(0,-1.0)
    node[anchor=north, font=\scriptsize, engred] {$T_1$ (hold)};
  \draw[frc] (3.2,1.0) -- ++(0,-1.0)
    node[anchor=north, font=\scriptsize, engred] {$T_2$ (load)};
  % wrap-angle marker
  \draw[draw=enggray, thin, <->] (0.55,1.4) arc (180:360:1.45);
  \node[font=\scriptsize, enggray] at (2.0,0.15) {wrap angle $\phi$};
  % small element
  \node[font=\scriptsize] at (2.0,3.55) {element $d\phi$};
  \draw[draw=engamber, thin] (2.0,3.2) -- (2.0,2.0);
  \node[font=\small, anchor=north] at (2.0,-0.3) {(a) rope over a drum};
\end{scope}
% ===== Right: tension ratio vs wrap angle =====
\begin{scope}[xshift=6.8cm, yshift=0.2cm]
  \begin{axis}[
    width=6.6cm, height=5.4cm,
    xlabel={wrap angle (turns)},
    ylabel={$T_2/T_1 = e^{\mu\phi}$},
    xmin=0, xmax=3, ymin=0.8, ymax=18,
    xtick={0,1,2,3},
    ytick={1,5,10,15},
    grid=both, grid style={enggray!25},
    label style={font=\scriptsize},
    tick label style={font=\scriptsize},
    every axis plot/.append style={line width=1.2pt},
  ]
    % phi in turns; mu*phi with phi=2*pi*turns
    \addplot[draw=engred, domain=0:3, samples=80] {exp(0.3*2*pi*x)};
    \addplot[draw=engblue, domain=0:3, samples=80] {exp(0.2*2*pi*x)};
    \node[font=\scriptsize, anchor=west, engred] at (axis cs:1.3,12) {$\mu=0.3$};
    \node[font=\scriptsize, anchor=west, engblue] at (axis cs:1.9,6) {$\mu=0.2$};
  \end{axis}
  \node[font=\small, anchor=north] at (3.0,-0.4) {(b) the capstan ratio};
\end{scope}
\end{tikzpicture}
\caption[A rope over a drum and the capstan tension ratio]{(a) A rope laid
over a fixed drum, holding a load tension $T_2$ with a much smaller grip
tension $T_1$. The free-body diagram of a small rope element subtending the
angle $d\phi$ balances the tension difference against the normal force from
the drum and the friction it supports; integrating around the wrap gives the
capstan relation $T_2 = T_1 e^{\mu\phi}$, where $\phi$ is the total wrap
angle in radians. (b) The ratio is exponential in wrap, so a few turns
around a bollard let a small hold force restrain a large load: at $\mu = 0.3$
a single full turn ($\phi = 2\pi$) gives a ratio of about $6.6$, and three
turns exceed a factor of $280$. The friction is distributed along the
contact rather than concentrated at one point, which is what turns the
linear block-on-incline picture into an exponential one.}
\label{fig:vol03ch01:capstan}
\end{figure}


Every friction example so far balanced friction at a single contact patch.
A rope wrapped over a fixed drum spreads its friction along the whole arc of
contact, and the spreading changes the arithmetic from linear to
exponential. The result is the relation a sailor uses without naming it when
a few turns of line around a bollard hold a vessel that no one could hold by
hand. Euler analysed the wrapped rope in 1762, and the relation carries his
name alongside Eytelwein's.

\textbf{FBD of an element.} The chosen body is a small element of rope
subtending the angle $d\phi$ at the drum centre, as in
figure~\ref{fig:vol03ch01:capstan}(a). Four forces act on it: the tension
$T$ at one end, the slightly larger tension $T + dT$ at the other, the
normal force $dN$ the drum pushes radially outward, and the friction $dF =
\mu\,dN$ the drum supplies along the contact, opposing the slip of the rope
toward the higher-tension side.

\textbf{Equilibrium of the element.} Resolve along the rope (tangentially)
and perpendicular to it (radially). For a small angle $d\phi$, the two
tension vectors each tilt by $d\phi/2$ from the tangent, so their radial
components are $T\sin(d\phi/2) \approx T\,d\phi/2$ each, summing to $T\,
d\phi$. The radial balance is $dN = T\,d\phi$. The tangential balance pits
the tension difference against the friction: $dT = \mu\,dN = \mu T\,d\phi$.
That is a separable differential equation, $dT/T = \mu\,d\phi$, integrated
across the full wrap angle $\phi$ from the slack tension $T_1$ to the taut
tension $T_2$:
\[
\int_{T_1}^{T_2} \frac{dT}{T} = \int_0^{\phi} \mu\,d\phi
\quad\Longrightarrow\quad
\ln\frac{T_2}{T_1} = \mu\phi
\quad\Longrightarrow\quad
T_2 = T_1\,e^{\mu\phi}.
\]
The capstan equation. The tension ratio is exponential in the wrap angle,
which is why wrap, not grip, does the work. Take $\mu = 0.3$ and a single
full turn, $\phi = 2\pi$: the ratio is $e^{0.3 \times 2\pi} = e^{1.885} =
6.6$, so a $100\,\text{N}$ hold restrains a $660\,\text{N}$ load. Three full
turns, $\phi = 6\pi$, give $e^{5.65} = 285$, so the same $100\,\text{N}$
hand restrains a $28{,}500\,\text{N}$ load. Figure~\ref{fig:vol03ch01:capstan}(b) plots the ratio against the number of turns for two
coefficients. The exponential is unforgiving in both directions: it is what
lets a winch operator hold a heavy line with one hand, and it is what makes
a rope that has taken three turns around a snagged drum nearly impossible to
pull free from the slack end. The capstan relation returns in Volume IV when
band brakes and belt drives are sized, where the same $e^{\mu\phi}$ sets the
torque a belt can transmit before it slips.

\subsection{A spatial force from its endpoints: direction cosines in use}

The planar examples never needed the third component. A force whose line
runs in space does, and the route is the position-vector recipe of section
1.1: write the vector along the member, normalise, scale by the magnitude.
A cable runs from an eye-bolt at $A = (3, 0, 0)\,\si{\meter}$ on the floor
to the top of a mast at $B = (0, 0, 6)\,\si{\meter}$ and is tensioned to
$T = 5{,}000\,\text{N}$. Find the three components of the force the cable
exerts on the mast top, and the angles the cable makes with the axes.

\textbf{The direction.} The cable pulls the mast top toward $A$, so the
relevant position vector runs from $B$ to $A$:
\[
\vec{r}_{BA} = A - B = (3 - 0,\; 0 - 0,\; 0 - 6) = (3, 0, -6)\,\si{\meter},
\quad
|\vec{r}_{BA}| = \sqrt{3^2 + 0^2 + (-6)^2} = \sqrt{45} = 6.708\,\si{\meter}.
\]
The unit vector along the pull is
\[
\hat{u} = \frac{(3, 0, -6)}{6.708} = (0.447,\; 0,\; -0.894),
\]
and its entries are the direction cosines: $\cos\theta_x = 0.447$,
$\cos\theta_y = 0$, $\cos\theta_z = -0.894$. The check is the unit-vector
identity, $0.447^2 + 0^2 + (-0.894)^2 = 0.200 + 0.800 = 1.000$, which holds
to rounding, so the direction is real.

\textbf{The force.} Scaling by the tension,
\[
\vec{F} = T\,\hat{u} = 5000 \times (0.447, 0, -0.894)
= (2{,}236,\; 0,\; -4{,}472)\,\text{N}.
\]
The cable pulls the mast top with $2{,}236\,\text{N}$ along $+x$ (toward the
anchor's horizontal position) and $4{,}472\,\text{N}$ straight down, with no
$y$ component because the anchor sits in the $xz$ plane. The angles follow
from the cosines: $\theta_x = \arccos 0.447 = 63.4^{\circ}$, $\theta_y =
90^{\circ}$, $\theta_z = \arccos(-0.894) = 153.4^{\circ}$, the last obtuse
because the force points downward against the upward $z$ axis. No angle was
measured; all three came from two endpoints and one magnitude. This is the
whole of spatial force resolution, and every leg of the tripod in section
1.6.7 and every cable in a guyed structure is this calculation repeated.

\subsection{A cranked handle: a force replaced by a force and a couple}

The force-couple reduction of section 1.1 is the tool that moves a force to
where the bearing actually feels it. A cranked handle turns a shaft: the
operator pushes down with $P = 80\,\text{N}$ at the end of a crank arm of
length $a = 0.25\,\si{\meter}$, the arm offset from the shaft axis. Find
the force and the couple the push delivers to the shaft axis.

\textbf{The reduction.} The push acts at the handle, a perpendicular
distance $a$ from the shaft axis. Move it to the axis by the force-couple
construction: the force on the axis is the same $P = 80\,\text{N}$ down,
and the couple is the push's moment about the axis,
\[
M = P\,a = 80 \times 0.25 = 20\,\text{N}\cdot\text{m},
\]
the torque that actually turns the shaft. The shaft axis feels a
$80\,\text{N}$ transverse force, carried by its bearings, together with a
$20\,\text{N}\cdot\text{m}$ twisting couple, carried by whatever the shaft
drives. The two effects are different in kind and go to different places:
the force is a bearing load, the couple is the useful output. A beginner
who reports only the torque has dropped the bearing load; a beginner who
reports only the $80\,\text{N}$ has dropped the torque the whole device
exists to deliver.

\textbf{Why the crank length sets the trade.} Doubling the crank arm to
$0.5\,\si{\meter}$ doubles the torque to $40\,\text{N}\cdot\text{m}$ for
the same hand push, while the bearing force stays at $80\,\text{N}$. The
longer arm buys torque at no cost in bearing load, which is why a breaker
bar turns a seized bolt that a short wrench cannot: the couple scales with
the arm, the transverse force does not. The force-couple split is the
statement that a force applied off-axis is a force on the axis plus a
torque about it, and the two scale by different rules.

\subsection{An eccentric load on a bolt group}

A bolt group loaded away from its centroid carries two demands at once: a
direct share of the load spread equally over the bolts, and a moment share
that grows with each bolt's distance from the group centroid. Separating
the two is the bolt-group method, and it is the force-couple reduction
applied to a fastener pattern. A bracket is bolted to a wall by four bolts
at the corners of a $0.2\,\si{\meter}$ square, and a vertical load
$P = 8{,}000\,\text{N}$ hangs from a point $e = 0.3\,\si{\meter}$ out from
the bolt-group centroid. Find the load on the most heavily loaded bolt.

\textbf{Reduce the load to the centroid.} The eccentric load is, by the
force-couple reduction, a load $P = 8{,}000\,\text{N}$ at the centroid plus
a couple $M = P\,e = 8000 \times 0.3 = 2{,}400\,\text{N}\cdot\text{m}$ about
the centroid. The two pieces are shared by different rules.

\textbf{The direct share.} The centroidal load splits equally over the four
bolts:
\[
F_{\text{direct}} = \frac{P}{4} = \frac{8000}{4} = 2{,}000\,\text{N}
\quad\text{down, on each bolt.}
\]

\textbf{The moment share.} The couple is resisted by bolt forces
perpendicular to the line from the centroid to each bolt, with magnitude
proportional to that distance. For the symmetric square, all four bolts sit
the same distance $d = \sqrt{0.1^2 + 0.1^2} = 0.1414\,\si{\meter}$ from the
centroid, so each carries an equal moment force $F_M$ with
\[
M = \sum F_{M,i}\,d_i = 4\,F_M\,d
\quad\Longrightarrow\quad
F_M = \frac{M}{4d} = \frac{2400}{4 \times 0.1414} = 4{,}243\,\text{N},
\]
directed perpendicular to the radius at each bolt.

\textbf{Combine on the worst bolt.} The most heavily loaded bolt is the one
where the direct and moment shares align most nearly. Resolve each into
components and add as vectors. The moment force on the top-right bolt points
up-and-right at $45^{\circ}$ from the radius; its components combine with
the $2{,}000\,\text{N}$ direct downward share to a resultant of
\[
F_{\max} = \sqrt{(F_{\text{direct},x} + F_{M,x})^2
+ (F_{\text{direct},y} + F_{M,y})^2} \approx 5{,}500\,\text{N},
\]
roughly two and three-quarter times the direct share alone. The eccentric
hang, not the load, governs the bolt: a designer who sized the bolts on
$P/4 = 2{,}000\,\text{N}$ would have under-sized them by nearly a factor of
three. The eccentricity is the term the force-couple reduction makes
visible and the equal-split intuition hides.

\subsection{A foot on a ladder rung: a three-dimensional reaction set}

A planar FBD has three reactions; a body that can twist out of plane has
six, and forgetting the out-of-plane three is how a real fitting fails a
load its planar analysis passed. A horizontal cantilever bracket is bolted
to a wall at one end and carries a load $W = 600\,\text{N}$ hung at its free
end, a distance $L = 0.4\,\si{\meter}$ out and offset $s = 0.15\,\si{\meter}$
to one side of the bolt line. Find the reactions the wall connection must
supply.

\textbf{FBD.} The chosen body is the bracket. The only external forces are
the load $W$ down at the free end and the wall reaction at the fixed end.
A fully fixed (built-in) spatial connection supplies all six reactions:
three force components $(R_x, R_y, R_z)$ and three moment components
$(M_x, M_y, M_z)$. We place the origin at the wall connection, with $x$ out
along the bracket, $y$ sideways, $z$ up.

\textbf{Force balance.} Only the vertical force is nonzero, so
\[
R_z = W = 600\,\text{N} \quad\text{up}, \qquad R_x = R_y = 0.
\]

\textbf{Moment balance.} The load acts at $\vec{r} = (0.4, 0.15, 0)\,
\si{\meter}$ with force $\vec{W} = (0, 0, -600)\,\text{N}$. Its moment about
the connection is the cross product,
\[
\vec{r}\times\vec{W}
= (0.15\times(-600) - 0,\;\; 0 - 0.4\times(-600),\;\; 0)
= (-90,\; 240,\; 0)\,\text{N}\cdot\text{m}.
\]
The connection must supply the opposite, $\vec{M} = (90, -240, 0)\,
\text{N}\cdot\text{m}$: a bending moment of $240\,\text{N}\cdot\text{m}$
about the $y$ axis (the cantilever bending the planar analysis expects) and
a twisting moment of $90\,\text{N}\cdot\text{m}$ about the $x$ axis, the
torsion the sideways offset $s$ produces. The torsion is the term a planar
FBD drawn in the vertical plane of the bracket never sees, because it acts
out of that plane. A connection sized for the $240\,\text{N}\cdot\text{m}$
bend alone is missing the $90\,\text{N}\cdot\text{m}$ twist, and a fitting
that carries bending well can be weak in torsion. The spatial FBD is not a
luxury here; it is the only diagram that shows the load that the offset
creates.

\subsection{A guyed pole: the base carries what the wire does not}

% Figure: a vertical pole held by a single guy wire, loaded at the top by
% a horizontal pull. The guy tension and the ground anchor carry the load;
% the base sees a large vertical thrust.
\begin{figure}[htbp]
\centering
\begin{tikzpicture}[
  pole/.style={draw=engblack, line width=2.0pt},
  wire/.style={draw=engblue, line width=1.2pt},
  force/.style={draw=engred, -{Stealth}, very thick},
  reac/.style={draw=englime!70!engblack, -{Stealth}, very thick},
  ground/.style={draw=enggray, thin},
  scale=1.0,
]
% ground line
\draw[ground] (-0.6,0) -- (5.2,0);
\foreach \x in {-0.4,0.0,0.4,0.8,1.2,1.6,2.0,2.4,2.8,3.2,3.6,4.0,4.4,4.8} {
  \draw[ground] (\x,0) -- ++(-0.18,-0.18);
}
% pole base at A, top at B
\coordinate (A) at (0,0);
\coordinate (B) at (0,4.0);
% anchor on the ground at C, to the right
\coordinate (C) at (3.0,0);
\draw[pole] (A) -- (B);
\draw[wire] (B) -- (C) node[midway, anchor=south west, font=\scriptsize, engblue] {$T$};
% base and anchor markers
\fill[engblack] (A) circle (0.07);
\node[font=\scriptsize, anchor=north east] at (A) {$A$};
\fill[engblack] (C) circle (0.07);
\node[font=\scriptsize, anchor=north] at (C) {$C$};
\node[font=\scriptsize, anchor=east] at (B) {$B$};
% applied horizontal load P at the top, pulling left
\draw[force] (B) -- ++(-1.4,0) node[anchor=east, font=\small, engred] {$P$};
% angle of the guy from the pole (vertical)
\draw[draw=enggray, thin] (0,3.2) arc (-90:-53:0.8);
\node[font=\scriptsize] at (0.55,3.05) {$\psi$};
% base reaction: large vertical thrust plus horizontal
\draw[reac] (A) -- ++(0,-1.2)
  node[anchor=north, font=\scriptsize, englime!60!engblack] {$A_y$ (down on pole)};
\end{tikzpicture}
\caption[A guyed pole loaded at the top]{A vertical pole pinned at its
base $A$ is held against a horizontal pull $P$ at the top $B$ by a single
guy wire running to a ground anchor $C$. The guy makes an angle $\psi$
with the pole. The guy can pull only along its own line, so its horizontal
component must balance $P$ while its vertical component presses the pole
down onto its base. A shallow guy (large $\psi$, anchor set far out)
balances $P$ cheaply in tension but with a modest downthrust; a steep guy
(small $\psi$, anchor close to the base) needs a large tension to find any
horizontal component, and drives a correspondingly large vertical thrust
into the base. The base, not the wire, often governs.}
\label{fig:vol03ch01:guy-wire}
\end{figure}


A single applied vignette closes the worked examples, the kind of small
load path a field engineer traces in their head before climbing. A vertical
utility pole of height $H = 8\,\si{\meter}$ is pinned at its base $A$ and
held against the horizontal pull of a tensioned line at the top, $P =
1{,}500\,\text{N}$, by one guy wire running from the top $B$ down to a
ground anchor $C$, as in figure~\ref{fig:vol03ch01:guy-wire}. The anchor
sits a horizontal distance $b = 5\,\si{\meter}$ from the base, so the guy
makes an angle $\psi = \arctan(b/H) = \arctan(5/8) = 32.0^{\circ}$ with the
pole. Find the guy tension and the force on the pole's base.

\textbf{FBD of the pole.} Three forces act: the horizontal line pull $P$ at
the top, the guy tension $T$ at the top along $BC$, and the base reaction
$(A_x, A_y)$ at the pin. We take the line pull and the guy on the same side,
so the guy's horizontal component opposes $P$.

\textbf{Moment about the base kills the pin.} Summing moments about $A$
leaves one equation in $T$. The line pull $P$ acts at height $H$ with lever
arm $H$; the guy tension's horizontal component $T\sin\psi$ acts at the same
height $H$ and opposes it, while the guy's vertical component passes through
nothing useful about $A$ until we resolve. Taking moments about $A$ with
horizontal forces at height $H$:
\[
\sum M_A = 0: \quad P\,H - (T\sin\psi)\,H = 0
\quad\Longrightarrow\quad
T\sin\psi = P.
\]
The guy's horizontal component alone balances the line pull, so
\[
T = \frac{P}{\sin\psi} = \frac{1500}{\sin 32.0^{\circ}}
= \frac{1500}{0.530} = 2{,}830\,\text{N}.
\]
The guy carries nearly twice the load it restrains, the price of its
shallow angle to the pole.

\textbf{The base reaction.} Force balance gives the pin force. Horizontally,
the guy's inward pull $T\sin\psi = 1{,}500\,\text{N}$ cancels $P$, so the
base supplies no net horizontal force: $A_x = 0$. Vertically, the guy's
downward component presses the pole onto its base,
\[
A_y = T\cos\psi = 2830 \times \cos 32.0^{\circ}
= 2830 \times 0.848 = 2{,}400\,\text{N} \quad\text{up, from the base.}
\]
The base carries a $2{,}400\,\text{N}$ vertical thrust, sixty percent more
than the $1{,}500\,\text{N}$ line pull that started the problem, and a
thrust that points nowhere the line pull did. The lesson is the chapter's
in miniature: the guy wire takes the horizontal load, but it pays for the
service by driving a vertical thrust into the base, and a pole footing sized
only against the line pull misses the larger load the wire creates. A
steeper guy (anchor closer in, smaller $\psi$) would raise both the tension
and the base thrust sharply; a shallower guy (anchor farther out) eases
both, which is why a generous guy radius is the cheap fix when the footing
is the worry. The force the wire saves the pole from is not the force the
base must carry.

\subsection{A wheel pulled over a step}

% Figure: a wheel being pulled over a step (curb). FBD shows the pull P at
% the axle, the weight W, and the single contact reaction at the step edge
% once the ground contact lifts. Numbers match the wheel-over-step worked
% example.
\begin{figure}[htbp]
\centering
\begin{tikzpicture}[
  force/.style={draw=engblue, -{Stealth}, very thick},
  rxn/.style={draw=engred, -{Stealth}, very thick},
  scale=1.0,
]
% ground and step
\draw[draw=engblack, very thick] (-0.2,0) -- (3.0,0);
\draw[draw=engblack, very thick] (3.0,0) -- (3.0,1.1) -- (6.2,1.1);
\foreach \x in {0.1,0.5,...,2.9} {\draw[draw=enggray, thin] (\x,0) -- (\x-0.2,-0.28);}
% wheel: radius r, resting against the step edge E at (3.0,1.1)
% centre chosen so wheel touches step corner
\coordinate (E) at (3.0,1.1);
\coordinate (C) at (2.05,1.6);
\def\rw{1.35}
\draw[draw=engblack, very thick] (C) circle (\rw);
\fill[engblack] (C) circle (0.05);
% line from centre to step edge (the contact-reaction line)
\draw[draw=enggray, thin, dashed] (C) -- (E);
% weight down at centre
\draw[force] (C) -- ++(0,-1.7) node[anchor=north, font=\small, engblue] {$W$};
% pull P horizontal at axle
\draw[force] (C) -- ++(1.9,0) node[anchor=west, font=\small, engblue] {$P$};
% contact reaction at step edge, along E->C
\draw[rxn] (E) -- ++($(C)-(E)$) node[anchor=south east, font=\small, engred] {$R$};
% geometry marks
\node[font=\scriptsize, anchor=east] at (2.05,1.6) {$C$};
\node[font=\scriptsize, anchor=north west] at (3.05,1.1) {$E$};
\draw[draw=enggray, thin, <->] (3.0,-0.05) -- (3.0,1.05);
\node[font=\scriptsize, enggray, anchor=west] at (3.08,0.5) {step $h$};
\end{tikzpicture}
\caption[A wheel pulled over a step]{The chosen body is the wheel at the
instant its ground contact lifts, so the only contact left is at the step
corner $E$. A smooth corner can push only along the line from the corner to
the wheel centre $C$, which fixes the direction of the step reaction $R$
before any arithmetic. Three forces act: the weight $W$ down at the axle, the
horizontal pull $P$ at the axle, and the corner reaction $R$ along $EC$. The
three lines meet at $C$, so the system is concurrent and a moment taken about
$E$ kills $R$ at a stroke, leaving one equation for the pull that just starts
the wheel over the step. A taller step needs a larger pull, and a pull
applied at the top of the wheel rather than the axle needs less, which is why
a hand cart is tilted back to climb a curb.}
\label{fig:vol03ch01:wheel-step}
\end{figure}


A round wheel meets a square curb at a corner, and the corner decides the pull
that lifts the wheel over. The contact that matters is not the flat ground but
the edge of the step, and the FBD at the instant of climbing has only that one
contact left. A wheel of radius $r = 0.30\,\si{\meter}$ carrying a load $W =
400\,\text{N}$ at its axle is to be rolled over a step of height $h = 0.10\,
\si{\meter}$ by a horizontal pull $P$ applied at the axle, as in
figure~\ref{fig:vol03ch01:wheel-step}. Find the pull that just starts the
wheel over.

\textbf{FBD.} The chosen body is the wheel at the instant its ground contact
lifts. Once the wheel begins to pivot over the corner $E$, the ground contact
carries no force, so only three forces remain: the weight $W$ down at the axle
$C$, the horizontal pull $P$ at the axle, and the reaction $R$ at the corner
$E$. A smooth corner can push only along the line from the corner to the
centre, so $R$ acts along $EC$, its direction fixed by geometry before any
arithmetic. Three forces, concurrent at $C$, and the corner is the point to
take moments about because $R$ passes through it.

\textbf{Geometry.} The horizontal distance from the corner $E$ to the vertical
through the centre is the moment arm the weight works on, and the vertical
distance from $E$ up to the axle is the arm the pull works on. With the axle a
height $r$ above the ground and the corner a height $h$, the corner sits a
vertical distance $r - h$ below the axle, and its horizontal offset from the
axle is
\[
x = \sqrt{r^2 - (r - h)^2} = \sqrt{h(2r - h)}
= \sqrt{0.10\,(0.60 - 0.10)} = \sqrt{0.05} = 0.2236\,\si{\meter}.
\]

\textbf{Equilibrium.} Take moments about the corner $E$. The weight drives the
wheel back down with moment $W\,x$; the pull lifts it over with moment $P\,(r -
h)$, since the pull is horizontal at a height $r - h$ above the corner. At the
verge of climbing the two balance:
\[
P\,(r - h) = W\,x
\quad\Longrightarrow\quad
P = W\,\frac{x}{r - h}
= 400 \times \frac{0.2236}{0.20} = 447\,\text{N}.
\]
The pull to climb the step exceeds the load itself, and it climbs steeply as
the step approaches the wheel radius: at $h = r$ the offset $x \to r$ while the
pull arm $r - h \to 0$, so $P \to \infty$ and no horizontal axle pull can lift
a wheel over a step as tall as its radius. Two design moves cut the pull. A
larger wheel lowers the ratio $x/(r-h)$ for a given step, which is why a cart
with large wheels crosses a curb a small-castored one jams against. Applying
the pull higher, at the top of the wheel rather than the axle, lengthens its
moment arm about the corner and drops the force needed, which is the mechanics
a person uses without naming it when they tip a hand truck back before
climbing a curb. The corner reaction, meanwhile, carries the vector sum of $W$
and $P$ through a single line of stone or steel, and a soft corner that
crushes turns the clean pivot into a jammed one.

\subsection{A bell-crank turning a control force}

% Figure: a bell-crank (right-angle lever) pinned at the corner, an input
% force on one arm and an output on the other, with the pin reaction. Numbers
% match the bell-crank worked example.
\begin{figure}[htbp]
\centering
\begin{tikzpicture}[
  force/.style={draw=engblue, -{Stealth}, very thick},
  rxn/.style={draw=engred, -{Stealth}, very thick},
  scale=1.0,
]
% pin at corner O
\coordinate (O) at (0,0);
\coordinate (A) at (0,2.6);   % end of vertical arm (input)
\coordinate (B) at (3.0,0);   % end of horizontal arm (output)
% the crank body: two arms meeting at O
\draw[draw=engblack, line width=2.4pt] (A) -- (O) -- (B);
% pin
\draw[draw=engblack, fill=enggray!20, thick] (O) circle (0.13);
\fill[engblack] (O) circle (0.04);
\node[font=\scriptsize, anchor=north east] at (O) {$O$};
% input force P at A, horizontal (pushing right)
\draw[force] (A) -- ++(1.7,0) node[anchor=west, font=\small, engblue] {$P$};
\node[font=\scriptsize, anchor=east] at (A) {$A$};
% arm lengths
\draw[draw=enggray, thin, <->] (-0.35,0) -- (-0.35,2.6);
\node[font=\scriptsize, enggray, anchor=east] at (-0.38,1.3) {$a$};
\draw[draw=enggray, thin, <->] (0,-0.35) -- (3.0,-0.35);
\node[font=\scriptsize, enggray, anchor=north] at (1.5,-0.4) {$b$};
% output force Q at B, vertical (up)
\draw[force] (B) -- ++(0,1.7) node[anchor=south, font=\small, engblue] {$Q$};
\node[font=\scriptsize, anchor=north west] at (B) {$B$};
% pin reaction (schematic, at O, up-left)
\draw[rxn] (O) -- ++(-1.2,1.0) node[anchor=south east, font=\small, engred] {$R_O$};
\end{tikzpicture}
\caption[A bell-crank lever]{The chosen body is the bell-crank, a rigid
right-angle lever pinned at the corner $O$. An input push $P$ acts
horizontally at the end $A$ of the arm of length $a$; the output $Q$ leaves
vertically at the end $B$ of the arm of length $b$. The pin supplies two
reaction components, drawn here as a single resultant $R_O$. Summing moments
about $O$ kills the pin reaction, because it passes through $O$, and leaves
one equation relating input to output: $P\,a = Q\,b$. The crank turns the
line of a force through a right angle and scales its magnitude by the arm
ratio $a/b$, which is why a bell-crank is the fitting of choice wherever a
control force must change direction, from a bicycle brake to an aircraft
control run. The pin then carries the vector sum of the two arm forces, a
load neither arm alone reveals.}
\label{fig:vol03ch01:bell-crank}
\end{figure}


A force sometimes has to leave a mechanism pointing the way it did not arrive,
and the bell-crank is the fitting that turns it. A bell-crank is a rigid lever
pinned at a corner with two arms at an angle, and it trades the direction of a
force for another direction while scaling the magnitude by the arm ratio. The
crank of figure~\ref{fig:vol03ch01:bell-crank} is pinned at $O$, with a
vertical input arm of length $a = 0.12\,\si{\meter}$ carrying a horizontal push
$P = 150\,\text{N}$ at its end $A$, and a horizontal output arm of length $b =
0.20\,\si{\meter}$ delivering a vertical force $Q$ at its end $B$. Find the
output force and the pin reaction.

\textbf{FBD.} The chosen body is the crank. Three forces act: the input $P$
horizontal at $A$, the output $Q$ vertical at $B$, and the pin reaction $(O_x,
O_y)$ at the corner. The pin supplies two components and no moment, as a pin
always does.

\textbf{The output from a single moment equation.} Sum moments about the pin
$O$, which kills both pin components at once. The input push $P$ acts at the
top of the vertical arm, a perpendicular distance $a$ from $O$, driving the
crank one way; the output $Q$ acts at the end of the horizontal arm, a
perpendicular distance $b$ from $O$, resisting:
\[
\sum M_O = 0: \quad P\,a - Q\,b = 0
\quad\Longrightarrow\quad
Q = P\,\frac{a}{b} = 150 \times \frac{0.12}{0.20} = 90\,\text{N}.
\]
The crank turns a horizontal $150\,\text{N}$ push into a vertical $90\,\text{N}$
output, scaled down by the arm ratio $a/b = 0.6$. A crank with a long output
arm reduces the force and multiplies the travel; one with a long input arm
does the reverse. The direction change is free, a right angle here, set by the
angle between the arms.

\textbf{The pin reaction.} The pin carries whatever the two arm forces do not
balance between themselves. Force balance gives
\[
O_x = P = 150\,\text{N}, \qquad O_y = Q = 90\,\text{N},
\]
so the pin resultant is $R_O = \sqrt{150^2 + 90^2} = 175\,\text{N}$, at
$\arctan(90/150) = 31^{\circ}$ above the horizontal. The pin load exceeds
either arm force, because the pin resists the vector sum of a horizontal input
and a vertical output that do not oppose each other. A designer who sizes the
pin on the larger arm force alone under-sizes it, the same lesson the pulley
axle taught: the mount for a force-turning device carries a resultant its two
rope or arm forces never show singly. The bell-crank runs through a bicycle
brake, a piano action, and the control runs of an aircraft, and in every case
the pin is the part the loads gang up on.

\subsection{An A-frame: two bodies and the force between them}

% Figure: an A-frame of two legs pinned at the apex, feet on the ground, a
% load at the apex. (a) the whole frame; (b) one leg isolated, showing the
% apex-pin interaction force that is internal to (a) and external to (b).
\begin{figure}[htbp]
\centering
\begin{tikzpicture}[
  force/.style={draw=engblue, -{Stealth}, very thick},
  rxn/.style={draw=engred, -{Stealth}, very thick},
  intn/.style={draw=engorange, -{Stealth}, very thick},
  scale=1.0,
]
% ===== (a) whole frame =====
\begin{scope}
  \coordinate (Ap) at (1.5,2.8);   % apex
  \coordinate (Lf) at (0,0);       % left foot
  \coordinate (Rf) at (3.0,0);     % right foot
  \draw[draw=engblack, very thick] (-0.4,0) -- (3.4,0);
  \foreach \x in {-0.3,0.1,...,3.3} {\draw[draw=enggray, thin] (\x,0) -- (\x-0.2,-0.26);}
  \draw[draw=engblack, line width=2.2pt] (Lf) -- (Ap) -- (Rf);
  \draw[draw=engblack, fill=enggray!20, thick] (Ap) circle (0.11);
  % load at apex
  \draw[force] (Ap) -- ++(0,-1.6) node[anchor=west, font=\small, engblue] {$W$};
  % foot reactions
  \draw[rxn] (Lf) -- ++(0,1.4) node[anchor=south, font=\small, engred] {$N_L$};
  \draw[rxn] (Rf) -- ++(0,1.4) node[anchor=south, font=\small, engred] {$N_R$};
  \node[font=\scriptsize, anchor=north east] at (Lf) {$L$};
  \node[font=\scriptsize, anchor=north west] at (Rf) {$R$};
  \node[font=\scriptsize, anchor=south] at (Ap) {$P$};
  \node[font=\small, anchor=north] at (1.5,-0.5) {(a) whole frame};
\end{scope}
% ===== (b) one leg isolated =====
\begin{scope}[xshift=6.2cm]
  \coordinate (Ap2) at (1.5,2.8);
  \coordinate (Lf2) at (0,0);
  \draw[draw=engblack, very thick] (-0.4,0) -- (2.2,0);
  \foreach \x in {-0.3,0.1,...,2.1} {\draw[draw=enggray, thin] (\x,0) -- (\x-0.2,-0.26);}
  \draw[draw=engblack, line width=2.2pt] (Lf2) -- (Ap2);
  \draw[draw=engblack, fill=enggray!20, thick] (Ap2) circle (0.11);
  % foot reaction
  \draw[rxn] (Lf2) -- ++(0,1.4) node[anchor=south, font=\small, engred] {$N_L$};
  % pin interaction from the other leg (external to this body)
  \draw[intn] (Ap2) -- ++(1.5,-0.6) node[anchor=west, font=\small, engorange] {$\vec{F}_{RL}$};
  \node[font=\scriptsize, anchor=north east] at (Lf2) {$L$};
  \node[font=\scriptsize, anchor=south] at (Ap2) {$P$};
  \node[font=\small, anchor=north] at (1.0,-0.5) {(b) left leg alone};
\end{scope}
\end{tikzpicture}
\caption[An A-frame and the apex-pin interaction]{(a) The whole A-frame, two
legs pinned at the apex $P$ with feet at $L$ and $R$ and a load $W$ hung at
the apex. Drawn as one body, the frame sees only external forces: the load
and the two foot reactions. The apex pin is internal and never appears. (b)
The left leg isolated as its own body. The cut through the apex pin now
exposes the force $\vec{F}_{RL}$ that the right leg exerts on the left, an
interaction force that was internal in (a) and is external in (b). By
Newton's third law the left leg pushes back on the right with $-\vec{F}_{RL}$,
which appears in the right leg's own diagram. The whole-frame diagram finds
the foot reactions; the single-leg diagram finds the apex-pin force the frame
diagram hides. A multi-body assembly is solved by choosing which cuts to make
and accepting the interaction forces those cuts expose.}
\label{fig:vol03ch01:a-frame}
\end{figure}


Every worked example so far chose a single body. A real assembly is several
bodies pinned together, and the discipline that carries across the joins is the
third law: the force one body exerts on its neighbour is equal and opposite to
the force the neighbour exerts back, and a cut through the joint exposes both.
An A-frame of two legs pinned at the apex, feet on the ground, carrying a load
at the apex, is the smallest problem that forces this bookkeeping, drawn in
figure~\ref{fig:vol03ch01:a-frame}. Two identical legs of length $\ell$ are
pinned at the apex $P$ and stand with their feet $L$ and $R$ a distance $2s =
2.0\,\si{\meter}$ apart on a smooth floor, the apex a height $H = 2.4\,\si{
\meter}$ up. A load $W = 1{,}600\,\text{N}$ hangs from the apex pin. Find the
foot reactions and the force in the apex pin, taking the legs as rigid and
weightless and the feet held from spreading by a floor tie.

\textbf{Whole frame first.} The chosen body is the whole frame. Drawn as one
body, the apex pin is internal and never appears; only external forces show.
By symmetry each smooth-floor foot carries a vertical reaction
\[
N_L = N_R = \frac{W}{2} = 800\,\text{N}.
\]
A smooth floor supplies no horizontal force, so the outward thrust of the
splayed legs must be held by the tie between the feet. Summing horizontal
forces on the whole frame gives nothing, because the two tie forces are
internal and cancel; the tie force is found only by cutting a single leg out.

\textbf{One leg alone.} Now isolate the left leg, figure~\ref{fig:vol03ch01:a-frame}(b).
The cut through the apex pin exposes the force $\vec{F}_{RL}$ that the right leg
exerts on the left, external to this body though internal to the frame. The
left leg sees four forces: its foot reaction $N_L = 800\,\text{N}$ up, the tie
force $F_t$ horizontal at the foot, the apex-pin force $\vec{F}_{RL}$, and (by
the load hanging at the shared apex pin) a share of the load. The clean route
sums moments about the apex, which kills the apex-pin force and leaves the tie.
The foot sits a horizontal distance $s = 1.0\,\si{\meter}$ and a vertical
distance $H = 2.4\,\si{\meter}$ from the apex. Moments about $P$ for the left
leg:
\[
\sum M_P = 0: \quad N_L\,s - F_t\,H = 0
\quad\Longrightarrow\quad
F_t = \frac{N_L\,s}{H} = \frac{800 \times 1.0}{2.4} = 333\,\text{N}.
\]
The tie between the feet carries $333\,\text{N}$ in tension, the outward thrust
each splayed leg would otherwise drive into the floor. A wider frame (feet
farther apart, $s$ larger) raises the tie force; a taller, steeper frame lowers
it. The apex-pin force follows from force balance on the leg: vertically the
pin supplies $800\,\text{N}$ down onto the leg to match the foot reaction less
its load share, and horizontally $333\,\text{N}$ to match the tie, so the pin
carries a resultant along the leg's own line, as a two-force idealisation would
predict once the tie is accounted for.

\textbf{The lesson of the cut.} The whole-frame diagram found the foot
reactions and could not find the tie or the pin, because both are internal to
it. The single-leg diagram found them, at the cost of carrying the apex-pin
interaction as an external force. Choosing which body to draw is choosing which
forces are internal (hidden) and which are external (exposed), and a
multi-body problem is solved by making the cuts that expose the forces the
question asks for. The trap is to draw one leg and forget that the apex pin now
pushes on it; the third law says the pin force is there whether or not the
beginner draws it, and the leg that omits it balances a diagram the real leg
does not obey.

\subsection{A hinged trapdoor held open by a cable}

% Figure: a hinged trapdoor held open by a cable. The door is a uniform slab
% hinged at one edge, propped open by a cable to a wall point above. FBD shows
% weight at the centroid, hinge reaction, cable tension. Numbers match the
% trapdoor worked example.
\begin{figure}[htbp]
\centering
\begin{tikzpicture}[
  force/.style={draw=engblue, -{Stealth}, very thick},
  rxn/.style={draw=engred, -{Stealth}, very thick},
  cable/.style={draw=engamber, line width=1.4pt},
  scale=1.0,
]
% wall (vertical)
\draw[draw=engblack, very thick] (0,-0.4) -- (0,3.4);
\foreach \y in {-0.3,0.1,...,3.3} {\draw[draw=enggray, thin] (0,\y) -- (-0.26,\y-0.2);}
% hinge at H at wall base
\coordinate (H) at (0,0);
\draw[draw=engblack, fill=enggray!20, thick] (H) circle (0.12);
\node[font=\scriptsize, anchor=north east] at (H) {$H$};
% door raised at angle phi above horizontal, length Ld
\def\phi{35}
\coordinate (T) at ({3.4*cos(\phi)},{3.4*sin(\phi)});
\draw[draw=engblack, line width=2.4pt] (H) -- (T);
% centroid
\coordinate (G) at ({1.7*cos(\phi)},{1.7*sin(\phi)});
\fill[engblack] (G) circle (0.06);
% cable from wall point C (up the wall) to door tip T
\coordinate (C) at (0,3.0);
\draw[cable] (C) -- (T);
\draw[draw=engblack, fill=enggray!20, thick] (C) circle (0.08);
\node[font=\scriptsize, anchor=east] at (C) {$C$};
% weight at centroid
\draw[force] (G) -- ++(0,-1.5) node[anchor=north, font=\small, engblue] {$W$};
% cable tension along door tip toward C
\draw[rxn] (T) -- ($(T)!0.42!(C)$) node[anchor=south west, font=\small, engred] {$T$};
% hinge reaction (schematic)
\draw[rxn] (H) -- ++(1.1,0.7) node[anchor=south west, font=\small, engred] {$R_H$};
% angle phi at hinge
\draw[draw=enggray, thin] (0.9,0) arc (0:\phi:0.9);
\node[font=\scriptsize, enggray] at (1.05,0.32) {$\phi$};
\node[font=\scriptsize, anchor=south west] at (T) {tip};
\end{tikzpicture}
\caption[A hinged trapdoor held open by a cable]{The chosen body is the
trapdoor, a uniform slab hinged at $H$ and propped open at angle $\phi$ above
horizontal by a cable running from the door tip to the wall point $C$. Three
external forces act: the weight $W$ down at the centroid, the cable tension
$T$ along the tip-to-$C$ line, and the hinge reaction $R_H$. Summing moments
about the hinge kills the hinge reaction and leaves one equation for the cable
tension. The weight's moment arm about $H$ shrinks with the horizontal
projection of the door, so a door raised nearer the vertical needs less cable
tension to hold, while a nearly horizontal door demands the most. The cable's
own geometry sets how much of its tension does useful lifting: a cable pulling
nearly along the door has a small moment arm and must carry a large tension to
supply the holding moment.}
\label{fig:vol03ch01:trapdoor}
\end{figure}


A slab hinged along one edge and propped open by a cable is a moment problem
with the cable geometry folded in, and it is the model for every gate, hatch,
and access panel held at an angle. A uniform trapdoor of mass $m = 30\,\si{
\kilogram}$ and length $\ell = 1.2\,\si{\meter}$ is hinged at its lower edge
$H$ to a wall and held open at $\phi = 35^{\circ}$ above horizontal by a cable
from the door tip to a wall point $C$ a height $d = 1.0\,\si{\meter}$ above the
hinge, as in figure~\ref{fig:vol03ch01:trapdoor}. Find the cable tension and
the hinge reaction.

\textbf{FBD.} The chosen body is the door. Three external forces: the weight
$W = mg = 30 \times 9.81 = 294\,\text{N}$ down at the centroid (mid-length),
the cable tension $T$ along the line from the tip to $C$, and the hinge
reaction $(H_x, H_y)$. Summing moments about the hinge kills the hinge reaction
and isolates $T$.

\textbf{Geometry of the two moment arms.} Place the hinge at the origin. The
tip sits at $(\ell\cos\phi, \ell\sin\phi) = (0.983, 0.688)\,\si{\meter}$, and
the wall point $C$ at $(0, 1.0)\,\si{\meter}$. The weight acts at the centroid,
a horizontal distance $\tfrac{\ell}{2}\cos\phi = 0.491\,\si{\meter}$ from the
hinge, so its moment about $H$ (closing the door) is
\[
M_W = W\,\frac{\ell}{2}\cos\phi = 294 \times 0.491 = 144\,\text{N}\cdot\text{m}.
\]
The cable runs from the tip $(0.983, 0.688)$ to $C\,(0, 1.0)$, a direction
$(-0.983, 0.312)$ of length $\sqrt{0.983^2 + 0.312^2} = 1.031\,\si{\meter}$, so
its unit vector is $(-0.953, 0.303)$. The cable's moment about the hinge is the
cross product of the tip position with the tension, giving a moment arm
\[
d_T = |x_{\text{tip}}\,u_y - y_{\text{tip}}\,u_x|
= |0.983 \times 0.303 - 0.688 \times (-0.953)| = |0.298 + 0.656| = 0.954\,\si{\meter}.
\]

\textbf{Equilibrium.} Moment balance about the hinge, the cable holding the
door up against its weight:
\[
T\,d_T = M_W
\quad\Longrightarrow\quad
T = \frac{144}{0.954} = 151\,\text{N}.
\]
The cable carries about half the door's weight in tension to hold it at this
angle, the geometry setting how efficiently the cable's pull becomes a holding
moment. A cable anchored higher up the wall pulls more nearly perpendicular to
the door and holds it with less tension; a cable anchored low pulls nearly
along the door, has a short moment arm, and must carry a large tension for the
same hold. The hinge reaction then follows from force balance, carrying the
weight less the cable's vertical component and the cable's horizontal pull
inward, a reaction the hinge leaf and its fasteners must be sized against. A
door that slams open and loads the cable dynamically, or a cable anchored to
save wall height at the cost of a shallow angle, is the field version of the
tension the static number understates.

\subsection{A screw jack: the wedge wrapped into a thread}

The wedge of section 1.6.7 lifted a load with a friction inequality that also
made it self-locking. Wrap that same inclined plane around a cylinder and it
becomes a screw thread, and the wedge analysis becomes the screw jack, the
device that lifts a car with a hand crank and holds it up when the crank is
released. A single-start square-thread jack has a mean thread radius $r_m =
20\,\si{\milli\meter}$, a lead (axial rise per turn) $L = 6\,\si{\milli\meter}$,
and a thread friction coefficient $\mu = 0.15$. Find the torque to raise a load
$W = 20\,\text{kN}$ and check whether the jack is self-locking.

\textbf{The thread as an unrolled incline.} Unroll one turn of the thread and
it is an inclined plane whose rise is the lead $L$ and whose run is the thread
circumference $2\pi r_m$. The helix angle $\lambda$, the incline's slope, is
\[
\tan\lambda = \frac{L}{2\pi r_m} = \frac{6}{2\pi \times 20} = \frac{6}{125.7}
= 0.0477, \quad \lambda = 2.73^{\circ}.
\]
The friction angle is $\phi = \arctan\mu = \arctan 0.15 = 8.53^{\circ}$. The
load on the thread is the axial force $W$; raising it is pushing a wedge up an
incline against friction, and the tangential force at the mean radius needed to
do so is $W\tan(\lambda + \phi)$, the friction angle added to the helix angle
because friction opposes the upward motion.

\textbf{The raising torque.} The torque is that tangential force times the mean
radius:
\[
M_{\text{raise}} = W\,r_m\,\tan(\lambda + \phi)
= 20{,}000 \times 0.020 \times \tan(2.73^{\circ} + 8.53^{\circ}).
\]
With $\tan(11.26^{\circ}) = 0.199$,
\[
M_{\text{raise}} = 20{,}000 \times 0.020 \times 0.199 = 79.6\,\text{N}\cdot\text{m}.
\]
A torque of about $80\,\text{N}\cdot\text{m}$ raises two tonnes, a force of
$400\,\text{N}$ on a $0.2\,\si{\meter}$ crank handle, well within one hand.

\textbf{The self-locking check.} Releasing the crank asks whether the load can
drive the thread back down, the mirror of the wedge's back-drive test. Lowering
against friction takes a torque $M_{\text{lower}} = W\,r_m\,\tan(\phi -
\lambda)$; when $\phi > \lambda$ this is positive, meaning a torque is needed
even to \emph{let} the load down, so the load cannot lower itself and the jack
is self-locking. Here $\phi = 8.53^{\circ}$ exceeds $\lambda = 2.73^{\circ}$
comfortably, so the jack holds its load with the crank released, the property
that lets a car sit on a jack without a hand on the handle. The self-locking
condition is $\mu > \tan\lambda$, a friction inequality identical in form to
the wedge's, and a thread cut with too steep a lead (a fast, coarse thread)
can slip below it and drop its load, which is why a lifting screw is cut fine
and a fast-traverse screw, cut coarse, is never trusted to hold.

\textbf{The efficiency it costs.} The price of self-locking is efficiency. The
work out per turn is $W L$; the work in is $2\pi M_{\text{raise}}$; the ratio
is
\[
\eta = \frac{W L}{2\pi M_{\text{raise}}}
= \frac{20{,}000 \times 0.006}{2\pi \times 79.6} = \frac{120}{500} = 0.24,
\]
so barely a quarter of the crank work reaches the load and the rest is spent in
thread friction. A self-locking screw is a low-efficiency screw by necessity:
the friction that holds the load when released is the same friction that wastes
three-quarters of the input while raising it, a trade the wedge's friction cone
foretold and the jack pays every lift.

\subsection{A differential band brake}

The capstan relation of section 1.6.16 governed a rope slipping over a fixed
drum. Anchor both ends of a wrapped band to a lever and the same $e^{\mu\phi}$
becomes a brake, one whose bite depends on which way the drum turns and where
the band ends are pinned. A band brake wraps a flexible band over an angle
$\phi = 270^{\circ}$ of a drum of radius $R = 0.15\,\si{\meter}$, with band
friction $\mu = 0.35$. The band's two ends are pinned to a lever pivoted at
$O$: the tight end a distance $c_2 = 0.05\,\si{\meter}$ from the pivot and the
slack end $c_1 = 0.15\,\si{\meter}$ on the far side. A brake torque $M_b =
200\,\text{N}\cdot\text{m}$ is required. Find the lever force and the effect of
reversing the drum's rotation.

\textbf{Band tensions from the capstan relation.} The wrapped band is a
capstan: the two end tensions differ by $e^{\mu\phi}$. In radians the wrap is
$\phi = 270^{\circ} = 3\pi/2 = 4.712\,\text{rad}$, so
\[
\frac{T_2}{T_1} = e^{\mu\phi} = e^{0.35 \times 4.712} = e^{1.649} = 5.20.
\]
The brake torque is the drum's radius times the friction drag, which is the
difference of the two tensions:
\[
M_b = (T_2 - T_1)\,R
\quad\Longrightarrow\quad
T_2 - T_1 = \frac{M_b}{R} = \frac{200}{0.15} = 1{,}333\,\text{N}.
\]
With $T_2 = 5.20\,T_1$, the difference is $4.20\,T_1 = 1{,}333\,\text{N}$, so
\[
T_1 = 317\,\text{N}, \qquad T_2 = 1{,}650\,\text{N}.
\]

\textbf{The lever force.} The lever is a body pinned at $O$ with the two band
tensions pulling on it at their offsets. For the drum turning so the tight end
$T_2$ is the one nearer the pivot, moments about $O$ balance the applied lever
force $F$ (at arm $a = 0.4\,\si{\meter}$) against the two band pulls:
\[
F\,a + T_1\,c_1 - T_2\,c_2 = 0
\quad\Longrightarrow\quad
F = \frac{T_2\,c_2 - T_1\,c_1}{a}
= \frac{1650 \times 0.05 - 317 \times 0.15}{0.4}
= \frac{82.5 - 47.6}{0.4} = 87\,\text{N}.
\]
A modest $87\,\text{N}$ on the lever develops a $200\,\text{N}\cdot\text{m}$
brake torque, the band's exponential wrap doing the amplifying.

\textbf{The direction dependence.} Reverse the drum, and the tight and slack
ends swap: the tension that was $T_2$ is now the slack side and the geometry
that subtracted the large tension's moment now adds it. The lever force needed
becomes
\[
F' = \frac{T_2\,c_1 - T_1\,c_2}{a}
= \frac{1650 \times 0.15 - 317 \times 0.05}{0.4}
= \frac{247.5 - 15.9}{0.4} = 579\,\text{N},
\]
nearly seven times the force in the first direction. A band brake is
direction-sensitive: in one sense of rotation the drum's own drag helps pull
the band tight (a self-energising or self-actuating action) and the operator's
force is small; in the other sense the operator fights the band and the force
climbs steeply. A brake that grabs hard in one direction and slips in the other
is the signature of this asymmetry, and a hoist band brake is arranged so the
self-energising direction is the one that resists a falling load, letting a
light hand hold a heavy descent. The capstan exponential, read through a lever,
is a brake whose character the geometry of the two anchor points sets.

\subsection{The moment of an oblique force, computed two ways}

Varignon's theorem earns its keep on any force that arrives at an awkward angle,
and the fastest way to trust it is to run one force through both routes and watch
the numbers agree. A force $F = 500\,\text{N}$ acts at the point $(0.3, 0.2)\,
\si{\meter}$ measured from a pivot $O$, directed $60^{\circ}$ above the positive
$x$ axis. Find its moment about $O$.

\textbf{The perpendicular-distance route.} The single moment arm is the
perpendicular distance from $O$ to the force's line of action. The line passes
through $(0.3, 0.2)$ with direction $(\cos 60^{\circ}, \sin 60^{\circ})$, so the
signed distance is $d = x\sin 60^{\circ} - y\cos 60^{\circ} = 0.3 \times 0.866 -
0.2 \times 0.5 = 0.260 - 0.100 = 0.160\,\si{\meter}$. The moment magnitude is
$M = F\,d = 500 \times 0.160 = 79.9\,\text{N}\cdot\text{m}$, counterclockwise.
The step a sketch gets wrong is that oblique distance $d$: it is neither the
$0.3$ nor the $0.2$ that the coordinates hand over, and reading it off a drawing
invites the exact error Varignon exists to remove.

\textbf{The component route.} Resolve the force first: $F_x = 500\cos 60^{\circ}
= 250\,\text{N}$ and $F_y = 500\sin 60^{\circ} = 433\,\text{N}$. Each component
has an obvious moment arm, the coordinate perpendicular to it, and Varignon adds
the two:
\[
M = x\,F_y - y\,F_x = 0.3 \times 433 - 0.2 \times 250 = 130.0 - 50.0
= 79.9\,\text{N}\cdot\text{m},
\]
the same number, reached with no oblique distance measured at all. The two
routes agree because they are the same cross product read two ways, the scalar
$M_z = r_x F_y - r_y F_x$ of section 1.2.2. The working lesson is that the
component route is the safe one whenever the geometry is given as coordinates,
because it replaces one distance a drawing can distort with two distances the
coordinates already carry. The perpendicular-distance route wins only when the
arm is obvious by inspection, a load hanging a clean horizontal reach from a
pivot, where reaching for components would be busywork.

\subsection{A ladder with a person standing on it}

The ladder of section 1.6.1 carried only its own weight. Put a person on it and
the analysis is unchanged in form but the numbers move, and where the person
stands moves them most. A uniform ladder of mass $m_L = 12\,\si{\kilogram}$ and
length $L = 5\,\si{\meter}$ leans at $\theta = 65^{\circ}$ to the floor against a
smooth wall, its foot on a rough floor. A person of mass $m_P = 75\,\si{
\kilogram}$ stands three-quarters of the way up, at $f = 0.75$ of the ladder
length from the foot. Find the wall and floor reactions and the coefficient of
friction the floor must supply.

\textbf{FBD.} The chosen body is the ladder with the person's weight applied to
it as an external load at the person's position. Five forces: the ladder weight
$W_L = m_L g = 117.7\,\text{N}$ at mid-length, the person's weight $W_P = m_P g =
735.8\,\text{N}$ at $f L = 3.75\,\si{\meter}$ up, the wall normal $N_W$
(horizontal, wall smooth), the floor normal $N_F$ (vertical), and the floor
friction $F_f$ (horizontal, toward the wall). We treat the person as a static
load, an assumption that fails the moment they move quickly.

\textbf{Equilibrium.} Vertical balance gives the floor normal, $N_F = W_L + W_P
= 117.7 + 735.8 = 853.5\,\text{N}$. Summing moments about the foot kills $N_F$
and $F_f$, both of which pass through it, and leaves the wall normal:
\[
N_W\,L\sin\theta = W_L\,\tfrac{L}{2}\cos\theta + W_P\,(fL)\cos\theta
\quad\Longrightarrow\quad
N_W = \frac{\big(W_L\,\tfrac{1}{2} + W_P\,f\big)\cos\theta}{\sin\theta}.
\]
Numerically $N_W = (117.7 \times 0.5 + 735.8 \times 0.75)\cos 65^{\circ}/\sin
65^{\circ} = (58.9 + 551.9)\times 0.466 = 284.8\,\text{N}$. Horizontal balance
then gives $F_f = N_W = 284.8\,\text{N}$, and the required floor coefficient is
\[
\mu_s \ge \frac{F_f}{N_F} = \frac{284.8}{853.5} = 0.334.
\]

\textbf{Where the person stands is the whole story.} The empty ladder at this
angle needed only $\mu_s \ge 1/(2\tan 65^{\circ}) = 0.233$. The person
three-quarters up has pushed the demand to $0.334$, close to the coefficient of
a dry wood-on-concrete contact and past a damp or dusty one. Had the person
stood at mid-height ($f = 0.5$), the demand would have been lower, because the
higher the load sits on the ladder the larger its moment about the foot and the
larger the wall reaction the friction must match. The rule a careful worker
follows without deriving it, that the danger grows as one climbs toward the top
of a ladder pitched too shallow, is this equation read for its $f$ dependence:
the friction demand climbs with the load's height on the ladder, and a ladder
safe to stand on at the bottom can slip out at the top.

\subsection{A sign panel on a post under wind}

A flat sign on a post is a distributed-pressure problem folded into a base
moment, and the panel's own area, not the post, sets the load. A rectangular
sign $w = 2.0\,\si{\meter}$ wide and $a = 1.0\,\si{\meter}$ tall is centred at a
height $H = 3.0\,\si{\meter}$ on a vertical post. A storm wind of dynamic
pressure $q = 600\,\text{Pa}$ (a gust near $31\,\text{m/s}$) blows square on the
face, with drag coefficient $C_d = 1.2$ for a flat plate. Find the wind force,
the moment at the post base, and the tension it drives into the base-plate bolts.

\textbf{The wind force and its point of application.} The pressure over the
panel is uniform to the model, so its resultant is the pressure times the area,
acting through the panel centroid at height $H$:
\[
F_w = q\,C_d\,A = 600 \times 1.2 \times (2.0 \times 1.0) = 1{,}440\,\text{N},
\]
a distributed load reduced to a single force by the mastery box of section 1.5.
The centroid of a uniform pressure on a rectangle is the rectangle's centre, so
$F_w$ acts at $H = 3.0\,\si{\meter}$ with no offset to compute.

\textbf{The base moment.} The post is a cantilever, and the wind force at height
$H$ works a moment arm $H$ about the base:
\[
M_b = F_w\,H = 1{,}440 \times 3.0 = 4{,}320\,\text{N}\cdot\text{m}.
\]

\textbf{The bolt tension.} The base plate is held by four bolts on a square of
side $s = 0.3\,\si{\meter}$. The overturning moment is resisted as a couple: the
two bolts on the windward side go into tension, the leeward edge into bearing,
and the couple's lever arm is roughly the bolt spacing $s$. Modelling the
tension pair at arm $s$ from the compression edge, the tension-side demand is
\[
T \approx \frac{M_b}{s} = \frac{4{,}320}{0.3} = 14{,}400\,\text{N},
\]
about $7{,}200\,\text{N}$ per bolt across the two windward bolts. The visible
load, a $1{,}440\,\text{N}$ push on the sign, has become a five-times-larger
pull in the bolts, the cantilever's height amplifying the force through the
short bolt lever arm. The lesson repeats the traffic-mast vignette at hand
scale: a sign that stands quietly for years under gravity, which loads the bolts
hardly at all, meets a wind that loads them past their capacity, and the base
plate, not the panel, is where a badly sized signpost fails.

\subsection{A spatial ring on three cables}

% Figure: a spatial three-cable ring. A load hangs from a ring at the origin;
% three cables run up to three anchor points at different heights and bearings.
% Drawn in a simple oblique projection. The FBD is the ring as a concurrent
% spatial node: three cable tensions plus the load must sum to zero.
\begin{figure}[htbp]
\centering
\begin{tikzpicture}[
  cable/.style={draw=engred, -{Stealth}, very thick},
  load/.style={draw=engblue, -{Stealth}, very thick},
  axis/.style={draw=enggray, -{Stealth}, thin},
  scale=1.0,
]
% oblique projection: x to the right, y up, z into the page as down-left.
% Projected coordinates precomputed as (x + z*zx, y + z*zy), zx=-0.5, zy=-0.35.
% origin (the ring)
\coordinate (O) at (0,0);
% three anchors in space (metres, matching the worked example)
% A=(0.0,4.0,3.0)  -> (-1.50, 2.95)
% B=(-3.2,4.5,-1.6) -> (-2.40, 5.06)
% C=(2.6,3.8,-1.2)  -> (3.20, 4.22)
\coordinate (A) at (-1.50, 2.95);
\coordinate (B) at (-2.40, 5.06);
\coordinate (C) at (3.20, 4.22);
% reference axes at O
\draw[axis] (O) -- ++(1.3,0) node[anchor=west, font=\scriptsize, enggray] {$x$};
\draw[axis] (O) -- ++(0,1.3) node[anchor=south, font=\scriptsize, enggray] {$y$};
\draw[axis] (O) -- ++(-0.65,-0.455) node[anchor=north east, font=\scriptsize, enggray] {$z$};
% cables from O up to anchors
\draw[cable] (O) -- (A) node[midway, anchor=north east, font=\small, engred] {$T_A$};
\draw[cable] (O) -- (B) node[midway, anchor=south east, font=\small, engred] {$T_B$};
\draw[cable] (O) -- (C) node[midway, anchor=south west, font=\small, engred] {$T_C$};
% anchor markers
\foreach \p in {A,B,C} {
  \fill[engblack] (\p) circle (0.06);
  \node[font=\scriptsize, anchor=south] at (\p) {$\p$};
}
% the ring at O and the hanging load
\draw[draw=engblack, fill=enggray!20, thick] (O) circle (0.11);
\draw[load] (O) -- ++(0,-1.6) node[anchor=north, font=\small, engblue] {$W$};
\node[font=\scriptsize, anchor=east] at (-0.15,-0.05) {ring};
\end{tikzpicture}
\caption[A spatial three-cable anchor point]{Three cables hold a ring against a
hanging load $W$. The cables run from the ring at the origin up to three
anchors $A$, $B$, and $C$ fixed at different heights and bearings in space. The
free-body diagram of the ring is a concurrent spatial node: the load and the
three cable tensions all pass through one point, so there is no moment equation
and the whole problem is the single vector balance $T_A\hat{u}_A + T_B\hat{u}_B
+ T_C\hat{u}_C + \vec{W} = \vec{0}$, three scalar equations in the three unknown
tensions. Each unit vector $\hat{u}$ is built from the ring-to-anchor position
vector, normalised, exactly the direction-cosine recipe: write the vector along
the cable, divide by its length, and the tension follows from a $3\times 3$
linear solve. Unlike the symmetric tripod, no symmetry collapses the system
here, so all three components must be carried, and a negative tension in the
solution would signal a cable that would have to push, meaning it goes slack and
the assumed geometry does not hold.}
\label{fig:vol03ch01:spatial-anchor}
\end{figure}


The tripod of section 1.6.7 leaned on symmetry to collapse three equations to
one. Strip the symmetry away and the spatial node shows its full shape: three
force equations in three unknowns, solved together. A ring hangs at the origin
under a load $W = 2{,}400\,\text{N}$ (about $245\,\si{\kilogram}$) directed
straight down, held by three cables running to anchors at $A = (0, 4.0, 3.0)$,
$B = (-3.2, 4.5, -1.6)$, and $C = (2.6, 3.8, -1.2)\,\si{\meter}$, with $y$
vertical, as in figure~\ref{fig:vol03ch01:spatial-anchor}. Find the tension in
each cable.

\textbf{FBD.} The chosen body is the ring, a point. Four forces meet there: the
load $\vec{W} = (0, -2{,}400, 0)\,\text{N}$ and the three cable tensions, each
directed from the ring toward its anchor along the unit vector $\hat{u} =
\vec{r}/|\vec{r}|$. Building each from its anchor coordinates,
\[
\hat{u}_A = \frac{(0, 4.0, 3.0)}{5.000} = (0, 0.800, 0.600), \quad
\hat{u}_B = \frac{(-3.2, 4.5, -1.6)}{5.749} = (-0.557, 0.783, -0.278),
\]
\[
\hat{u}_C = \frac{(2.6, 3.8, -1.2)}{4.758} = (0.546, 0.799, -0.252),
\]
each direction cosine passing the check $\ell^2 + m^2 + n^2 = 1$ to rounding.

\textbf{Equilibrium.} The single vector balance $T_A\hat{u}_A + T_B\hat{u}_B +
T_C\hat{u}_C + \vec{W} = \vec{0}$ is three scalar equations, one per axis:
\begin{align*}
x: &\quad -0.557\,T_B + 0.546\,T_C = 0, \\
y: &\quad 0.800\,T_A + 0.783\,T_B + 0.799\,T_C - 2{,}400 = 0, \\
z: &\quad 0.600\,T_A - 0.278\,T_B - 0.252\,T_C = 0.
\end{align*}
No axis has a single unknown, so the three solve together. The result, from the
$3\times 3$ system, is
\[
T_A = 927\,\text{N}, \qquad T_B = 1{,}039\,\text{N}, \qquad T_C = 1{,}058\,\text{N}.
\]
All three tensions are positive, so every cable pulls, and the geometry holds.
The check is the vertical sum: $0.800 \times 927 + 0.783 \times 1{,}039 + 0.799
\times 1{,}058 = 742 + 814 + 845 = 2{,}401\,\text{N}$, balancing the $2{,}400\,
\text{N}$ load to rounding. The script \texttt{code/resolve\_forces.py} assembles
the three unit vectors from the endpoints and solves the system, the spatial
analogue of the planar force-polygon closure. A negative tension in such a
solution is the signal that a cable has been asked to push: it would go slack,
the assumed three-cable geometry would not hold, and the load would redistribute
onto the two that remain, a case the linear solve flags by its sign before any
cable is cut.

\subsection{A toggle clamp near dead centre}

% Figure: a toggle (over-centre) clamp near its locked position. A hand force
% on the handle drives two links whose common pin approaches the line joining
% their outer pins; as the pin crosses that line the clamping force runs up
% toward the singular value and the mechanism locks. Panel (a) the linkage;
% panel (b) the FBD of the middle pin as a three-force node.
\begin{figure}[htbp]
\centering
\begin{tikzpicture}[
  force/.style={draw=engblue, -{Stealth}, very thick},
  rxn/.style={draw=engred, -{Stealth}, very thick},
  link/.style={draw=engblack, line width=2.2pt},
  scale=1.0,
]
% ----- panel (a): the two-link toggle -----
\begin{scope}[shift={(0,0)}]
\coordinate (A) at (0,0);      % fixed pin (clamp arm side)
\coordinate (B) at (4.0,0);    % fixed pin (handle pivot side)
\coordinate (K) at (2.0,0.55); % knee pin, just above line AB
% base line AB (dashed reference)
\draw[draw=enggray, dashed, thin] (A) -- (B) node[midway,below=1pt,font=\scriptsize,enggray] {toggle line};
% two links
\draw[link] (A) -- (K);
\draw[link] (B) -- (K);
% pins
\foreach \p in {A,B,K} {\draw[draw=engblack, fill=enggray!20, thick] (\p) circle (0.10);}
\node[font=\scriptsize, anchor=north east] at (A) {$A$};
\node[font=\scriptsize, anchor=north west] at (B) {$B$};
\node[font=\scriptsize, anchor=south] at (K) {$K$};
% handle continues past B
\draw[link] (B) -- ++(1.5,0.9);
\draw[force] ($(B)+(1.5,0.9)$) -- ++(0.35,0.75) node[anchor=south west, font=\small, engblue] {$F_h$};
\node[font=\scriptsize, enggray, anchor=west] at ($(B)+(1.5,0.9)+(0.1,-0.2)$) {handle};
% clamp output at A (the arm presses the workpiece down)
\draw[force] (A) -- ++(0,-1.0) node[anchor=north, font=\small, engblue] {$Q$};
\node[font=\scriptsize, anchor=north east] at (1.9,-0.9) {(a)};
% small angle epsilon at knee
\node[font=\scriptsize, enggray, anchor=south] at (2.0,0.68) {$\varepsilon$};
\end{scope}

% ----- panel (b): FBD of the knee pin K -----
\begin{scope}[shift={(7.6,0.1)}]
\coordinate (Kp) at (0,0);
\fill[engblack] (Kp) circle (0.05);
\node[font=\scriptsize, anchor=south] at (Kp) {$K$};
% link forces along each link (nearly collinear, small angle to horizontal)
\draw[rxn] (Kp) -- ++(-2.3,-0.6) node[anchor=east, font=\small, engred] {$F_{AK}$};
\draw[rxn] (Kp) -- ++(2.3,-0.6) node[anchor=west, font=\small, engred] {$F_{BK}$};
% driving force from the handle, pushing the knee down toward the toggle line
\draw[force] (Kp) -- ++(0,1.3) node[anchor=south, font=\small, engblue] {$D$};
% angle markers
\draw[draw=enggray, thin] (-0.9,-0.235) arc (180+14.6:180:0.9);
\node[font=\scriptsize, enggray, anchor=east] at (-0.95,-0.13) {$\varepsilon$};
\node[font=\scriptsize, anchor=north] at (0,-0.95) {(b)};
\end{scope}
\end{tikzpicture}
\caption[A toggle (over-centre) clamp near lock]{An over-centre toggle clamp
approaching its locked position. Panel (a): two links join the fixed pins $A$
and $B$ at a knee pin $K$ that sits a small distance above the straight line
$AB$. A hand force $F_h$ on the handle drives the knee down toward that line,
and the clamp arm at $A$ presses the work with the output force $Q$. Panel
(b): the free-body diagram of the knee pin $K$, a three-force node carrying the
driving force $D$ from the handle and the two nearly collinear link forces
$F_{AK}$ and $F_{BK}$, each inclined at the small angle $\varepsilon$ to the
toggle line. Force balance at the knee gives $F_{AK}=F_{BK}=D/(2\sin
\varepsilon)$, which runs away as $\varepsilon\to 0$: the clamping force grows
without bound as the knee nears the toggle line, so a modest hand force
develops a large clamp force at the instant of lock. Driving the knee just past
the line puts the mechanism over centre, where the work's reaction can no
longer push it back, and the clamp holds with the hand removed.}
\label{fig:vol03ch01:toggle-clamp}
\end{figure}


The screw jack traded efficiency for a large force from a small input by
friction; a toggle trades it by geometry, developing a large clamp force from a
small hand force as a linkage nears a straight line. A toggle clamp joins two
fixed pins $A$ and $B$ through a knee pin $K$ that sits a small angle
$\varepsilon$ above the line $AB$, as in
figure~\ref{fig:vol03ch01:toggle-clamp}. A driving force $D = 200\,\text{N}$
from the handle pushes the knee toward that line. Find the force in the links,
and its behaviour as the knee approaches dead centre.

\textbf{FBD of the knee.} The chosen body is the knee pin $K$, a point. Three
forces meet there: the driving force $D$ pushing the knee toward the line $AB$,
and the two link forces $F_{AK}$ and $F_{BK}$ directed along the nearly
collinear links, each inclined at the small angle $\varepsilon$ to the line. By
the symmetry of the two links, $F_{AK} = F_{BK} = F$.

\textbf{Equilibrium.} Resolve along the direction of $D$. Each link's force has
a component $F\sin\varepsilon$ opposing $D$, so the balance is
\[
2F\sin\varepsilon = D
\quad\Longrightarrow\quad
F = \frac{D}{2\sin\varepsilon}.
\]
The link force is the driving force divided by twice the sine of the knee angle,
and it runs away as $\varepsilon \to 0$. At $\varepsilon = 10^{\circ}$, $F =
200/(2\sin 10^{\circ}) = 576\,\text{N}$, a mechanical advantage of about three;
at $\varepsilon = 5^{\circ}$, $F = 1{,}147\,\text{N}$, advantage near six; at
$\varepsilon = 2^{\circ}$, $F = 2{,}865\,\text{N}$, advantage above fourteen; at
$\varepsilon = 1^{\circ}$ it is $5{,}730\,\text{N}$. The clamp force, carried by
the link into the work, climbs toward the singular value the instant the knee
reaches the line, which is why a toggle clamp bites hardest just before it
locks.

\textbf{Over centre, and why it holds.} Pushing the knee a hair past the line
puts the mechanism over centre. The work's reaction on the clamp now tends to
drive the knee further past the line rather than back through it, so the linkage
holds with the hand removed, the same self-locking behaviour the wedge and the
screw showed, reached here by geometry rather than friction. A small stop limits
the over-travel, and the release takes a deliberate pull that unseats the knee
back across the line. The toggle is the clamp of choice on a welding fixture or
a machine-tool hold-down for exactly this pair of properties: a large clamp
force from a light hand, and a locked hold that needs no continued force, both
of them the single relation $F = D/(2\sin\varepsilon)$ read near its singular
limit.

\subsection{Applied vignette: a two-point climbing anchor}

% Figure: a two-point climbing anchor. Two placements share a load through a
% sling to a shared master point; the included angle sets the leg tension.
% (a) the anchor geometry; (b) leg tension vs included angle.
\begin{figure}[htbp]
\centering
\begin{tikzpicture}[
  force/.style={draw=engblue, -{Stealth}, very thick},
  rxn/.style={draw=engred, -{Stealth}, very thick},
  sling/.style={draw=engamber, line width=1.4pt},
  scale=1.0,
]
% ===== (a) two-point anchor =====
\begin{scope}
  % rock face at top
  \draw[draw=engblack, very thick] (-0.4,3.2) -- (3.8,3.2);
  \foreach \x in {-0.3,0.1,...,3.7} {\draw[draw=enggray, thin] (\x,3.2) -- (\x-0.2,3.46);}
  % two placements
  \coordinate (P1) at (0.4,3.2);
  \coordinate (P2) at (3.0,3.2);
  \fill[engblack] (P1) circle (0.07);
  \fill[engblack] (P2) circle (0.07);
  \node[font=\scriptsize, anchor=south] at (P1) {placement 1};
  \node[font=\scriptsize, anchor=south] at (P2) {placement 2};
  % master point below
  \coordinate (M) at (1.7,1.2);
  \draw[sling] (P1) -- (M);
  \draw[sling] (P2) -- (M);
  \draw[draw=engblack, fill=enggray!20, thick] (M) circle (0.11);
  % load down at master point
  \draw[force] (M) -- ++(0,-1.4) node[anchor=north, font=\small, engblue] {$W$};
  % leg tensions
  \draw[rxn] (M) -- ($(M)!0.42!(P1)$) node[anchor=south east, font=\scriptsize, engred] {$T_1$};
  \draw[rxn] (M) -- ($(M)!0.42!(P2)$) node[anchor=south west, font=\scriptsize, engred] {$T_2$};
  % included angle
  \draw[draw=enggray, thin] ($(M)!0.7!(P1)$) arc (128:52:0.85);
  \node[font=\scriptsize, enggray] at (1.7,2.05) {$2\theta$};
  \node[font=\small, anchor=north] at (1.7,-0.4) {(a) two-point anchor};
\end{scope}
% ===== (b) leg tension vs included angle =====
\begin{scope}[xshift=6.4cm, yshift=0.1cm]
  \begin{axis}[
    width=6.4cm, height=5.2cm,
    xlabel={included angle $2\theta$ (deg)},
    ylabel={leg tension $/\,W$},
    xmin=0, xmax=170, ymin=0.4, ymax=3.2,
    xtick={0,60,120,180},
    ytick={0.5,1,1.5,2,2.5,3},
    grid=both, grid style={enggray!25},
    label style={font=\scriptsize},
    tick label style={font=\scriptsize},
    every axis plot/.append style={line width=1.2pt},
  ]
    % T/W = 1/(2 cos theta), theta = x/2 degrees
    \addplot[draw=engred, domain=0:165, samples=120] {1/(2*cos(x/2))};
    \node[font=\scriptsize, anchor=west, engred] at (axis cs:10,0.65) {$T/W = \tfrac{1}{2\cos\theta}$};
  \end{axis}
  \node[font=\small, anchor=north] at (3.0,-0.4) {(b) the angle penalty};
\end{scope}
\end{tikzpicture}
\caption[A two-point climbing anchor and its angle penalty]{(a) Two rock
placements share a load through a sling to a single master point, the fitting
a climber clips into. The two legs meet at the master point with an included
angle $2\theta$. Each leg carries a tension $T = W/(2\cos\theta)$, the same
hanging-mass result read as a rig the climber builds. (b) The leg tension per
unit load climbs slowly to $60^{\circ}$ of included angle, then sharply: at a
$120^{\circ}$ included angle each leg already carries the full load, and past
that the tension runs away. The working rule that an anchor's included angle
should stay under about $60^{\circ}$ is this curve read for its knee, and a
wide sling on a narrow ledge is the field mistake that overloads both
placements at once.}
\label{fig:vol03ch01:climbing-anchor}
\end{figure}


A climber building an anchor solves the hanging-mass problem with their life on
the answer, and the geometry they control is the angle between the two legs.
The device is two placements in the rock, a cam or a nut in each, joined by a
sling to a single master point the climber clips into, as in
figure~\ref{fig:vol03ch01:climbing-anchor}. The load $W$ is the climber's
weight amplified by any fall, and it hangs from the master point where the two
sling legs meet at an included angle $2\theta$. Each leg carries the
hanging-mass tension read for a symmetric pair,
\[
T = \frac{W}{2\cos\theta},
\]
the same relation the ceiling ropes obeyed in section 1.6.3, now a rig the
climber assembles from slings and rock.

The number that matters is how the leg tension climbs with the included angle.
At a narrow angle, legs nearly parallel and hanging straight down, $\theta \to
0$ and each leg carries half the load, the best case. Widen the angle and
$\cos\theta$ falls, so the tension rises: at $2\theta = 60^{\circ}$ each leg
carries $W/(2\cos 30^{\circ}) = 0.577\,W$, a modest fifteen percent over the
even split; at $2\theta = 90^{\circ}$ it is $0.707\,W$; at $2\theta =
120^{\circ}$ each leg carries the \emph{full} load $W$, because $\cos
60^{\circ} = 0.5$ cancels the factor of two. Past $120^{\circ}$ the tension
runs away, and a nearly straight sling would load each placement far beyond the
weight it holds. The working rule taught in climbing courses, that the included
angle should stay under about $60^{\circ}$ and never approach $120^{\circ}$, is
figure~\ref{fig:vol03ch01:climbing-anchor}(b) read for its knee: below the knee
the penalty is small and slowly growing, above it the curve turns up and each
degree costs more than the last.

The field mistake the mechanics predicts is the wide-angle anchor. A climber
with two placements far apart and a short sling is forced to a large included
angle, and the geometry that looks like a solid wide stance is in fact loading
each placement harder than a single piece directly below would. The fix is a
longer sling, which drops the master point and narrows the angle, or placements
closer together. The same equation warns against the American triangle, a
tying method that routes the sling so the angle at the master point is forced
wide: the knot looks tidy and the load path is the worst case. A climber who
has internalised $T = W/(2\cos\theta)$ reads an anchor's angle the way an
engineer reads a sling plan, and the two are the same free-body diagram with
the same divergence at shallow angles, one drawn on a workshop floor and the
other on a cliff.

\subsection{Applied vignette: a cantilevered traffic-signal mast}

A traffic signal hung out over the middle of an intersection stands on a single
post at the kerb, and the whole load path is a cantilever: everything the arm
carries becomes a moment at the base plate, and the base plate's anchor bolts
are where the mast lives or falls. A mast arm reaches $R = 9\,\si{\meter}$ out
from a vertical post, carrying a signal head of weight $W_s = 500\,\text{N}$ at
the tip and a self-weight of the arm we lump as $W_a = 1{,}200\,\text{N}$
acting at its mid-length, $4.5\,\si{\meter}$ out. The post base is a plate held
by anchor bolts on a bolt circle of diameter $D = 0.5\,\si{\meter}$. We trace
the load to the bolts, first under gravity alone, then with the wind that
usually governs.

\textbf{The gravity moment at the base.} The two vertical loads each work a
horizontal moment arm about the base:
\[
M_g = W_s\,R + W_a\,\frac{R}{2} = 500 \times 9 + 1{,}200 \times 4.5
= 4{,}500 + 5{,}400 = 9{,}900\,\text{N}\cdot\text{m}.
\]
This overturning moment is resisted by the anchor bolts as a couple: the bolts
on the far side of the base from the arm go into tension, those on the near
side into bearing, and the couple's lever arm is roughly the bolt-circle
diameter $D$. Modelling the tension as carried by the far bolts at an arm $D$,
the tension demand is
\[
T_{\text{bolt}} \approx \frac{M_g}{D} = \frac{9{,}900}{0.5}
= 19{,}800\,\text{N},
\]
split among the tension-side bolts. The signal head, the visible load, supplied
under half the moment; the arm's own weight, spread along its reach, supplied
the rest, which is the self-weight lesson of the judgment exercise made
concrete: a long cantilever's own weight is not a rounding error but a leading
term.

\textbf{The wind case that governs.} A traffic mast spends its life in air, and
the wind on the broad signal head and the long arm usually sets the design. A
storm wind of dynamic pressure $q = 700\,\text{Pa}$ (a gust near $34\,\text{m/s}$)
on a signal head of frontal area $A_s = 0.6\,\si{\meter\squared}$ with drag
coefficient $C_d = 1.2$ gives a horizontal wind force
\[
F_w = q\,C_d\,A_s = 700 \times 1.2 \times 0.6 = 504\,\text{N},
\]
applied at the tip height and, worse for the base, at the full $9\,\si{\meter}$
reach if the wind blows along the arm, or driving a twisting moment about the
post if it blows across. Blowing across the arm, the wind force at the tip
drives a torsional moment about the vertical post of $F_w\,R = 504 \times 9 =
4{,}536\,\text{N}\cdot\text{m}$, a twist the base plate must carry in addition
to the gravity bending. The mast base therefore sees bending from gravity,
bending from along-arm wind, and torsion from across-arm wind, three moment
components at one connection, and the governing bolt is loaded by their
combination, not by any one alone.

The lesson is the spatial-bracket example of section 1.6.20 at civil scale: a
load offset from a connection produces a moment, and a load offset in two
directions produces both bending and torsion. A base plate sized for the
gravity bend alone misses the wind torsion, and a traffic mast that stands for
years under its own weight can fail in a storm that adds a moment the static
diagram never drew. The engineer who signs the base plate draws the wind cases,
resolves each into the bolt group, and sizes the worst bolt against the
combination, because the intersection below the mast is full of cars and the
cantilever hides none of its load once the free-body diagram is drawn in full.

\subsection{Applied vignette: a canal lock mitre gate}

% Figure: a canal mitre lock gate seen in plan. Two leaves meet at a mitre
% pointing upstream; water pressure on the upstream face drives each leaf
% against its neighbour at the mitre and against the quoin (heel) post in the
% wall. Panel (a) plan of the pair; panel (b) FBD of one leaf as a three-force
% body. Numbers match the lock-gate vignette.
\begin{figure}[htbp]
\centering
\begin{tikzpicture}[
  force/.style={draw=engblue, -{Stealth}, very thick},
  rxn/.style={draw=engred, -{Stealth}, very thick},
  water/.style={draw=engblue!45, -{Stealth}, thick},
  scale=1.0,
]
% ----- panel (a): plan of the mitred pair -----
\begin{scope}[shift={(0,0)}]
% chamber walls
\draw[draw=enggray, line width=2pt] (-0.2,2.2) -- (-0.2,0.6);
\draw[draw=enggray, line width=2pt] (5.2,2.2) -- (5.2,0.6);
\node[font=\scriptsize, enggray, anchor=east] at (-0.25,1.9) {wall};
% quoin (heel) posts in the wall corners
\fill[enggray!35] (-0.2,0.6) rectangle (0.15,0.95);
\fill[enggray!35] (4.85,0.6) rectangle (5.2,0.95);
\node[font=\scriptsize, enggray, anchor=south east] at (0.15,0.95) {quoin};
% two leaves meeting at the mitre point M (pointing upstream, downward here)
\coordinate (Lh) at (0.0,0.8);   % left heel
\coordinate (Rh) at (5.0,0.8);   % right heel
\coordinate (M)  at (2.5,-0.15);  % mitre point, upstream
\draw[draw=engblack, line width=2.4pt] (Lh) -- (M);
\draw[draw=engblack, line width=2.4pt] (Rh) -- (M);
\fill[engblack] (M) circle (0.05);
\node[font=\scriptsize, anchor=north] at (M) {$M$};
\node[font=\scriptsize, anchor=south] at (Lh) {heel};
% upstream / downstream labels
\node[font=\scriptsize\itshape, engblue, anchor=north] at (2.5,-0.55) {upstream (high water)};
\node[font=\scriptsize\itshape, enggray, anchor=south] at (2.5,1.15) {downstream (low water)};
% water pressure arrows on the upstream face of the left leaf (pushing downstream)
\draw[water] (1.55,0.05) -- ++(0.0,0.55);
\draw[water] (1.05,0.28) -- ++(0.0,0.55);
\draw[water] (0.55,0.52) -- ++(0.0,0.55);
\node[font=\scriptsize, anchor=north east] at (0.55,0.05) {(a)};
\end{scope}

% ----- panel (b): FBD of one leaf as a three-force body -----
\begin{scope}[shift={(8.4,0.4)}]
\coordinate (H) at (0,0);      % heel (quoin) reaction point
\coordinate (Mp) at (2.5,-0.9); % mitre contact point
\draw[draw=engblack, line width=2.4pt] (H) -- (Mp);
\fill[engblack] (H) circle (0.05);
\fill[engblack] (Mp) circle (0.05);
\node[font=\scriptsize, anchor=south east] at (H) {heel $H$};
\node[font=\scriptsize, anchor=north west] at (Mp) {mitre $M$};
% resultant water thrust P, acting at the leaf mid-length, perpendicular to leaf
\coordinate (C) at (1.25,-0.45);
\draw[force] (C) -- ++(0.62,0.83) node[anchor=south west, font=\small, engblue] {$P$};
\node[font=\scriptsize, engblue, anchor=north east] at (C) {mid-span};
% heel reaction at H (component form): horizontal H_x and along-gate H_y
\draw[rxn] (H) -- ++(-1.0,0.35) node[anchor=east, font=\small, engred] {$R_H$};
% mitre reaction: normal to the contact face at M
\draw[rxn] (Mp) -- ++(1.05,0.30) node[anchor=west, font=\small, engred] {$R_M$};
\node[font=\scriptsize, anchor=north] at (1.25,-1.45) {(b)};
\end{scope}
\end{tikzpicture}
\caption[A mitre lock gate as a three-force body]{A canal lock's mitre gate,
plan view. Panel (a): two leaves meet at the mitre point $M$, which points
upstream so that the higher water presses the leaves together rather than
apart. Each leaf pivots about its heel (quoin) post set into the wall. The
upstream water pressure grows with depth, and its resultant $P$ acts
perpendicular to the leaf at the centroid of the triangular pressure
distribution. Panel (b): the free-body diagram of one leaf. Three forces act:
the water thrust $P$ at mid-span, the heel reaction $R_H$ where the leaf bears
in its quoin, and the mitre reaction $R_M$ where one leaf presses on the
other. With only three non-parallel forces, the leaf is a three-force body and
its lines of action are concurrent: extending the lines of $P$ and $R_M$ to
their crossing fixes the direction of $R_H$ before any component equation is
written. The gate is closed and held by water alone, which is why a mitre gate
needs no latch against the head it faces and why it must be opened against
almost no load once the levels are equalised.}
\label{fig:vol03ch01:lock-gate}
\end{figure}


A lock gate holds back a river with no latch, no motor engaged, and no bolt in
tension against the head it faces, and the reason is a free-body diagram a canal
builder settled three centuries ago. A mitre gate is a pair of leaves that meet
at a point $M$ aimed upstream, so the higher water presses the two leaves
together against each other and against the walls, as in
figure~\ref{fig:vol03ch01:lock-gate}. Each leaf pivots in a quoin (heel) recess
set into the chamber wall. We size the forces on one leaf of a lock whose
chamber is $b = 5\,\si{\meter}$ wide, with the upstream water $h = 4\,\si{
\meter}$ deep and the downstream side drained, and the leaves mitred at
$\theta = 20^{\circ}$ to the line square across the chamber.

\textbf{The water thrust.} The pressure on the upstream face grows linearly with
depth, from zero at the surface to $\rho g h$ at the sill. Its resultant per
metre of gate width is the area of that triangular pressure profile, $\tfrac{1}{2}
\rho g h^2 = \tfrac{1}{2} \times 1000 \times 9.81 \times 4^2 = 78{,}480\,\text{N}$
per metre, acting at one-third of the depth up from the sill, the centroid of the
triangle. The leaf spans a length $L_{\text{leaf}} = (b/2)/\cos\theta = 2.5/\cos
20^{\circ} = 2.66\,\si{\meter}$ along the mitre, so the total thrust on one leaf
is
\[
P = \tfrac{1}{2}\rho g h^2 \times L_{\text{leaf}}
= 78{,}480 \times 2.66 = 208{,}800\,\text{N} \approx 209\,\text{kN},
\]
a distributed pressure reduced to one resultant $P$ acting perpendicular to the
leaf at its mid-length, the same reduction the sign panel used, now on a
triangular profile rather than a uniform one.

\textbf{The three-force leaf.} The leaf sees three forces: the water thrust $P$,
the heel reaction $R_H$ where it bears in its quoin, and the mitre reaction
$R_M$ where the two leaves press on each other. Three non-parallel forces on a
rigid body are concurrent, so the lines of $P$ and $R_M$ cross at a point through
which $R_H$ must also pass, fixing the heel-reaction direction before a component
is written. The mitre reaction is the thrust the two leaves exchange, directed
along the chamber. Resolving the leaf's balance with the mitre reaction along the
line the two leaves share gives the standard result
\[
R_M = \frac{P}{2\sin\theta} = \frac{208{,}800}{2\sin 20^{\circ}}
= \frac{208{,}800}{0.684} = 305{,}000\,\text{N} \approx 305\,\text{kN},
\]
half again the water thrust it carries. The heel reaction comes out the same
magnitude by symmetry of the force triangle, so the quoin and the mitre each
carry about $305\,\text{kN}$ into the wall and the neighbouring leaf.

\textbf{Why the mitre points upstream.} The sign of $R_M$ is the whole design.
Aimed upstream, the water thrust drives the leaves into each other and into
their quoins, and the mitre reaction is a compression the timber or steel of the
contact face carries easily. Reverse the mitre to point downstream and the same
thrust would try to force the leaves apart, turning the mitre contact into a
tension no simple bearing can hold and throwing the entire head onto whatever
latched the leaves together. The mitre gate needs no latch against its design
head precisely because the free-body diagram puts the closing force where the
structure is strongest, and it opens against almost nothing once the levels are
equalised and $P$ falls to zero. A lock keeper who has never drawn the diagram
still trusts it every time the paddles are lifted, and the shallow mitre angle
$\theta$ is chosen small enough to keep the leaf short and long enough to keep
$R_M = P/(2\sin\theta)$ from running away, the same divergence at shallow angle
that the hanging ropes and the climbing anchor showed, read here in stone and
water.

\section{Case study: the load path of a wall-mounted jib crane}
\pathtag{standard}

The worked examples isolated one body at a time. A real assembly forces
the harder discipline: carrying numbers from the load, through each body
and each joint in turn, to the one connection that governs. We work a
small wall-mounted jib crane end to end, the kind bolted to a workshop
wall to swing a few hundred kilograms off a delivery truck. The point is
the sequence: define the system, draw the free-body diagram of each body,
solve the joints in dependency order, and read off which connection is the
weakest link. No single equation in what follows is harder than the ladder
example. The difficulty, and the realism, is the bookkeeping across
bodies.

% Figure: the wall-mounted jib crane of the case study. The assembled
% scene on the left; the chosen-body decomposition (boom and tie rod) on
% the right. Numbers match the worked case study in the text.
\begin{figure}[htbp]
\centering
\begin{tikzpicture}[
  member/.style={draw=engblack, line width=1.6pt},
  tie/.style={draw=engblue, line width=1.4pt},
  load/.style={draw=engred, -{Stealth}, very thick},
  scale=1.0,
]
% ===== Left: the assembled crane =====
\begin{scope}[xshift=0cm]
  % wall
  \draw[draw=engblack, very thick] (0,0) -- (0,4.4);
  \foreach \y in {0.3,0.7,...,4.2} {\draw[draw=enggray, thin] (0,\y) -- (-0.28,\y-0.2);}
  % pin at A (boom root)
  \coordinate (A) at (0,1.0);
  \coordinate (B) at (4.0,1.0);
  \coordinate (C) at (0,3.4);
  % boom A==B (horizontal)
  \draw[member] (A) -- (B);
  % tie rod C==B
  \draw[tie] (C) -- (B);
  % pin markers
  \fill[engblack] (A) circle (0.08);
  \draw[draw=engblack, fill=white] (A) circle (0.13);
  \fill[engblack] (C) circle (0.08);
  \draw[draw=engblack, fill=white] (C) circle (0.13);
  \fill[engblack] (B) circle (0.08);
  \draw[draw=engblack, fill=white] (B) circle (0.13);
  % node labels
  \node[font=\scriptsize, anchor=east] at (A) {$A$};
  \node[font=\scriptsize, anchor=east] at (C) {$C$};
  \node[font=\scriptsize, anchor=south west] at (B) {$B$};
  % hoisted load at the boom tip
  \draw[load] (B) -- ++(0,-1.4) node[anchor=north, font=\small, engred] {$W$};
  % tie-rod angle marker at B
  \draw[draw=enggray, thin] (3.4,1.0) arc (180:149:0.6);
  \node[font=\scriptsize] at (3.25,1.22) {$\beta$};
  % dimension hints
  \node[font=\scriptsize, anchor=north] at (2.0,1.0) {boom $L$};
  \node[font=\scriptsize, anchor=east, rotate=31] at (2.0,2.5) {tie rod};
  \node[font=\small, anchor=north] at (2.0,-0.1) {(a) the assembly};
\end{scope}
% ===== Right: FBD of the boom AB as a single body =====
\begin{scope}[xshift=7.2cm]
  \coordinate (A2) at (0,1.0);
  \coordinate (B2) at (4.0,1.0);
  \draw[member] (A2) -- (B2);
  \fill[engblack] (A2) circle (0.06);
  \fill[engblack] (B2) circle (0.06);
  \node[font=\scriptsize, anchor=north east] at (A2) {$A$};
  \node[font=\scriptsize, anchor=north] at (B2) {$B$};
  % pin reaction components at A
  \draw[draw=engred, -{Stealth}, very thick] (A2) -- ++(1.5,0)
    node[anchor=south, font=\scriptsize, engred] {$A_x$};
  \draw[draw=engred, -{Stealth}, very thick] (A2) -- ++(0,1.5)
    node[anchor=east, font=\scriptsize, engred] {$A_y$};
  % tie-rod force at B: directed up and to the left along the rod (tension)
  \draw[draw=engblue, -{Stealth}, very thick] (B2) -- ++(-2.0,1.2)
    node[anchor=south, font=\scriptsize, engblue] {$T$};
  \draw[draw=enggray, thin] (3.4,1.0) arc (180:149:0.6);
  \node[font=\scriptsize] at (3.22,1.22) {$\beta$};
  % hoisted load down at B
  \draw[load] (B2) -- ++(0,-1.4) node[anchor=north, font=\small, engred] {$W$};
  \node[font=\small, anchor=north] at (2.0,-0.1) {(b) FBD of the boom};
\end{scope}
\end{tikzpicture}
\caption[A wall-mounted jib crane and the free-body diagram of its
boom]{The wall-mounted jib crane of the case study. (a) The assembly: a
horizontal boom $AB$ is pinned to the wall at $A$ and held up by a tie
rod from a higher wall anchor $C$ to the boom tip $B$, where a load $W$
hangs. (b) The free-body diagram of the boom treated as one body. The
pin at $A$ supplies two reaction components $A_x$ and $A_y$; the tie rod,
a two-force member, pulls along its own axis with tension $T$ at angle
$\beta$ above the boom; the load $W$ hangs from the tip. Summing moments
about $A$ kills both pin components and solves for $T$ directly, after
which force balance returns $A_x$ and $A_y$. The governing connection is
whichever of the pin, the tie-rod thread, or the wall anchor first
reaches its capacity under the resolved forces.}
\label{fig:vol03ch01:jib-crane}
\end{figure}


\textbf{The system.} A horizontal steel boom $AB$ of length $L = 2.0\,
\si{\meter}$ is pinned to the wall at its root $A$. A tie rod runs from a
higher wall anchor $C$ down to the boom tip $B$, holding the boom up. The
anchor $C$ sits $H = 1.2\,\si{\meter}$ directly above $A$ on the wall, so
the tie rod $CB$ makes an angle $\beta = \arctan(H/L) = \arctan(1.2/2.0) =
31.0^{\circ}$ with the boom. A load $W = 4{,}000\,\text{N}$ (about $408\,
\si{\kilogram}$, a generous truck pallet) hangs from the tip $B$. We take
the boom and tie rod as weightless against this load, an assumption we
check at the end. The bodies in the assembly are three: the boom, the tie
rod, and the wall anchor plate at $C$. The joints are three pins, at $A$,
$B$, and $C$.

\textbf{Step one: classify the members.} The tie rod $CB$ is pinned at
both ends, carries no transverse load, and has negligible weight, so it is
a two-force member: it carries pure axial force $T$, directed along $CB$.
That single fact removes two unknowns before any equation is written,
because the tie-rod force has a known direction. The boom $AB$ is loaded
at three points (the pin at $A$, the tie-rod pin at $B$, the hanging load
at $B$) and is not a two-force member; its root pin supplies two reaction
components $A_x$ and $A_y$.

\textbf{Step two: FBD of the boom and the governing moment equation.}
Figure~\ref{fig:vol03ch01:jib-crane}(b) is the boom isolated. The forces
on it are the pin reaction $(A_x, A_y)$ at $A$, the tie-rod tension $T$ at
$B$ pulling up and toward the wall at angle $\beta$ above the boom, and the
load $W$ hanging at $B$. Summing moments about $A$ kills both pin
components, because both pass through $A$:
\[
\sum M_A = 0: \quad (T\sin\beta)\,L - W\,L = 0
\quad\Longrightarrow\quad T\sin\beta = W.
\]
The vertical component of the tie-rod tension alone holds the load. So
\[
T = \frac{W}{\sin\beta} = \frac{4000}{\sin 31.0^{\circ}}
= \frac{4000}{0.515} = 7{,}767\,\text{N}.
\]
The tie rod carries nearly twice the hung load, the direct cost of its
shallow angle. A steeper tie rod (anchor higher up the wall) would lower
$T$; the geometry, fixed here by the available wall height, sets the
penalty.

\textbf{Step three: the boom-root pin.} With $T$ known, force balance on
the boom returns the pin reaction. The tie-rod tension at $B$ has
components $T\cos\beta$ toward the wall (negative $x$) and $T\sin\beta$ up:
\begin{align*}
\sum F_x = 0: & \quad A_x - T\cos\beta = 0
  \;\Rightarrow\; A_x = 7767\cos 31.0^{\circ} = 6{,}658\,\text{N}, \\
\sum F_y = 0: & \quad A_y + T\sin\beta - W = 0
  \;\Rightarrow\; A_y = W - T\sin\beta = 4000 - 4000 = 0.
\end{align*}
The vertical pin reaction vanishes, which is the geometry telling us the
tie rod carries the entire vertical load and the pin at $A$ carries only
the horizontal thrust. The pin reaction is therefore $A_x = 6{,}658\,
\text{N}$ horizontal, $A_y = 0$, a pure pull away from the wall of $6.66\,
\text{kN}$. The boom is in compression along its length at $6.66\,
\text{kN}$ (the horizontal squeeze between the wall pin and the tie-rod's
horizontal pull), a fact Volume V will need when it sizes the boom against
buckling.

% Figure: FBD of the upper wall anchor C of the jib crane, isolating the
% bracket plate so the bolt group load path is explicit.
\begin{figure}[htbp]
\centering
\begin{tikzpicture}[
  load/.style={draw=engred, -{Stealth}, very thick},
  tie/.style={draw=engblue, -{Stealth}, very thick},
  scale=1.0,
]
% wall
\draw[draw=engblack, very thick] (0,-0.4) -- (0,3.6);
\foreach \y in {-0.2,0.2,...,3.4} {\draw[draw=enggray, thin] (0,\y) -- (-0.26,\y-0.18);}
% anchor plate
\draw[draw=engblack, fill=engblue!10, thick] (0,0.6) rectangle (0.7,2.6);
\node[font=\scriptsize, anchor=west] at (0.75,2.3) {anchor plate};
% two bolts into the wall
\draw[draw=engblack, fill=engamber!50] (0.05,1.0) rectangle (0.35,1.25);
\draw[draw=engblack, fill=engamber!50] (0.05,1.95) rectangle (0.35,2.2);
\node[font=\scriptsize, anchor=west] at (0.4,1.12) {bolt 2};
\node[font=\scriptsize, anchor=west] at (0.4,2.07) {bolt 1};
% tie-rod pin at the plate (point C)
\coordinate (C) at (0.7,1.6);
\fill[engblack] (C) circle (0.06);
\node[font=\scriptsize, anchor=south west] at (C) {$C$};
% the tie rod pulls the plate down and to the right (reaction to T): the
% rod runs down to the boom tip, so on the anchor it pulls toward the tip
\draw[tie] (C) -- ++(2.0,-1.2) node[anchor=west, font=\small, engblue] {$T$};
% bolt tension on the top bolt (the prying couple) and bearing at bottom
\draw[load] (0.2,2.2) -- ++(-1.3,0) node[anchor=east, font=\scriptsize, engred] {$T_{b1}$};
\draw[load] (-1.0,1.0) -- ++(1.3,0) node[anchor=west, font=\scriptsize, engred] {$C_{b2}$};
\end{tikzpicture}
\caption[Free-body diagram of the jib-crane wall anchor]{The upper wall
anchor $C$ isolated as its own body. The tie rod pulls on the plate with
force $T$, directed toward the boom tip. The plate is held by a two-bolt
group: the moment of the off-axis pull about the bottom edge is resisted
as a couple, putting the top bolt in tension $T_{b1}$ and the bottom of
the plate in bearing $C_{b2}$ against the wall. Tracing the load through
the plate to the individual bolts is the step that decides whether the
anchor or the tie rod governs, and it is the step the original Hyatt
connection skipped.}
\label{fig:vol03ch01:jib-anchor}
\end{figure}


\textbf{Step four: the tie rod and the wall anchor.}
Figure~\ref{fig:vol03ch01:jib-anchor} isolates the anchor plate at $C$.
The tie rod pulls on the plate with the same $T = 7{,}767\,\text{N}$ it
carries (a two-force member transmits its axial force unchanged to both
ends), directed down the rod toward $B$, that is at $31.0^{\circ}$ below
the horizontal away from the wall. Resolved at the plate, this is a pull
of $T\cos\beta = 6{,}658\,\text{N}$ away from the wall and $T\sin\beta =
4{,}000\,\text{N}$ downward. The plate is bolted to the wall by two bolts.
The downward and outward pull, applied at the tie-rod pin offset from the
bolt line, is a combined direct pull plus a prying moment. Taking moments
about the lower bolt, with the bolts spaced $s = 0.15\,\si{\meter}$
vertically and the tie-rod pin centred between them at $d_v = s/2 = 0.075\,
\si{\meter}$ above the lower bolt and $e = 0.05\,\si{\meter}$ out from the
bolt plane, the top bolt tension is
\[
T_{b1} = \frac{(T\cos\beta)\,d_v + (T\sin\beta)\,e}{s}
= \frac{6658 \times 0.075 + 4000 \times 0.05}{0.15}
= \frac{499 + 200}{0.15} = 4{,}660\,\text{N},
\]
with the lower region of the plate in bearing against the wall. Each bolt
also shares the direct outward pull of $6{,}658\,\text{N}$, so the top
bolt's total demand combines $4{,}660\,\text{N}$ of prying tension with
about $3{,}329\,\text{N}$ of direct pull share, a combined load the bolt
must be checked against by an interaction rule rather than by either
component alone.

\textbf{Step five: the verdict.} Lay the connection demands side by side.
The boom-root pin at $A$ carries $6.66\,\text{kN}$ in shear across a single
large pin. The tie rod carries $7.77\,\text{kN}$ in pure tension along a
rod whose threaded ends are the weak section. The top anchor bolt carries
a combined $4.66\,\text{kN}$ of prying tension plus its share of the
direct pull. The governing connection is the one whose capacity is lowest
relative to its demand, and for a crane of this kind it is the tie-rod
thread or the anchor bolts, not the boom pin: a threaded rod loses a third
or more of its plain-shank area at the root of the thread, and a small
anchor bolt in a prying connection is working against both tension and the
amplifying lever arm $e$. The boom pin, sized as a structural pin, has the
most margin. A designer reads this load path and puts the attention, and
the inspection schedule, on the tie-rod ends and the anchor bolts, the
connections the numbers single out.

\textbf{The check we promised.} The weightless-member assumption. A
$2.0\,\si{\meter}$ steel boom of modest section masses perhaps $15\,\si{
\kilogram}$, a weight of $147\,\text{N}$ at its midpoint, against a hung
load of $4{,}000\,\text{N}$ at the tip. Its moment about $A$ is $147
\times 1.0 = 147\,\text{N}\cdot\text{m}$ against the load's $4000 \times
2.0 = 8{,}000\,\text{N}\cdot\text{m}$, under two percent. Neglecting the
boom weight changes $T$ by under two percent, below the other
uncertainties in the problem, so the assumption is defensible here; on a
longer or heavier boom it would not be, and the estimate is what tells the
two cases apart. This is the discipline the failure section names, run
forward: the assumption is stated, then checked with a number, not carried
on faith.

\begin{mastery}
\textbf{Why the free-body diagram alone cannot close an indeterminate
crane.} Suppose the jib crane above is built with a second tie rod, from a
second wall anchor $C'$, also reaching the boom tip $B$. The boom is now
held by two rods and the root pin. Count the unknowns: two rod tensions
$T_1, T_2$ and two pin components $A_x, A_y$, four in all. Count the
equations: the boom is one planar rigid body, so three, $\sum F_x = 0$,
$\sum F_y = 0$, $\sum M_A = 0$. Four unknowns, three equations: the system
is statically indeterminate to the first degree, and no choice of
moment-summing point recovers the missing equation, because the
moment-transfer rule of section 1.1 shows that all moment equations about
all points are linear combinations of the same three. The free-body
diagram has done everything it can do and the problem is still open. What
closes it is a compatibility condition: the two rods, sharing the same tip
point $B$, must stretch by amounts consistent with a single displacement
of $B$. That condition involves the rods' stiffnesses ($EA/\ell$ for each)
and brings the deformation of the bodies into a problem the rigid-body FBD
treated as undeformable. The stiffer rod takes the larger share, in
proportion the equilibrium equations never see. This is the boundary of
the present chapter: the FBD is necessary for every statics problem and
sufficient for determinate ones, and the count $r > 3e$ (reactions
exceeding the equilibrium equations, here $4 > 3$) is the alarm that
sufficiency has failed. Chapter 3 of this volume supplies the missing
compatibility equations and with them the force distribution in
indeterminate frames.
\end{mastery}

\section{Case study: a multi-leg sling lifting plan}
\pathtag{standard}

The jib-crane study carried numbers across the bodies of a fixed structure.
A lifting plan carries them across the geometry of a rig that the rigger
chooses, and the choices, sling angle and attachment placement, decide
whether a sling rated for the job is actually safe for the lift in front of
it. We work a single lift end to end: a machine of mass $1{,}200\,\si{
\kilogram}$ to be lifted by a two-leg sling from a single crane hook, and we
ask the question a competent rigger asks before the lift, which is not ``can
the sling hold the weight'' but ``can each leg hold its share at the angle
this lift forces.'' The arithmetic is the hanging-mass example of section
1.6.3 run with a rated capacity standing behind it.

% Figure: a two-leg sling lifting a load. (a) the rigging scene with the
% included angle and the leg angle from vertical; (b) the leg-tension
% multiplier as a function of the angle each leg makes with the vertical.
% Numbers match the multi-leg-sling case study in the text.
\begin{figure}[htbp]
\centering
\begin{tikzpicture}[
  leg/.style={draw=engblue, line width=1.4pt},
  ten/.style={draw=engred, -{Stealth}, very thick},
  wt/.style={draw=engblack, -{Stealth}, very thick},
  scale=1.0,
]
% ===== Left: the rigging scene =====
\begin{scope}[xshift=0cm]
  % hook / master link at top
  \coordinate (H) at (2.0,4.0);
  \draw[draw=engblack, very thick] (H) -- ++(0,0.5);
  \draw[draw=engblack, fill=white, thick] (2.0,4.65) circle (0.18);
  % attachment points on the load
  \coordinate (L) at (0.6,1.4);
  \coordinate (R) at (3.4,1.4);
  % sling legs
  \draw[leg] (H) -- (L);
  \draw[leg] (H) -- (R);
  % the load (a crate)
  \draw[draw=engblack, fill=engblue!10, thick] (0.4,0.4) rectangle (3.6,1.4);
  \node[font=\scriptsize] at (2.0,0.9) {load $W$};
  \fill[engblack] (L) circle (0.06);
  \fill[engblack] (R) circle (0.06);
  % vertical reference through the hook
  \draw[draw=enggray, thin, dashed] (H) -- ++(0,-2.6);
  % leg angle from vertical at the hook
  \draw[draw=enggray, thin] (2.0,3.3) arc (270:236:0.7);
  \node[font=\scriptsize] at (1.55,3.25) {$\theta$};
  % included angle marker
  \draw[draw=enggray, thin] (1.3,3.1) arc (236:304:0.7);
  \node[font=\scriptsize] at (2.0,2.95) {$2\theta$};
  \node[font=\small, anchor=north] at (2.0,0.1) {(a) two-leg sling};
\end{scope}
% ===== Right: tension multiplier vs leg angle =====
\begin{scope}[xshift=6.6cm, yshift=0.2cm]
  \begin{axis}[
    width=6.6cm, height=5.4cm,
    xlabel={leg angle from vertical $\theta$ (deg)},
    ylabel={leg tension $/$ (half load)},
    xmin=0, xmax=75, ymin=0.9, ymax=4.2,
    xtick={0,15,30,45,60,75},
    ytick={1,2,3,4},
    grid=both, grid style={enggray!25},
    label style={font=\scriptsize},
    tick label style={font=\scriptsize},
    every axis plot/.append style={line width=1.2pt},
  ]
    % multiplier = 1/cos(theta)
    \addplot[draw=engred, domain=0:74, samples=80] {1/cos(x)};
    \addplot[draw=enggray, dashed, domain=0:75] {2};
    \node[font=\scriptsize, anchor=west, engred] at (axis cs:6,1.7) {$1/\cos\theta$};
    \node[font=\scriptsize, anchor=south east, enggray] at (axis cs:60,2.0) {tension $= W$};
  \end{axis}
  \node[font=\small, anchor=north] at (3.0,-0.4) {(b) tension penalty};
\end{scope}
\end{tikzpicture}
\caption[A two-leg sling and the tension penalty of a wide angle]{(a) A
symmetric two-leg sling lifts a load $W$; each leg makes an angle $\theta$
with the vertical, so the legs subtend an included angle $2\theta$ at the
hook. Vertical balance of the master link gives $2T\cos\theta = W$, so each
leg carries $T = W/(2\cos\theta)$. (b) The leg tension, in units of half
the load, rises as $1/\cos\theta$: at $\theta = 0$ each leg carries exactly
half the load; at $\theta = 60^{\circ}$ each leg carries the full load $W$;
beyond that the penalty climbs without bound. The horizontal pull that the
wide angle introduces is the silent cost a rigger pays for spreading the
legs to clear the load, and it is why sling charts derate sharply below a
$45^{\circ}$ included angle.}
\label{fig:vol03ch01:sling-plan}
\end{figure}


\textbf{The system.} The machine weighs $W = 1{,}200 \times 9.81 = 11{,}772
\,\text{N}$, call it $11.8\,\text{kN}$. Two lifting lugs on the machine are
spaced $1.6\,\si{\meter}$ apart. A single crane hook sits above the load. The
two sling legs run from the hook down to the two lugs. The rigger has a
choice of sling length, and the choice fixes the angle $\theta$ each leg
makes with the vertical, because the legs must span from the lugs up to the
single hook. With the hook centred $1.5\,\si{\meter}$ above the lug line,
each leg makes an angle $\theta = \arctan(0.8/1.5) = 28.1^{\circ}$ with the
vertical. The bodies are two: the master link or hook ring where the legs
meet, and the load itself. The unknowns are the two leg tensions.

\textbf{Step one: the symmetric case and the leg tension.} Take the centre
of gravity at the geometric midpoint first, so the lift is symmetric and the
two legs share the weight equally. The FBD of the hook ring carries the
upward crane pull $P = W$, balanced by the downward pull of the two legs at
angle $\theta$. Vertical balance:
\[
2T\cos\theta = W \quad\Longrightarrow\quad
T = \frac{W}{2\cos\theta} = \frac{11{,}772}{2\cos 28.1^{\circ}}
= \frac{11{,}772}{1.764} = 6{,}674\,\text{N}.
\]
Each leg carries $6.67\,\text{kN}$, not the $5.89\,\text{kN}$ that an even
split of the weight alone would suggest. The extra is the sling-angle
penalty: the factor $1/\cos\theta = 1.13$ is the cost of the legs pulling
inward as well as up. Figure~\ref{fig:vol03ch01:sling-plan}(b) shows the
penalty climbing as the angle widens; at this $28^{\circ}$ it is a modest
thirteen percent, but a rigger forced to a $60^{\circ}$ leg angle (a short
sling on a wide load) would double the leg tension to the full weight.

% Figure: a two-leg sling on a load whose centre of gravity is offset from
% the geometric midpoint. The offset throws unequal tension into the two
% legs (non-equalising sling); the FBD of the load shows the moment balance
% that sets the split. Numbers match the case study.
\begin{figure}[htbp]
\centering
\begin{tikzpicture}[
  leg/.style={draw=engblue, line width=1.4pt},
  ten/.style={draw=engred, -{Stealth}, very thick},
  wt/.style={draw=engblack, -{Stealth}, very thick},
  scale=1.0,
]
% ===== Left: the offset-CG scene =====
\begin{scope}[xshift=0cm]
  \coordinate (H) at (2.3,4.0);
  \draw[draw=engblack, very thick] (H) -- ++(0,0.5);
  \draw[draw=engblack, fill=white, thick] (2.3,4.65) circle (0.18);
  \coordinate (L) at (0.6,1.4);
  \coordinate (R) at (4.0,1.4);
  \draw[leg] (H) -- (L);
  \draw[leg] (H) -- (R);
  \draw[draw=engblack, fill=engblue!10, thick] (0.4,0.4) rectangle (4.2,1.4);
  \fill[engblack] (L) circle (0.06);
  \fill[engblack] (R) circle (0.06);
  \node[font=\scriptsize, anchor=south] at (L) {$L$};
  \node[font=\scriptsize, anchor=south] at (R) {$R$};
  % centre of gravity, offset toward R
  \coordinate (G) at (2.8,0.9);
  \fill[engred] (G) circle (0.07);
  \node[font=\scriptsize, anchor=north, engred] at (G) {$G$};
  \draw[wt] (G) -- ++(0,-1.0) node[anchor=north, font=\scriptsize] {$W$};
  % geometric midpoint marker
  \draw[draw=enggray, thin, dashed] (2.3,1.4) -- (2.3,-0.1);
  \node[font=\scriptsize, anchor=north, enggray] at (2.3,-0.1) {midpoint};
  \node[font=\small, anchor=north] at (2.3,-0.7) {(a) offset centre of gravity};
\end{scope}
% ===== Right: FBD of the load (moment balance) =====
\begin{scope}[xshift=7.0cm]
  \coordinate (L2) at (0,1.4);
  \coordinate (R2) at (3.4,1.4);
  \draw[draw=engblack, fill=engblue!10, thick] (-0.2,0.4) rectangle (3.6,1.4);
  \fill[engblack] (L2) circle (0.06);
  \fill[engblack] (R2) circle (0.06);
  \node[font=\scriptsize, anchor=north east] at (L2) {$L$};
  \node[font=\scriptsize, anchor=north west] at (R2) {$R$};
  % leg tensions (vertical components shown for the moment balance)
  \draw[ten] (L2) -- ++(-0.6,1.5) node[anchor=south, font=\scriptsize, engred] {$T_L$};
  \draw[ten] (R2) -- ++(0.6,1.5) node[anchor=south, font=\scriptsize, engred] {$T_R$};
  % weight at offset G
  \coordinate (G2) at (2.2,0.9);
  \fill[engred] (G2) circle (0.06);
  \draw[wt] (G2) -- ++(0,-1.0) node[anchor=north, font=\scriptsize] {$W$};
  % distance markers
  \draw[draw=enggray, thin, <->] (0,-0.1) -- (2.2,-0.1);
  \node[font=\scriptsize, anchor=north, enggray] at (1.1,-0.1) {$a$};
  \draw[draw=enggray, thin, <->] (2.2,-0.55) -- (3.4,-0.55);
  \node[font=\scriptsize, anchor=north, enggray] at (2.8,-0.55) {$b$};
  \node[font=\small, anchor=north] at (1.6,-1.0) {(b) FBD of the load};
\end{scope}
\end{tikzpicture}
\caption[A two-leg sling on a load with an offset centre of gravity]{(a) A
crate slung from two legs, with its centre of gravity $G$ displaced toward
the right attachment. (b) The free-body diagram of the load: each leg
applies its tension at its own attachment, and the weight $W$ acts at $G$, a
distance $a$ from the left attachment and $b$ from the right. Moments about
the right attachment give the left leg's vertical share, $T_{L,v} = Wb/(a+b)$,
and about the left attachment the right leg's, $T_{R,v} = Wa/(a+b)$. The leg
nearer the centre of gravity takes the larger share. A non-equalising sling
locks each leg to its computed length and lets the split stand; an
equalising sling (a single line reeved through a ring at the hook) forces
$T_L = T_R$ and instead tilts the load until the geometry restores moment
balance. The choice governs whether the load hangs level or hangs safe.}
\label{fig:vol03ch01:sling-cg-offset}
\end{figure}


\textbf{Step two: the centre-of-gravity offset.} Real machines do not carry
their weight at the geometric midpoint. Suppose this machine's centre of
gravity sits $0.2\,\si{\meter}$ off centre toward the right lug, so it is a
distance $a = 1.0\,\si{\meter}$ from the left lug and $b = 0.6\,\si{\meter}$
from the right, with $a + b = 1.6\,\si{\meter}$. The FBD of the load,
figure~\ref{fig:vol03ch01:sling-cg-offset}(b), takes moments about each lug
in turn to split the vertical load. Moments about the right lug give the
left leg's vertical share, and about the left lug the right leg's:
\[
T_{L,v} = W\frac{b}{a+b} = 11{,}772 \times \frac{0.6}{1.6} = 4{,}415\,\text{N},
\qquad
T_{R,v} = W\frac{a}{a+b} = 11{,}772 \times \frac{1.0}{1.6} = 7{,}357\,\text{N}.
\]
The leg nearer the centre of gravity, the right leg, carries the larger
share, $7.36\,\text{kN}$ of vertical load against the left leg's $4.42\,
\text{kN}$. Dividing each by $\cos\theta$ to recover the axial leg force, the
right leg carries $7357/\cos 28.1^{\circ} = 8{,}340\,\text{N}$ and the left
$4415/\cos 28.1^{\circ} = 5{,}005\,\text{N}$. The offset has thrown the
right leg up to $8.34\,\text{kN}$, a quarter more than the $6.67\,\text{kN}$
the symmetric calculation reported, and the symmetric calculation is the one
a careless plan would have used.

\textbf{Step three: equalising versus non-equalising.} The split above
assumes a non-equalising sling, two separate legs of fixed length, each
locked to carry whatever its geometry assigns. An equalising sling reeves a
single line through a ring or block at the hook, so the line is free to slide
until the tension is the same in both legs, $T_L = T_R$. Equal tension does
not mean equal vertical share, though: with the centre of gravity off centre,
forcing $T_L = T_R$ makes the load tilt until the two legs reach unequal
angles whose vertical components restore moment balance about the centre of
gravity. The equalising sling trades a known tension split for an unknown
tilt, and a load that must hang level (a vessel of liquid, a machine with a
spirit-level datum) cannot use one. The non-equalising sling holds the load
level and accepts the unequal tension, which is why the lift plan must size
the heavy leg, not the average leg.

\textbf{Step four: the capacity verdict.} A two-leg chain sling is sold with
a rated working load limit per leg, derated for the included angle, on the
tag the standard requires. Suppose this sling's legs are each rated at
$8\,\text{kN}$ at a vertical pull, with the manufacturer's chart derating to
about $7.5\,\text{kN}$ per leg at a $28^{\circ}$ leg angle. The symmetric
calculation, $6.67\,\text{kN}$ per leg, clears the $7.5\,\text{kN}$ rating
with margin and would pass the lift. The honest calculation, which carries
the centre-of-gravity offset, puts $8.34\,\text{kN}$ into the right leg,
which exceeds the $7.5\,\text{kN}$ derated rating. The verdict reverses: the
lift that passed on the symmetric number fails on the real one. The fix is
the rigger's to choose, a longer sling to lower $\theta$ and shave the angle
penalty, a higher-capacity sling, or a moved hook so the upper point sits
over the centre of gravity and restores a near-even split. The case is the
chapter's discipline in miniature: the free-body diagram that answers the
symmetric question does not answer the offset question, and a lift planned on
the wrong diagram is overloaded with a confident face, the same sentence the
failure section will write at building scale.

\begin{mastery}
\textbf{The funicular polygon and graphic statics.} Before cheap
computation, engineers solved force systems by drawing them, and the drawing
method survives because it makes the load path visible in a way the algebra
hides. Two diagrams do the work. The \engterm{force polygon} draws every
force tip to tail, scaled by length and oriented by direction; if the forces
balance, the polygon closes on itself. The \engterm{funicular polygon} (the
``funicular,'' from the Latin for rope) draws the line a weightless rope
would take if the same loads hung from it, and its segments are parallel to
the rays of the force polygon drawn from a chosen pole. The two diagrams are
reciprocal: a closed force polygon together with a closed funicular is the
graphic statement of $\sum\vec{F} = 0$ and $\sum\vec{M} = 0$ at once. The
method's power is that the funicular's shape is the line of thrust for the
load set, so an arch designed to follow its own funicular carries its load in
pure compression with no bending, the principle behind a well-shaped masonry
arch and the inverted-catenary forms used from Gaud\'i's models to thin-shell
roofs. Reading a funicular backward diagnoses a structure: where the line of
thrust strays outside the masonry, the arch wants to hinge and crack, and the
crack pattern of a failing arch traces the funicular the loads actually
found. The cable examples of this chapter are funiculars seen directly, since
a loaded cable hangs in exactly the funicular shape for its loads; an arch is
that same shape flipped to put the rope's tension into the stone's
compression. Graphic statics is studied today less as a calculation tool than
as a way of seeing, because a designer who can sketch a funicular can judge a
load path before any number is computed.
\end{mastery}

\section{Case study: the tipping stability of a tower crane}
\pathtag{standard}

The jib-crane study followed forces through a fixed structure; the sling
study followed them through a rig the operator chooses. A tower crane adds a
third question that neither raised, because both were bolted or hung from
something stronger than themselves: whether the machine stays upright at all.
A tower crane is a tall mast carrying a long horizontal jib, with a
counterweight on a short counter-jib behind the mast to balance the load
hung out front. The machine does not fail by a member breaking so much as by
the whole crane rotating about an edge of its base, and the calculation that
governs is a moment balance about that edge. We work the stability of a small
tower crane end to end: the forward tipping case under load, the backward
tipping case with the load released, the ballast that makes both safe, and
the out-of-service wind case that often governs the design.

% Figure: the tower crane of the third case study. The assembled machine
% on the left with the two tipping cases marked; the simplified stability
% free-body diagram on the right, reduced to the forces and lever arms
% that enter the overturning balance. Numbers match the worked case study.
\begin{figure}[htbp]
\centering
\begin{tikzpicture}[
  mast/.style={draw=engblack, line width=1.8pt},
  jib/.style={draw=engblack, line width=1.4pt},
  load/.style={draw=engred, -{Stealth}, very thick},
  ballast/.style={draw=engblue, -{Stealth}, very thick},
  dim/.style={draw=enggray, thin, {Stealth}-{Stealth}},
  scale=1.0,
]
% ===== Left: the assembled tower crane =====
\begin{scope}[xshift=0cm]
  % ground
  \draw[draw=engblack, very thick] (-2.6,0) -- (2.8,0);
  \foreach \x in {-2.4,-1.9,...,2.6} {\draw[draw=enggray, thin] (\x,0) -- (\x-0.2,-0.25);}
  % base / foundation block
  \draw[draw=engblack, fill=enggray!18, thick] (-0.9,0) rectangle (0.9,0.5);
  % mast
  \draw[mast] (-0.35,0.5) -- (-0.35,4.0);
  \draw[mast] (0.35,0.5) -- (0.35,4.0);
  \foreach \y in {0.8,1.3,...,3.8} {\draw[draw=enggray, thin] (-0.35,\y) -- (0.35,\y-0.3);}
  % jib (working) to the right, counter-jib to the left
  \draw[jib] (0,4.0) -- (2.4,4.0);
  \draw[jib] (0,4.0) -- (-1.5,4.0);
  % A-frame apex
  \draw[jib] (-0.35,4.0) -- (0,4.6) -- (0.35,4.0);
  \draw[draw=enggray, thin] (0,4.6) -- (2.2,4.05);
  \draw[draw=enggray, thin] (0,4.6) -- (-1.35,4.05);
  % counterweight slab
  \draw[draw=engblack, fill=engblue!18, thick] (-1.5,3.7) rectangle (-1.0,4.0);
  \node[font=\scriptsize, anchor=east] at (-1.5,3.85) {$W_c$};
  % hoist load near jib tip
  \draw[load] (2.2,4.0) -- ++(0,-1.5) node[anchor=north, font=\small, engred] {$W_L$};
  % tipping edges (front and rear base corners)
  \fill[engred] (0.9,0) circle (0.07);
  \node[font=\scriptsize, anchor=north west, engred] at (0.9,0) {$O_f$};
  \fill[engblue] (-0.9,0) circle (0.07);
  \node[font=\scriptsize, anchor=north east, engblue] at (-0.9,0) {$O_b$};
  % radius hints
  \draw[dim] (0,4.25) -- (2.2,4.25);
  \node[font=\scriptsize, anchor=south] at (1.1,4.25) {$R_L$};
  \draw[dim] (0,4.25) -- (-1.25,4.25);
  \node[font=\scriptsize, anchor=south] at (-0.6,4.25) {$R_c$};
  \node[font=\small, anchor=north] at (0.0,-0.45) {(a) the machine};
\end{scope}
% ===== Right: simplified stability FBD about the front tipping edge =====
\begin{scope}[xshift=6.6cm]
  % ground line
  \draw[draw=engblack, very thick] (-2.4,0) -- (2.6,0);
  \foreach \x in {-2.2,-1.7,...,2.4} {\draw[draw=enggray, thin] (\x,0) -- (\x-0.2,-0.25);}
  % schematic crane reduced to a horizontal arm at working height
  \coordinate (Of) at (0.9,0);
  \coordinate (Ob) at (-0.9,0);
  \draw[mast] (0,0) -- (0,3.4);
  \draw[jib] (-1.4,3.4) -- (2.2,3.4);
  % tipping edges
  \fill[engred] (Of) circle (0.07);
  \node[font=\scriptsize, anchor=north, engred] at (Of) {$O_f$};
  \fill[engblue] (Ob) circle (0.07);
  \node[font=\scriptsize, anchor=north, engblue] at (Ob) {$O_b$};
  % load down at jib tip
  \draw[load] (2.2,3.4) -- ++(0,-1.4) node[anchor=west, font=\scriptsize, engred] {$W_L$};
  % counterweight down at counter-jib
  \draw[ballast] (-1.4,3.4) -- ++(0,-1.4) node[anchor=east, font=\scriptsize, engblue] {$W_c$};
  % machine self-weight down at mast centreline
  \draw[draw=engamber, -{Stealth}, very thick] (0,3.4) -- ++(0,-2.0)
    node[anchor=west, font=\scriptsize, engamber] {$W_m$};
  % lever arms from the front edge O_f, drawn along the ground
  \draw[dim] (Of) -- (2.2,0);
  \node[font=\scriptsize, anchor=north] at (1.55,-0.05) {$a_L$};
  \draw[dim] (Of) -- (-1.4,0);
  \node[font=\scriptsize, anchor=north] at (-0.25,-0.05) {$a_c$};
  \node[font=\small, anchor=north] at (0.1,-0.5) {(b) overturning balance about $O_f$};
\end{scope}
\end{tikzpicture}
\caption[A tower crane and its overturning free-body diagram]{The tower
crane of the case study. (a) The machine: a mast carries a working jib to
the right and a shorter counter-jib to the left, with a counterweight
$W_c$ at radius $R_c$ balancing the hoisted load $W_L$ at radius $R_L$.
The base rests on two tipping edges, the front edge $O_f$ (toward the
load) and the back edge $O_b$ (toward the counterweight). (b) The
stability free-body diagram, reduced to the three vertical forces and
their horizontal lever arms about a tipping edge. Forward tipping turns
about $O_f$: the load moment $W_L a_L$ drives the overturn and the
counterweight and machine weight resist it. Backward tipping turns about
$O_b$ with the load removed: the counterweight now drives the overturn
and must itself be resisted. Stability is the statement that the
resisting moment exceeds the driving moment about each edge, with a
margin the standard sets.}
\label{fig:vol03ch01:tower-crane}
\end{figure}


\textbf{The system.} Figure~\ref{fig:vol03ch01:tower-crane} shows the
machine. A mast rises to a horizontal jib. A load $W_L$ hangs from a trolley
on the working jib at a radius $R_L$ from the mast centreline; a
counterweight $W_c$ sits on the counter-jib at a radius $R_c$ behind the
mast. The crane stands on a square base of width $2c$, so the two tipping
edges, the front edge $O_f$ toward the load and the back edge $O_b$ toward
the counterweight, are each a distance $c$ from the mast centreline. The
machine's own weight $W_m$ acts down the mast centreline, midway between the
edges. We take, for this crane, a base half-width $c = 2.0\,\si{\meter}$, a
machine self-weight $W_m = 120\,\text{kN}$, a maximum load $W_L = 40\,\text{kN}$
at a working radius $R_L = 18\,\si{\meter}$, and a counterweight radius $R_c
= 5\,\si{\meter}$. The unknown is the counterweight $W_c$ that keeps the crane
upright in every case, and the question is what value of $W_c$ the two tipping
balances and the wind case together demand.

\textbf{Step one: forward tipping under load.} The dangerous case in service
is the loaded crane tipping forward, over the front edge $O_f$, dragged down
by the load on the long jib. Take moments about $O_f$. The load drives the
overturn: its lever arm about $O_f$ is its radius less the half-width,
$R_L - c = 18 - 2 = 16\,\si{\meter}$, so its overturning moment is
\[
M_{\text{over}} = W_L\,(R_L - c) = 40 \times 16 = 640\,\text{kN}\cdot\text{m}.
\]
Two weights resist. The machine self-weight acts at the mast centreline, a
distance $c = 2.0\,\si{\meter}$ behind $O_f$, and the counterweight acts at
$R_c + c = 5 + 2 = 7\,\si{\meter}$ behind $O_f$. Their resisting moment is
\[
M_{\text{resist}} = W_m\,c + W_c\,(R_c + c) = 120 \times 2 + W_c \times 7
= 240 + 7 W_c \quad(\text{kN}\cdot\text{m}).
\]
A crane standard does not ask merely that resistance equal overturning; it
asks for a margin. Taking a representative stability factor of $1.5$ on the
load moment (the value depends on the governing standard and the load case),
the requirement is $M_{\text{resist}} \ge 1.5\,M_{\text{over}}$:
\[
240 + 7 W_c \ge 1.5 \times 640 = 960
\quad\Longrightarrow\quad
W_c \ge \frac{960 - 240}{7} = 103\,\text{kN}.
\]
Forward stability alone calls for a counterweight of at least $103\,\text{kN}$.

\textbf{Step two: backward tipping with the load released.} A crane is not
loaded all the time, and the counterweight sized to resist the load does not
disappear when the load does. With the hook empty, the same counterweight now
tends to tip the crane backward, over the rear edge $O_b$. Take moments about
$O_b$. The counterweight drives this overturn: its lever arm about $O_b$ is
$R_c - c = 5 - 2 = 3\,\si{\meter}$, so
\[
M_{\text{over,back}} = W_c\,(R_c - c) = 3 W_c \quad(\text{kN}\cdot\text{m}).
\]
The machine self-weight resists, acting $c = 2.0\,\si{\meter}$ in front of
$O_b$, with moment $W_m\,c = 120 \times 2 = 240\,\text{kN}\cdot\text{m}$.
There is no load to help, by assumption, and a prudent backward check takes
the hook empty. With the same $1.5$ margin applied to the backward overturn,
\[
W_m\,c \ge 1.5\,W_c\,(R_c - c)
\quad\Longrightarrow\quad
240 \ge 1.5 \times 3 W_c = 4.5 W_c
\quad\Longrightarrow\quad
W_c \le \frac{240}{4.5} = 53\,\text{kN}.
\]
The backward case sets an \emph{upper} bound on the counterweight, and it
contradicts the forward case. Forward stability wants $W_c \ge 103\,\text{kN}$;
backward stability wants $W_c \le 53\,\text{kN}$. No single counterweight
satisfies both at once with the geometry as given, which is the calculation
telling the designer something physical: a crane cannot carry a heavy load on
a long jib with a base this narrow and a counter-jib this short. The base, the
counter-jib radius, or both must grow.

\textbf{Step three: sizing the geometry.} Widen the base to $c = 3.0\,\si{
\meter}$ and lengthen the counter-jib so the counterweight sits at $R_c = 6\,
\si{\meter}$, and rerun both balances. Forward, about $O_f$ with load lever arm
$R_L - c = 18 - 3 = 15\,\si{\meter}$:
\[
W_m\,c + W_c\,(R_c + c) \ge 1.5\,W_L\,(R_L - c),
\quad
360 + 9 W_c \ge 1.5 \times 40 \times 15 = 900,
\quad
W_c \ge 60\,\text{kN}.
\]
Backward, about $O_b$ with counterweight lever arm $R_c - c = 6 - 3 = 3\,\si{
\meter}$ and machine arm $c = 3.0\,\si{\meter}$:
\[
W_m\,c \ge 1.5\,W_c\,(R_c - c),
\quad
360 \ge 4.5 W_c,
\quad
W_c \le 80\,\text{kN}.
\]
The two bounds now overlap: any counterweight between $60$ and $80\,\text{kN}$
keeps the crane upright in both cases with the required margin. A designer
picks a value inside the window, say $W_c = 70\,\text{kN}$, which leaves room
on both sides for the inevitable departures from the nominal numbers. The
narrow first geometry had an empty window; widening the base and reaching the
counterweight out opened one. The window between the two bounds is the design
freedom, and a crane whose window has closed is a crane that has been asked to
do more than its base allows.

\textbf{Step four: the out-of-service wind case.} A tower crane spends most of
its life not lifting. Parked overnight or through a storm, it is left to
weathervane, its jib free to swing downwind so the wind blows along the jib
rather than across it, which is the lightest case and the reason the slewing
brake is released out of service. The wind still presses on the exposed mast,
jib, and counter-jib, and the resulting horizontal force, applied high up,
produces an overturning moment about the base edge that can exceed the
in-service load case. Model the wind as a horizontal force $F_w$ acting at an
effective height $h_w$ above the base. Take a storm wind giving an effective
$F_w = 30\,\text{kN}$ at $h_w = 20\,\si{\meter}$ (the area-weighted centroid of
the structure's drag), so its overturning moment about the downwind base edge
is
\[
M_{w} = F_w\,h_w = 30 \times 20 = 600\,\text{kN}\cdot\text{m}.
\]
The resisting moment is the dead weight of the whole machine, counterweight
included, about that edge: with $W_m = 120\,\text{kN}$ and $W_c = 70\,\text{kN}$
and the wind blowing toward the front edge $O_f$, the counterweight at $R_c +
c = 9\,\si{\meter}$ and the machine at $c = 3\,\si{\meter}$ give
\[
M_{\text{resist}} = W_m\,c + W_c\,(R_c + c) = 120 \times 3 + 70 \times 9 = 990\,
\text{kN}\cdot\text{m},
\]
a ratio $M_{\text{resist}}/M_w = 990/600 = 1.65$, clearing a typical
out-of-service stability margin of about $1.2$ to $1.5$. Had the wind moment
come out larger than the dead-weight resistance allows, the remedy is not more
counterweight (which the backward case caps) but a ballasted base, ground
anchors, or tying the mast into an adjacent structure, which is why tall cranes
are bolted to the buildings they erect. The wind case is run last and often
governs, because the force that decides whether a parked crane survives a storm
is not the one it was built to lift.

\textbf{The verdict.} The stability of a tower crane is three moment balances,
not one. Forward tipping under load sets a floor on the counterweight; backward
tipping with the hook empty sets a ceiling; the out-of-service wind case sets a
demand on the dead-weight resistance that neither lifting case sees. The first
geometry failed because its two bounds did not overlap, and the failure showed
up in the arithmetic before it could show up on a site. The discipline is the
chapter's, scaled to a machine that weighs more than the loads it carries: draw
the body, choose the edge to take moments about, and let the sign of the margin,
not the look of the machine, decide whether it stands.

\begin{mastery}
\textbf{Tipping and sliding as competing failure modes.} The pushed block of
section 1.6.10 and the tower crane above share one structure: a body held by a
contact can fail by sliding, governed by a force inequality, or by tipping,
governed by a moment inequality, and the FBD must test both because either can
govern. Write them side by side for a body of weight $mg$ and base half-width
$c$, pushed by a horizontal force $P$ at height $h$ on a floor of coefficient
$\mu_s$. Sliding impends at $P_{\text{slide}} = \mu_s mg$, a pure force
balance that ignores $h$. Tipping impends at $P_{\text{tip}} = mg\,c/h$, a
pure moment balance about the leading edge that ignores $\mu_s$. The two are
equal at the critical push height
\[
h^\ast = \frac{c}{\mu_s},
\]
above which the body tips first and below which it slides first. The
competition is general. A high, narrow body on a grippy floor tips (a
refrigerator on rubber feet, a tall crane); a low, wide body on a slick floor
slides (a flat crate on a polished hall). The design moves are opposite: to
stop sliding, raise $\mu_s$ or lower $P$; to stop tipping, widen the base
(raise $c$), lower the centre of mass, or lower the push (reduce $h$). A fix
for one mode can worsen the other, since a grippier floor that cures sliding
makes tipping the new governor. The honest analysis names which mode governs
the geometry in hand before reaching for a cure, because the cheap remedy for
one failure is no remedy at all for the other, and a body braced against the
wrong mode fails in the mode that was never checked.
\end{mastery}

\begin{mastery}
\textbf{The centre of pressure: where a contact reaction really acts.} The
tipping analyses drew the floor's reaction as a single normal force $N$ at a
point, and slid that point to the leading edge at the tipping threshold. The
single point is a resultant, and its location is the \engterm{centre of
pressure}, the analogue for a contact of the centroid of a distributed load. A
flat foot on a floor presses through a distributed pressure $p(x)$ over its
contact area; the reaction the FBD draws is that distribution's resultant, of
magnitude $N = \int p(x)\,dx$ acting at the centre of pressure
\[
x_{\text{cp}} = \frac{\int x\,p(x)\,dx}{\int p(x)\,dx},
\]
the pressure-weighted mean position. Under a centred load the pressure is
symmetric and the centre of pressure sits under the load; add an overturning
moment and the pressure crowds toward the leading edge, dragging $x_{\text{cp}}$
with it. The tipping condition has a clean reading in these terms: the body is
stable exactly while the centre of pressure stays inside the contact patch, and
it tips at the instant $x_{\text{cp}}$ reaches the edge, because the floor
cannot push the reaction past the material it bears on. A floor cannot pull
down, so the pressure cannot go negative, and once the moment demands a centre
of pressure outside the patch the contact lifts on the trailing side and the
body rotates. This is why the tipping calculation moves $N$ to the edge: not by
convention but because the edge is the farthest the centre of pressure can
migrate before the contact opens. The same idea sizes the bearing area of a
footing (keep the resultant inside the middle third of the base and the whole
base stays in compression, the classic no-tension rule for masonry), sets the
stability of a parked aircraft on its gear, and explains why a forklift's rated
capacity falls with load-centre distance: the load's centre of pressure marches
toward the front axle and the counterweight's job is to hold the machine's
overall centre of pressure inside the wheelbase. The contact reaction is never
really a point force; the point is where its distribution balances, and the
tipping edge is where that balance point runs out of floor.
\end{mastery}

\section{Failure: free-body diagrams that hid the real load path}
\pathtag{core}

Every chapter in this work has a failure section. We begin Volume III
with the failure mode that the chapter's own subject can hide: an FBD
drawn for the original design, never redrawn when the design changed.

% Figure: the Hyatt Regency hanger-rod connection. Original single-rod
% load path (left) versus the as-built two-rod load path (right) that
% doubled the upper connection load.
\begin{figure}[htbp]
\centering
\begin{tikzpicture}[
  rod/.style={draw=enggray, very thick},
  beam/.style={draw=engblack, fill=engblue!10, thick},
  load/.style={draw=engred, -{Stealth}, very thick},
  scale=1.0,
]
% ===== Left: original design, single continuous rod =====
\begin{scope}[xshift=0cm]
  % ceiling
  \draw[draw=engblack, very thick] (-0.4,5.4) -- (3.0,5.4);
  \foreach \x in {-0.2,0.2,...,2.8} {\draw[draw=enggray, thin] (\x,5.4) -- (\x-0.2,5.65);}
  % single rod from ceiling all the way down
  \draw[rod] (1.3,5.4) -- (1.3,0.6);
  % upper (4th-floor) box beam at y=3.6
  \draw[beam] (0.2,3.4) rectangle (2.4,3.85);
  \node[font=\scriptsize] at (1.3,3.62) {4th-floor beam};
  % nut at upper beam
  \draw[draw=engblack, fill=engamber!50] (1.15,3.4) rectangle (1.45,3.55);
  % lower (2nd-floor) box beam at y=1.0
  \draw[beam] (0.2,0.8) rectangle (2.4,1.25);
  \node[font=\scriptsize] at (1.3,1.02) {2nd-floor beam};
  \draw[draw=engblack, fill=engamber!50] (1.15,0.6) rectangle (1.45,0.75);
  % loads
  \draw[load] (0.5,3.85) -- (0.5,4.55) node[anchor=south, font=\scriptsize, engred] {$W_4$};
  \draw[load] (0.5,1.25) -- (0.5,1.95) node[anchor=south, font=\scriptsize, engred] {$W_2$};
  % connection load callout
  \node[font=\scriptsize, anchor=west, engblue] at (2.6,3.6) {carries $W_4$};
  \node[font=\small, anchor=north] at (1.3,0.0) {(a) original: single rod};
\end{scope}
% ===== Right: as-built, two rods, doubled connection load =====
\begin{scope}[xshift=7.4cm]
  \draw[draw=engblack, very thick] (-0.4,5.4) -- (3.0,5.4);
  \foreach \x in {-0.2,0.2,...,2.8} {\draw[draw=enggray, thin] (\x,5.4) -- (\x-0.2,5.65);}
  % upper rod: ceiling to 4th-floor beam
  \draw[rod] (1.3,5.4) -- (1.3,3.85);
  % lower rod: 4th-floor beam to 2nd-floor beam (OFFSET to show it is separate)
  \draw[rod] (1.7,3.4) -- (1.7,1.25);
  % upper beam
  \draw[beam] (0.2,3.4) rectangle (2.4,3.85);
  \node[font=\scriptsize] at (0.9,3.62) {4th-floor};
  \draw[draw=engblack, fill=engamber!50] (1.15,3.4) rectangle (1.45,3.55);
  \draw[draw=engblack, fill=engamber!50] (1.55,3.4) rectangle (1.85,3.55);
  % lower beam
  \draw[beam] (0.2,0.8) rectangle (2.4,1.25);
  \node[font=\scriptsize] at (0.9,1.02) {2nd-floor};
  \draw[draw=engblack, fill=engamber!50] (1.55,1.25) rectangle (1.85,1.4);
  % loads
  \draw[load] (0.5,3.85) -- (0.5,4.55) node[anchor=south, font=\scriptsize, engred] {$W_4$};
  \draw[load] (2.2,1.25) -- (2.2,1.95) node[anchor=south, font=\scriptsize, engred] {$W_2$};
  % connection load callout: now W4 + W2
  \node[font=\scriptsize, anchor=west, engred] at (2.6,3.6) {carries $W_4 + W_2$};
  \node[font=\small, anchor=north] at (1.3,0.0) {(b) as-built: two rods};
\end{scope}
\end{tikzpicture}
\caption[Hyatt Regency: the load path that the redrawn FBD would have
caught]{The hanger-rod connection at the heart of the 1981 Hyatt
Regency walkway collapse \cite{acc:nbs-hyatt-1982}. (a) In the original
design a single continuous rod runs from the ceiling through both
walkway box-beams; the upper connection carries only the fourth-floor
walkway load $W_4$, because the second-floor load $W_2$ reaches the
ceiling through the rod below the upper nut. (b) The as-built design,
changed during fabrication to ease assembly, used two rod segments: the
lower walkway now hangs from the upper walkway's box-beam, so the upper
connection carries $W_4 + W_2$, roughly double the design intent. The
change looks local on a drawing. The free-body diagram of the upper
connection, redrawn for the as-built geometry, makes the doubled load
visible at once. That diagram was never drawn.}
\label{fig:vol03ch01:hyatt-loadpath}
\end{figure}


On 17 July 1981, two suspended walkways in the atrium of the Hyatt
Regency Hotel in Kansas City, Missouri, collapsed during a tea
dance. The fourth-floor walkway fell onto the second-floor walkway;
both then fell to the atrium floor; 114 people were killed and more
than 200 injured. The National Bureau of Standards investigation by
Pfrang and Marshall, published as NBS Building Science Series 143 in
1982, identified the cause as a connection failure at the
fourth-floor walkway's hanger-rod attachment
\cite{acc:nbs-hyatt-1982}.

The original design used a single continuous threaded rod from the
atrium ceiling, passing through the box-beam ends of both the
fourth-floor and second-floor walkways. The FBD of the upper box-beam
connection, drawn for that original design, shows the load of one
walkway pulling down through the rod from above and the load of one
walkway pulling down through the box-beam itself, with the connection
carrying the weight of the fourth-floor walkway only. The second-floor
walkway's weight reaches the ceiling through its own segment of the
rod, hanging from the upper walkway's nut, but the nut at the upper
walkway transmits the second-floor walkway's load to the upper rod,
not to the upper box-beam.

During fabrication, the design was modified. Threading a single
continuous rod through both walkways was difficult: the second-floor
nut would have had to be threaded along the rod's full length. The
fabricator proposed, and the structural engineer accepted, a
two-segment rod: one segment from the ceiling to the fourth-floor
walkway, a second segment from the fourth-floor walkway to the
second-floor walkway. The change appears, on a drawing, to be local. A
single rod becomes two. The connection's bookkeeping changes silently.

The FBD of the upper box-beam connection in the modified design has a
different load path. The second-floor walkway's weight now reaches the
ceiling by pulling down on the lower segment of rod, which pulls down
on the upper walkway's box-beam, which is held up by the upper
segment. The upper box-beam connection now carries both walkways'
loads. The NBS investigation calculated that the modified connection
had a factor of safety of approximately $0.6$ against the design
code: the connection was overloaded under the design load alone, and
the added tea-dance crowd produced loads that exceeded the
connection's strength by a substantial margin
\cite{acc:nbs-hyatt-1982}.

The structural engineer of record stamped the modified shop drawings
without recalculating the connection's capacity. The FBD of the upper
connection in the modified configuration was never drawn. The change
crossed the boundary from designer to fabricator; the design-load
recalculation that the change required did not. The connection failed
by tear-out of the box-beam's lower flange; the upper walkway dropped
onto the lower; the lower's connection then failed by overload; both
walkways fell. The structural engineer's licence was subsequently
revoked. The case became a canonical example in engineering ethics
curricula of the limits of the engineer's stamp as a verification
mechanism.

The mechanical lesson is the chapter's. A free-body diagram is not a
document drawn once for a design and filed. A free-body diagram is a
recalculation triggered by every change in geometry, load, or
connection. The FBD that exists in the original drawing answers a
question about the original design. The FBD that the fabricator's
shop drawing requires is a new diagram, with a new load path, and a
new equilibrium equation. The engineer who skips the redrawing has
assumed that the change is local; the diagram does not enforce that
assumption; the laws of mechanics do not honour it.

We will return to Hyatt Regency in chapter 8 of this volume (its
box-beam-connection mechanics), in chapter 8 of Volume V (the
tear-out failure mode), in Volume X (the organisational mechanism),
and in chapter 11 of Volume XI (engineering ethics under sign-off
authority). The connection deserves a multi-volume treatment because
its failure is multi-scale: a fabrication-time decision, a missing
recalculation, a code-compliance gap, an ethics failure of stamp
authority, and a real building with real people in it. The chapter's
narrow lesson is that the FBD is not a document. The FBD is a
discipline of redrawing.

\begin{mastery}
\textbf{Redundant load paths: when indeterminacy is the safe choice.} The
indeterminate crane of the jib-crane study was presented as a problem the
free-body diagram cannot close, and it is, but indeterminacy is also a design
virtue, and the Hyatt connection is the case that shows why. A determinate
structure has exactly as many reactions as equilibrium equations, so every
member is load-bearing and the loss of any one turns the structure into a
mechanism: there is no second path for the load to take. The single hanger rod
at each Hyatt connection was determinate in exactly this sense, a lone tension
path from walkway to ceiling, and when the modified connection at the top
overloaded, the load had nowhere else to go and both walkways fell. A
\engterm{redundant} structure carries more reactions than equilibrium needs, so
a load has more than one path to ground, and the failure of one path sheds its
share onto the others rather than dropping the load. The cost is that the FBD
alone no longer finds the force distribution, since the shares depend on the
relative stiffnesses of the competing paths, the compatibility condition of
chapter 3. The benefit is that a single failure is arrested rather than
propagated, the property aviation calls \engterm{fail-safe}: a structure sized
so that any one member can fail and the remainder still carry the load, at least
long enough for the failure to be found. The alternative philosophy,
\engterm{safe-life}, keeps the determinate simplicity and instead retires the
part before its statistical failure life is reached, trusting inspection and a
service limit rather than a second path. The two philosophies trade the same
quantity: safe-life buys analysability and low weight and pays with the
consequence of any missed crack, while fail-safe buys tolerance to a missed
crack and pays with the weight and the indeterminate analysis of the extra path.
A redundant Hyatt connection, two independent rods or a rod plus a bearing seat,
would have made the fabrication change survivable, because the load that
overwhelmed the single path would have found the second. The rigid-body FBD of
this chapter is the tool for the determinate structure and the alarm bell,
through the count $r > e$, for the redundant one, and the choice between the two
counts is a safety decision made before any diagram is drawn: how many paths
should the load have to ground, and what happens the day one of them is gone.
\end{mastery}

\section{Project}\pathtag{core}

\begin{project}[Free-body diagrams of ten mechanical scenes]
The reader photographs ten mechanical scenes in their immediate
environment over one week, and produces a free-body diagram of each.
The deliverable is a thirty- to fifty-page workbook.

\textbf{Track}. Physical observation plus analysis. No physical
build; no work above hand-tool scale. \textbf{Hazard class}: none.

\textbf{Scenes}. The reader photographs ten mechanical situations.
Suggested first-day list: a person carrying a bag in one hand
(asymmetric load); a bicycle leaning against a wall; a hanging
plant on a hook; a kitchen scale with a heavy pot on it; a doorstop
wedged under a door; a shelf bracket holding a row of books; a
person walking up a flight of stairs; a parked car on a sloped
street; a person on a chair leaning back on two legs; a clothes line
with a wet towel. The reader chooses ten total; at least three must
involve a non-trivial constraint (pin, roller, or rough contact at
the friction limit), and at least one must involve a rope or cable.

\textbf{Procedure}. For each scene:

\begin{itemize}[leftmargin=*]
  \item A photograph (or a clear sketch from a photograph).
  \item An identification of the chosen body for the FBD.
  \item The FBD itself, drawn clearly, with every force labelled with
    magnitude (estimated), direction, and point of application.
  \item The equilibrium equations the diagram implies.
  \item A solution, with the estimated forces propagated through the
    equations to a numerical answer for the unknowns. Estimation is
    expected; calibrated values from a kitchen scale or bathroom
    scale are encouraged.
  \item A failure-mode sentence: how would this body fail if one of
    the constraints were removed or one of the forces increased
    significantly?
\end{itemize}

\textbf{Deliverable}. The ten diagrams plus the analysis, bound or
PDF, with the photographs alongside each FBD. The general editor's
reference workbook for this project will be published online; the
reader is encouraged to compare \emph{after} completing their own.

\textbf{Time}. Twenty to forty minutes per scene, plus two to three
hours of writing and bookkeeping. The project is best done over one
week rather than one sitting; the daily photograph and diagram is
the habit the project exists to install.
\end{project}

\section{Exercises}

The exercises mix calculation, derivation, estimation, simulation,
design, diagnosis, failure analysis, and judgment per the editorial
settlement on exercise types. Solutions are provided for calculation
and derivation exercises; estimation exercises receive published
reference estimates with the editor's working; design, diagnosis,
failure analysis, and judgment exercises receive rubrics rather than
single answers. Exercises are tagged with their reader-path tier.

\subsection*{Calculation}\pathtag{core}

\begin{exercise}
A $20\,\si{\kilogram}$ box rests on a horizontal floor with $\mu_s =
0.4$. A person applies a horizontal force of $50\,\text{N}$. State
the value of the friction force on the box and whether the box
moves. Repeat for an applied force of $100\,\text{N}$.
\paragraph{Solution.} The normal force is $N = mg = 20 \times 9.81 =
196.2\,\text{N}$, so the static-friction ceiling is $\mu_s N = 0.4
\times 196.2 = 78.5\,\text{N}$. At an applied $50\,\text{N}$: this is
below the ceiling, so the box does not move and static friction matches
the applied force, $F_f = 50\,\text{N}$. At an applied $100\,\text{N}$:
this exceeds the $78.5\,\text{N}$ ceiling, so the box slides; friction
then takes its kinetic value, not $100\,\text{N}$. The trap is to write
$F_f = \mu_s N$ in the first case; below threshold, friction is whatever
balances the applied force.
\end{exercise}

\begin{exercise}
A uniform plank of mass $8\,\si{\kilogram}$ and length
$3\,\si{\meter}$ rests on two supports, one at each end. A person of
mass $70\,\si{\kilogram}$ stands $1\,\si{\meter}$ from one end. Find
the reaction at each support. State which support carries the
larger load and by what factor.
\paragraph{Solution sketch.} Call the near support (the one the person
stands closer to) the left support at $x = 0$ and the far one at $x =
3\,\si{\meter}$. The plank weight $W_p = 8 \times 9.81 = 78.5\,\text{N}$
acts at the midpoint $x = 1.5$; the person's weight $W_h = 70 \times
9.81 = 686.7\,\text{N}$ acts at $x = 1$. Moments about the left support:
$R_{\text{right}} \times 3 = W_p \times 1.5 + W_h \times 1$, giving
$R_{\text{right}} = (117.7 + 686.7)/3 = 268\,\text{N}$. Then
$R_{\text{left}} = W_p + W_h - R_{\text{right}} = 765.2 - 268 =
497\,\text{N}$. The near support carries the larger load by a factor of
$497/268 \approx 1.85$. The person's position, not the plank's weight,
sets the split.
\end{exercise}

\begin{exercise}
A point mass of $50\,\si{\kilogram}$ hangs in static equilibrium from
two ropes attached to the ceiling, making angles of $30^{\circ}$ and
$60^{\circ}$ with the vertical. Find the tension in each rope.
\paragraph{Solution.} With $\alpha = 30^{\circ}$ and $\beta =
60^{\circ}$, the two angles sum to $90^{\circ}$, so $\sin(\alpha+\beta)
= 1$ and the formulae from section 1.6.3 simplify. Here $mg = 50 \times
9.81 = 490.5\,\text{N}$. The rope at $30^{\circ}$ from vertical carries
$T_{30} = mg\sin 60^{\circ} = 490.5 \times 0.866 = 424.8\,\text{N}$; the
rope at $60^{\circ}$ carries $T_{60} = mg\sin 30^{\circ} = 490.5 \times
0.5 = 245.2\,\text{N}$. The rope nearer the vertical takes the larger
share, as it should: it is doing more of the lifting and less of the
sideways pulling. The two tensions check against the data table
\texttt{data/two\_rope\_tensions.csv}.
\end{exercise}

\begin{exercise}
A traffic light of mass $25\,\si{\kilogram}$ hangs at the midpoint of
a horizontal cable strung between two poles $20\,\si{\meter}$ apart.
The cable sags $0.5\,\si{\meter}$ at the midpoint. Find the tension
in the cable just to the left of the light's attachment point.
Comment on the implication for cable rating.
\paragraph{Solution.} The cable makes a small angle $\phi$ below
horizontal on each side, with $\tan\phi = \text{sag} / (\text{half-
span}) = 0.5 / 10 = 0.05$, so $\phi = 2.86^{\circ}$. The light's weight
$W = 25 \times 9.81 = 245.2\,\text{N}$ is shared by the vertical
components of the two cable tensions: $2T\sin\phi = W$, hence $T = W /
(2\sin\phi) = 245.2 / (2 \times 0.0499) = 2{,}456\,\text{N}$. The
tension is about ten times the weight it supports, the direct
consequence of the shallow sag. Halving the sag would roughly double the
tension. The rating implication: a cable run nearly level pays for its
flatness in tension, and the designer trades visible sag against cable
size and pole bending moment.
\end{exercise}

\begin{exercise}
A wooden block of mass $5\,\si{\kilogram}$ sits on a ramp inclined
at $20^{\circ}$ above horizontal. The static friction coefficient is
$\mu_s = 0.35$. Show that the block does not slide. Find the
friction force on the block.
\paragraph{Solution.} The slip criterion compares the demanded friction
to the available friction. Resolving along the incline, the demanded
friction is the down-slope weight component $F_f = mg\sin\theta = 5
\times 9.81 \times \sin 20^{\circ} = 16.8\,\text{N}$. The available
friction is $\mu_s N = \mu_s mg\cos\theta = 0.35 \times 5 \times 9.81
\times \cos 20^{\circ} = 16.1\,\text{N}$. Equivalently, sliding is
impending when $\tan\theta = \mu_s$; here $\tan 20^{\circ} = 0.364$ and
$\mu_s = 0.35$, so the block sits right at the slip threshold,
with the demand a hair above the limit. The honest reading is that this
geometry is marginal: a coefficient of $0.36$ or better holds it, $0.35$
leaves it on the verge. If one takes the block as held, the friction
force is the demanded $16.8\,\text{N}$, not $\mu_s N$, because below
threshold friction supplies exactly the balance the tangential equation
needs. The example is a reminder that $F_f = mg\sin\theta$ until slip,
and that a margin of a few percent on $\mu_s$ is not a margin a designer
should rely on.
\end{exercise}

\begin{exercise}
A $1{,}500\,\si{\kilogram}$ car is parked on a $5^{\circ}$ slope.
Find the normal force at the front and rear axles, given that the
car's centre of mass is $1.2\,\si{\meter}$ behind the front axle and
the wheelbase is $2.6\,\si{\meter}$.
\paragraph{Solution.} The total weight is $W = 1500 \times 9.81 =
14{,}715\,\text{N}$, of which the component pressing into the road is
$W\cos 5^{\circ} = 14{,}659\,\text{N}$ and the component down the slope
is $W\sin 5^{\circ} = 1{,}283\,\text{N}$ (held by the brakes). With the
centre-of-mass height not given, take it at road level so the moment of
the slope component about an axle drops out. Moments about the front
axle: $N_r \times 2.6 = W\cos 5^{\circ} \times 1.2$, so $N_r =
14{,}659 \times 1.2 / 2.6 = 6{,}766\,\text{N}$. Then $N_f = W\cos
5^{\circ} - N_r = 14{,}659 - 6{,}766 = 7{,}893\,\text{N}$. The front
axle carries more because the centre of mass sits closer to it ($1.2$ of
the $2.6\,\si{\meter}$ wheelbase). Were the CM height included, the
$h\sin\theta$ term would shift a little load to the downhill axle.
\end{exercise}

\begin{exercise}
A truss joint connects four members at angles $0^{\circ}$,
$90^{\circ}$, $180^{\circ}$, and $270^{\circ}$ (a cross). The joint
is loaded by an external force of $5\,\text{kN}$ at $45^{\circ}$
above horizontal. Find the forces in each member, assuming tension
positive.
\paragraph{Solution sketch.} Resolve the load: $P_x = P_y = 5\cos
45^{\circ} = 3.54\,\text{kN}$. Label the members $F_0$ (along $+x$),
$F_{90}$ ($+y$), $F_{180}$ ($-x$), $F_{270}$ ($-y$), tension positive
(force on the joint directed away along each member). The two joint
equations are $\sum F_x: F_0 - F_{180} + P_x = 0$ and $\sum F_y:
F_{90} - F_{270} + P_y = 0$. Four unknowns, two equations: the cross
joint is internally indeterminate, because the two collinear members on
each axis cannot be separated by equilibrium alone. Only the axial
\emph{net} on each line is fixed: the horizontal pair must together
supply $-P_x = -3.54\,\text{kN}$ and the vertical pair $-P_y =
-3.54\,\text{kN}$. If the back members ($F_{180}, F_{270}$) are absent
or carry zero (a common idealisation), then $F_0 = -3.54\,\text{kN}$
(compression) and $F_{90} = -3.54\,\text{kN}$ (compression). The lesson
is the determinacy count of section 1.4: a single joint resolves only
two unknowns, so a four-member joint must be closed by the neighbouring
joints, exactly the stepping the method of joints performs.
\end{exercise}

\begin{exercise}
A pulley redirects a rope from horizontal to vertical. A
$30\,\si{\kilogram}$ mass hangs from the vertical segment. Assuming
the pulley is frictionless, find the tension in the horizontal
segment and the magnitude of the resultant force on the pulley's
axle.
\paragraph{Solution.} The vertical segment carries the weight $T = mg =
30 \times 9.81 = 294.3\,\text{N}$. A frictionless pulley does not change
the tension, so the horizontal segment carries the same $294.3\,\text{N}$.
The axle must balance the vector sum of the two rope pulls, which are
perpendicular and equal: $R = \sqrt{T^2 + T^2} = T\sqrt{2} = 294.3
\times 1.414 = 416\,\text{N}$, directed along the bisector of the
right-angle turn. The axle load exceeds either rope tension by the
factor $\sqrt{2}$, the worked result of section 1.6.5.
\end{exercise}

\subsection*{Derivation}\pathtag{standard}

\begin{exercise}
Prove that the equilibrium conditions $\sum F_x = 0$, $\sum F_y = 0$,
$\sum M_O = 0$ for a planar rigid body are independent of the choice
of moment-summing point $O$. State the analogous result in three
dimensions.
\paragraph{Solution sketch.} Let the forces $\vec{F}_i$ act at points
$\vec{r}_i$. The moment about a point $O'$ relates to that about $O$ by
the transfer rule $\vec{M}_{O'} = \vec{M}_O + (\vec{r}_O - \vec{r}_{O'})
\times \vec{R}$, where $\vec{R} = \sum_i \vec{F}_i$ (derived in the
mastery box of section 1.1). Suppose equilibrium holds about $O$: $\vec{R}
= \vec{0}$ and $\vec{M}_O = \vec{0}$. Then the correction term
$(\vec{r}_O - \vec{r}_{O'}) \times \vec{R} = \vec{0}$, so $\vec{M}_{O'}
= \vec{0}$ for every $O'$. Hence the choice of moment point is free once
the force resultant vanishes. In two dimensions this means the three
scalar equations may be replaced by $\sum F_x = 0$ with two moment
equations about distinct points, or by three moment equations about
three non-collinear points. In three dimensions the result is identical
in vector form: $\sum \vec{F} = \vec{0}$ and $\sum \vec{M}_O = \vec{0}$
about any one point imply zero net moment about every point, giving six
scalar equations whose moment-point choice is the engineer's.
\end{exercise}

\begin{exercise}
Derive the formula for the tensions in two ropes supporting a hanging
mass at angles $\alpha$ and $\beta$ with the vertical (the formula
worked in section 1.6.3). Identify the limit in which the formula
predicts a tension larger than the weight, and the geometric meaning
of that limit.
\paragraph{Solution sketch.} Place the mass at the origin with rope $1$
up and to the left at angle $\alpha$ from vertical, rope $2$ up and to
the right at $\beta$. The horizontal balance is $T_1\sin\alpha =
T_2\sin\beta$ and the vertical balance is $T_1\cos\alpha + T_2\cos\beta
= mg$. From the first, $T_2 = T_1\sin\alpha/\sin\beta$; substitute into
the second: $T_1(\cos\alpha + \sin\alpha\cos\beta/\sin\beta) = mg$, i.e.
$T_1(\cos\alpha\sin\beta + \sin\alpha\cos\beta)/\sin\beta = mg$. The
numerator is $\sin(\alpha+\beta)$, so $T_1 = mg\sin\beta/\sin(\alpha+
\beta)$ and symmetrically $T_2 = mg\sin\alpha/\sin(\alpha+\beta)$. A
tension exceeds $mg$ when its numerator-to-denominator ratio exceeds one,
which happens as $\sin(\alpha+\beta) \to 0$, that is $\alpha + \beta \to
\pi$. Geometrically that is the ropes becoming collinear and horizontal:
the weight hangs from a nearly straight line, and the only way a nearly
straight rope can carry a transverse load is with very large axial
tension. The clothesline and the shallow rigging sling are the same
limit.
\end{exercise}

\begin{exercise}
Show that, for a uniform ladder against a smooth wall with rough
floor at angle $\theta$ to the horizontal, the minimum coefficient of
static friction at the floor for the ladder not to slip is $\mu_s =
1/(2 \tan\theta)$.
\paragraph{Solution sketch.} Forces on the ladder: weight $mg$ at the
midpoint, wall normal $N_W$ (horizontal, the wall is smooth), floor
normal $N_F$ (vertical), floor friction $F_f$ (horizontal, toward the
wall). Equilibrium: $\sum F_x: F_f = N_W$; $\sum F_y: N_F = mg$. Take
moments about the foot to eliminate $N_F$ and $F_f$: the weight acts at
horizontal lever arm $\tfrac{L}{2}\cos\theta$ and the wall normal at
vertical lever arm $L\sin\theta$, so $N_W L\sin\theta = mg\,\tfrac{L}{2}
\cos\theta$, giving $N_W = mg\cos\theta/(2\sin\theta) = mg/(2\tan\theta)$.
Then $F_f = N_W = mg/(2\tan\theta)$. Slipping is prevented while $F_f
\le \mu_s N_F = \mu_s mg$, that is $mg/(2\tan\theta) \le \mu_s mg$, which
rearranges to $\mu_s \ge 1/(2\tan\theta)$. The minimum coefficient is
the equality $\mu_s = 1/(2\tan\theta)$, independent of mass and length,
as the cancellation of $mg$ and $L$ shows.
\end{exercise}

\begin{exercise}
A weightless rigid link is pinned at both ends and carries no
transverse load. Prove that the force in the link is along its axis.
State the additional assumption needed for the equivalent statement
when the link has weight.
\paragraph{Solution sketch.} The link touches the world only at its two
pins, $A$ and $B$. A frictionless pin transmits force but no moment, so
the only external forces are $\vec{F}_A$ at $A$ and $\vec{F}_B$ at $B$
(no distributed load, no weight). Force balance gives $\vec{F}_A +
\vec{F}_B = \vec{0}$, so $\vec{F}_B = -\vec{F}_A$: the two forces are
equal and opposite. Moment balance about $A$ gives $\vec{r}_{AB} \times
\vec{F}_B = \vec{0}$, so $\vec{F}_B$ is parallel (or anti-parallel) to
$\vec{r}_{AB}$, the line $AB$. Hence both forces lie along the link's
axis: a two-force member carries pure axial force, tension or
compression. If the link has weight, the weight is a third external
force acting at the centre of mass, and the conclusion fails in general.
The additional assumption that restores it is that the link's weight is
negligible compared to the pin forces, the standard truss idealisation;
otherwise the member carries bending as well as axial force and must be
analysed as a loaded beam.
\end{exercise}

\begin{exercise}[\textit{deferred to Chapter 5}]
Derive the principle of virtual work for a planar rigid body in
static equilibrium. Show that for any virtual displacement
consistent with the constraints, the virtual work of all forces
vanishes. Chapter 5 of this volume develops the energy methods that
generalise this result.
\paragraph{Solution sketch.} A rigid-body virtual displacement in the
plane is a virtual translation $\delta\vec{u}$ of a reference point plus
a virtual rotation $\delta\phi$ about it; the virtual displacement of a
point at $\vec{r}_i$ is $\delta\vec{u} + \delta\phi\,\hat{k} \times
\vec{r}_i$. The virtual work of all forces is $\delta W = \sum_i
\vec{F}_i \cdot (\delta\vec{u} + \delta\phi\,\hat{k} \times \vec{r}_i)
= \left(\sum_i \vec{F}_i\right) \cdot \delta\vec{u} + \delta\phi\,
\hat{k} \cdot \left(\sum_i \vec{r}_i \times \vec{F}_i\right)$, where the
scalar-triple-product rearrangement moves the cross product onto the
forces. The two bracketed sums are the resultant force $\vec{R}$ and the
resultant moment $\vec{M}_O$. In equilibrium both vanish, so $\delta W =
0$ for every admissible $(\delta\vec{u}, \delta\phi)$. The converse
holds too: if $\delta W = 0$ for all admissible virtual displacements,
then $\vec{R} = \vec{0}$ and $\vec{M}_O = \vec{0}$, because
$\delta\vec{u}$ and $\delta\phi$ are independent and arbitrary. Chapter 5
generalises this to systems with internal degrees of freedom, where
virtual work eliminates the constraint reactions and yields equilibrium
in the active forces alone.
\end{exercise}

\subsection*{Estimation}\pathtag{core}

\begin{exercise}
Estimate the tension in the upper segment of a rope on which a
single shoe hangs from a tree branch by its laces. State assumptions
about the shoe's mass and the rope's geometry; bracket the answer
to within a factor of three.
\paragraph{Reference estimate.} Take a shoe at $0.4\,\si{\kilogram}$, so
weight $\approx 4\,\text{N}$. If the laces hang nearly straight down,
the tension is just the weight, $\approx 4\,\text{N}$. The interesting
case is a shoe thrown over a branch so both lace ends drape and the shoe
hangs at the bottom of a shallow V: the two lace segments share the
weight, $2T\cos\psi = mg$ where $\psi$ is each segment's angle from
vertical. For a typical drape, $\psi \approx 30^{\circ}$ to
$60^{\circ}$, giving $T$ between $mg/(2\cos 30^{\circ}) \approx
2.3\,\text{N}$ and $mg/(2\cos 60^{\circ}) \approx 4\,\text{N}$. A bound
of $2$ to $6\,\text{N}$ brackets the answer comfortably within a factor
of three. The dominant uncertainty is the shoe mass (a boot is several
times a sandal), not the geometry.
\end{exercise}

\begin{exercise}
Estimate the force a person of mass $80\,\si{\kilogram}$ exerts on
each handle of a wheelbarrow carrying a load of $40\,\si{\kilogram}$
of soil, distributed roughly uniformly in the wheelbarrow. State
assumptions about the wheelbarrow's geometry.
\paragraph{Reference estimate.} The person's mass is a distractor; what
matters is the lever geometry of the barrow. Model the barrow as a body
pivoting at the wheel axle, with the soil's centre of mass a distance
$d_1$ ahead (toward the wheel) and the hands a distance $L$ behind the
axle. A common layout puts the load roughly a third of the way from
wheel to hands, $d_1 \approx L/3$. Moments about the axle: $F_{\text{
hands}} L = W d_1$, with $W = 40 \times 9.81 = 392\,\text{N}$, gives
$F_{\text{hands}} \approx W/3 \approx 131\,\text{N}$ total, or about
$65\,\text{N}$ per handle. This is the lift, perhaps $13\,\si{\kilogram}$
-equivalent per hand, manageable, which is the point of putting the
wheel near the load. If the load sat at mid-span ($d_1 = L/2$), each
hand would carry $\approx 98\,\text{N}$; pushing the load forward over
the wheel is what makes the barrow effective. The empty barrow's own
weight, ignored here, adds perhaps $20$ to $40\,\text{N}$ per hand.
\end{exercise}

\begin{exercise}
Estimate the tension in the suspension cables of a footbridge of
span $30\,\si{\meter}$ and deck mass per unit length
$200\,\si{\kilogram\per\meter}$, with two cables and a sag-to-span
ratio of $1{:}10$. Comment on whether your estimate is consistent
with everyday cable sizes.
\paragraph{Reference estimate.} Deck weight per unit length $w = 200
\times 9.81 = 1{,}962\,\text{N/m}$, total $W = wL = 58.9\,\text{kN}$
over the $30\,\si{\meter}$ span. For a parabolic cable carrying a uniform
load, the horizontal tension component is $H = wL^2/(8s)$ with sag $s =
L/10 = 3\,\si{\meter}$, giving $H \approx 73.6\,\text{kN}$ for both
cables together. The maximum tension is at the supports, $T_{\max} =
\sqrt{H^2 + (W/2)^2} \approx \sqrt{73.6^2 + 29.4^2} \approx
79.2\,\text{kN}$ total, or about $40\,\text{kN}$ per cable. A
$40\,\text{kN}$ working tension calls for a steel cable of roughly $20$
to $25\,\si{\milli\meter}$ diameter at a sensible factor of safety,
which is an entirely ordinary footbridge cable size. The estimate is
consistent with everyday hardware, which is the sanity check the question
asks for.
\end{exercise}

\begin{exercise}
Estimate the friction force on the soles of a person standing
stationary on a smooth wooden floor. Then estimate the force when
the person is walking at moderate pace. State whether the static
or kinetic coefficient is relevant in each case, and why.
\paragraph{Reference estimate.} Standing still on level ground, there is
no horizontal applied force, so the friction force is zero: static
friction supplies exactly what equilibrium demands, and equilibrium
demands nothing horizontal. The static coefficient is the relevant one,
but its value is irrelevant to the magnitude (which is zero); it matters
only as the ceiling that is not being approached. Walking is different.
Each step pushes the floor backward to propel the body forward; the peak
horizontal ground-reaction force on a level walk is roughly $0.15$ to
$0.25$ of body weight, so for an $80\,\si{\kilogram}$ person, $W = 785\,
\text{N}$, the friction is $\approx 120$ to $200\,\text{N}$. This is
still static friction, because the foot does not slide relative to the
floor during a normal step; the sole is momentarily at rest on the
ground. Kinetic friction enters only when the foot slips, the event the
whole gait is arranged to avoid. The relevant ceiling is $\mu_s W$; a
slip occurs when the demanded $0.15$ to $0.25\,W$ exceeds $\mu_s W$,
which is why a $\mu_s$ below about $0.3$ (ice, polished floor) makes
ordinary walking treacherous.
\end{exercise}

\begin{exercise}
Estimate the maximum mass of a load that a person can hold at arm's
length, horizontally, without dropping. State the constraint
(muscle, joint, balance, fatigue) you believe is dominant, and
justify in one sentence.
\paragraph{Reference estimate.} Holding a weight at arm's length is a
moment problem about the shoulder. The deltoid acts on a very short
lever arm (a few centimetres) while the load acts on the full arm length
(roughly $0.6\,\si{\meter}$), so the muscle must generate a force perhaps
ten to twenty times the load weight just to hold position. For a held
mass of $10\,\si{\kilogram}$ ($\approx 98\,\text{N}$) the shoulder
muscle force is on the order of $1$ to $2\,\text{kN}$. A trained adult
can hold maybe $10$ to $20\,\si{\kilogram}$ at full horizontal extension
briefly; sustained, far less. The dominant constraint is muscle moment
at the shoulder, because the load's long lever arm against the muscle's
short one multiplies the required muscle force into the range where the
deltoid saturates well before the elbow or wrist joints reach their load
limits. Balance is secondary for a load this size, and fatigue sets the
\emph{sustained} limit below the brief one.
\end{exercise}

\subsection*{Simulation}\pathtag{mastery}

\begin{exercise}
In a language of the reader's choice, simulate the motion of a block
on a rough inclined plane as the angle of inclination increases
slowly from zero. Plot the friction force on the block as a function
of angle, and show the kink at the slipping threshold. State the
threshold angle and compare to the static-friction prediction.
\paragraph{Solution sketch.} The reference implementation is
\texttt{code/block\_on\_incline.py}. The logic at each angle $\theta$:
compute the demanded down-slope force $mg\sin\theta$ and the available
limit $\mu_s mg\cos\theta$; while demand $\le$ limit, the friction equals
the demand and the block is static; once demand exceeds limit, the
friction saturates at the kinetic value $\mu_k mg\cos\theta$ and the
block slides. The plot of $F_f$ against $\theta$ rises along the
$mg\sin\theta$ curve, then kinks down at the threshold. The threshold is
$\theta^\ast = \arctan\mu_s$, derived by setting $mg\sin\theta = \mu_s
mg\cos\theta$. For $\mu_s = 0.5$ the script reports $\theta^\ast =
26.57^{\circ}$, matching $\arctan 0.5$ exactly. A strong submission also
notes that the post-slip friction drops because $\mu_k < \mu_s$, so the
block accelerates immediately after breaking free rather than creeping.
\end{exercise}

\begin{exercise}
Implement a numerical FBD solver for a planar truss with up to ten
joints. Input: joint coordinates, member connectivity, applied loads,
support types. Output: member forces and support reactions. Test on
the truss of exercise 1.7.
\paragraph{Solution sketch.} The reference implementation is
\texttt{code/truss\_solver.py}. The method of joints becomes a single
linear system $A\vec{x} = \vec{b}$: each joint contributes two scalar
equations ($\sum F_x = 0$, $\sum F_y = 0$), and the unknown vector
$\vec{x}$ holds the member axial forces followed by the support
reactions. A member from joint $a$ to joint $b$ contributes, in joint
$a$'s rows, the unit vector $\hat{u}_{ab}$ (tension pulls the joint
toward $b$), and the negative of it in joint $b$'s rows. Support columns
place a $1$ in the row of the constrained component. Applied loads go to
$\vec{b}$ with a sign flip. The system is square and solvable when the
truss is statically determinate, $2n_{\text{joints}} = n_{\text{members}}
+ n_{\text{reactions}}$; the demo triangle truss returns member forces
$+5$, $-7.07$, $-7.07\,\text{kN}$ (one tension, two compression) under a
$10\,\text{kN}$ downward load, with pin and roller reactions of $5\,
\text{kN}$ each. A negative result means the member is in compression,
the sign convention worth checking against a hand calculation on the
exercise-1.7 cross before trusting the solver on larger inputs. A
careful implementation also checks the determinacy count and reports
singularity (a mechanism) rather than returning garbage.
\end{exercise}

\subsection*{Design}\pathtag{standard}

\begin{exercise}
Design a wall-mounted shelf to carry a row of hardcover books of
total mass $30\,\si{\kilogram}$, with a $300\,\si{\milli\meter}$
overhang. Specify the bracket geometry, the bolt size into a wall
stud, and the maximum bracket spacing. State the FBD of one
bracket and the equilibrium equations it satisfies.
\paragraph{Model-answer sketch.} The load case is a cantilever: each
bracket carries a downward load at the overhang and resists it through a
moment couple at the wall. A strong answer draws the bracket FBD as a
vertical load $P$ at distance $\approx 150\,\si{\milli\meter}$ (the load
centroid of a $300\,\si{\milli\meter}$ shelf), a top bolt in tension
$T$, and a bottom bearing reaction $C$ against the wall, separated by the
bracket height $h$. Equilibrium: $\sum F_y: C_v + (\text{bolt shear}) =
P$; $\sum M_{\text{bottom}}: T h = P d$, so $T = Pd/h$. With two
brackets sharing $30\,\si{\kilogram}$, each carries $\approx 147\,
\text{N}$ at $d = 0.15\,\si{\meter}$; a bracket height $h = 0.1\,
\si{\meter}$ gives a top-bolt tension $T \approx 221\,\text{N}$, well
within an M6 or larger lag screw into a stud. The deliverable should
state: bracket height and the load arm, the bolt tension from the moment
equation, a stated factor of safety (the books are static, but impact
and overload justify $\times 3$ or more), and a maximum bracket spacing
set by the shelf board's own bending (typically $400$ to $600\,
\si{\milli\meter}$). The grading rubric rewards the moment couple being
drawn correctly, the lever arm $d/h$ amplification being named, and a
safety factor being stated rather than assumed.
\end{exercise}

\begin{exercise}
Design a temporary brace to hold a $50\,\si{\kilogram}$ door upright
in a frame during installation, with the worker free to use both
hands. State the FBD of the brace and identify the failure mode
(slipping at the floor, tipping, buckling) you have designed
against.
\paragraph{Model-answer sketch.} The brace is a two-force compression
strut running from a foot on the floor to a contact on the door, at an
angle. A strong answer draws the strut FBD as an axial compression $F$
along its length, with the floor end resisted by a normal force and
friction, and the door end pushing horizontally against the door's
out-of-plane lean. The dominant failure mode for a temporary diagonal
strut on a smooth floor is \emph{slipping at the floor}: the horizontal
component of the strut force must be held by floor friction, $F\cos\phi
\le \mu_s N$, where $\phi$ is the strut angle from horizontal and $N$ the
vertical reaction. The design choices that follow: a strut angle steep
enough that the horizontal component stays under the friction limit, a
non-slip foot (rubber pad, or a cleat against a wall) to raise the
effective $\mu_s$, and a length-to-thickness ratio low enough to avoid
buckling under the compression (a Volume V topic, but worth naming). The
rubric rewards naming slipping as the governing mode for a smooth floor,
writing the friction inequality, and choosing the foot detail that
addresses it rather than hand-waving "make it sturdy."
\end{exercise}

\begin{exercise}
Design a hand-held lever to lift one corner of a $1{,}000\,
\si{\kilogram}$ pallet by $50\,\si{\milli\meter}$. State the lever
length, the fulcrum position, and the operator's hand force. State
the FBD of the lever and its equilibrium equations.
\paragraph{Model-answer sketch.} Lifting one corner carries roughly a
quarter of the pallet weight (the other three corners share the rest),
so the load at the lever tip is $\approx 0.25 \times 1000 \times 9.81
\approx 2{,}450\,\text{N}$. A class-1 lever (fulcrum between load and
hand) FBD shows the load $W$ down at the short arm $a$, the hand force
$F$ down at the long arm $b$, and the fulcrum reaction $R = W + F$ up.
Moment about the fulcrum: $F b = W a$, so the mechanical advantage is
$b/a$. To bring a $2{,}450\,\text{N}$ load down to a comfortable
$\approx 250\,\text{N}$ hand force needs $b/a \approx 10$: a
$1.1\,\si{\meter}$ bar with the fulcrum $0.1\,\si{\meter}$ from the load
end. The $50\,\si{\milli\meter}$ lift at the load corresponds to a
$\approx 500\,\si{\milli\meter}$ hand sweep, consistent with the same
ratio. The deliverable states the load fraction, the lever ratio, the
hand force, and the fulcrum reaction (which the fulcrum block must
withstand). The rubric rewards taking only the corner load, computing
the ratio from the moment equation, and noting the lift-distance trade
that the same ratio imposes.
\end{exercise}

\begin{exercise}
Design a hanging fixture to support a $25\,\si{\kilogram}$ light from
a wooden ceiling joist, with two attachment points spaced
$600\,\si{\milli\meter}$ apart. Identify the load on each attachment
and the failure mode the design eliminates.
\paragraph{Model-answer sketch.} If the light hangs at the midpoint
between two symmetric attachments, each carries half the weight, $0.5
\times 25 \times 9.81 \approx 123\,\text{N}$, plus any prying moment if
the fixture is rigid. A strong answer draws the fixture FBD: weight $mg$
down at the light's centre of mass, two upward attachment forces, and,
if the light hangs below the attachment plane, the geometry that turns a
horizontal offset into unequal attachment loads. Spacing two attachments
$600\,\si{\milli\meter}$ apart eliminates the single-point failure mode:
a single hook concentrates all $245\,\text{N}$ and any swing-induced
moment at one screw, whereas two points share the load and resist
rotation about the vertical. The deliverable identifies the per-
attachment load ($\approx 123\,\text{N}$ if centred, more if offset),
states the screw or lag size into the joist with a factor of safety, and
names the eliminated failure mode (single-fastener pull-out and uncontrol
-led swing). The rubric rewards the symmetry argument, the recognition
that two points resist moment as well as carry load, and a stated margin.
\end{exercise}

\subsection*{Diagnosis and reverse engineering}\pathtag{core}

\begin{exercise}
A wooden bookshelf bows visibly in the middle under load. The
bookshelf is supported only at the two ends. Draw the FBD of the
shelf and identify which loading configuration produces the maximum
deflection. State whether the failure mode is yielding, fracture, or
deflection-without-failure.
\paragraph{Model-answer sketch.} The shelf is a simply supported beam:
FBD shows the distributed book load down along its length and an upward
reaction at each end support. Among loadings of equal total weight, a
single point load at mid-span produces the largest deflection, larger
than the same weight spread uniformly (the mid-span deflection scales as
$PL^3/48EI$ for a central point load versus $5wL^4/384EI$ for a uniform
load, and for equal total the point case is larger by a factor of
$8/5$). A strong answer identifies the mode as deflection-without-failure
in most domestic cases: wood is stiff enough that a sagging shelf is
serviceable-limit-governed, bowing visibly long before the bending
stress reaches the fracture or yield level. The diagnosis is therefore a
stiffness ($EI$) problem, fixed by a thicker board, a stiffer material,
a centre support, or a stiffening lip, not a strength problem. The
rubric rewards naming the simply supported FBD, identifying the central
point load as worst case, and distinguishing the serviceability (sag)
limit from a strength (fracture) limit.
\end{exercise}

\begin{exercise}
A garden gate sags at its latch side. The gate is hinged at two
points along one vertical edge. Draw the FBD of the gate and
identify which of the two hinges carries the larger load. State
what failure mode would produce the observed sag.
\paragraph{Model-answer sketch.} The gate's weight acts at its centre of
mass, offset horizontally from the hinge line toward the latch side. The
FBD shows weight $mg$ down at the centroid, and two hinge reactions on
the hinge edge. The weight's offset creates a moment that the hinges
resist as a couple: the \emph{top} hinge is pulled away from the post
(tension, outward) and the \emph{bottom} hinge is pushed toward the post
(compression, inward), with horizontal forces of magnitude $mg \cdot
d/h$ where $d$ is the centroid offset and $h$ the hinge spacing. The top
hinge carries the larger \emph{total} load because it adds this
horizontal pull to its share of the vertical weight. The observed sag at
the latch side is produced by the top hinge pulling out: screw
withdrawal from the post, hinge-leaf bending, or post rot at the top
fastener lets the gate rotate about the bottom hinge, dropping the latch
corner. The rubric rewards the couple decomposition (top in tension,
bottom in compression), naming the top hinge as the more loaded, and
tying the sag to top-hinge pull-out rather than to the gate "being
heavy."
\end{exercise}

\begin{exercise}
A swing seat is suspended from two chains attached to a ceiling beam.
The beam shows a crack at the attachment point of one chain. Draw the
FBD of the beam (treated as a body in its own right) and identify
the load that the cracked attachment was carrying and the load
distribution after the crack.
\paragraph{Model-answer sketch.} Treat the beam as a simply supported
member with two downward chain loads applied at the attachment points
and upward reactions at the beam's own supports. Before the crack, each
attachment carries roughly half the swing-plus-occupant weight (more
during pumping, when dynamic loads can reach two to three times static).
The cracked attachment was carrying its share, a downward pull of order
$0.5 mg$ plus dynamic amplification. After the crack, the question is
whether the attachment still transmits load: a partial crack may still
carry reduced load while shedding the rest to the beam's bending and to
the other chain through the seat; a through crack drops that chain's
share entirely, throwing the full weight onto the remaining attachment
and converting the symmetric two-point load into an eccentric single-
point load. A strong answer notes that the redistribution roughly
doubles the surviving attachment's load and introduces a twisting/
eccentric load on the seat, and flags the dynamic (pumping) amplification
as the likely trigger for the original crack. The rubric rewards the
two-point beam FBD, the before/after load accounting, and the
recognition that one failed attachment overloads the survivor.
\end{exercise}

\begin{exercise}
A bicycle locked to a railing by a U-lock has shifted overnight so
the front wheel rests on the ground. Draw the FBD of the bicycle in
its current position, then in its locked-upright position. Identify
the constraint failure responsible for the shift.
\paragraph{Model-answer sketch.} In the locked-upright position the bike
is held by three contacts: the two tyres on the ground (normal plus
friction) and the U-lock at the railing (a constraint that can pull and
push within the lock's reach). The FBD shows weight at the bike's centre
of mass, ground reactions at the tyres, and the lock reaction at the
frame. The lock was providing the constraint that kept the bike from
tipping or rolling. After the shift, the FBD has the bike leaning, the
front wheel newly in contact, and the lock now bearing against a
different point with a different reaction line. The constraint failure is
that the lock did not fully fix the bike's position: a U-lock typically
prevents translation away from the railing but allows the frame to
rotate or slide within the lock's clearance, so the bike settled under
gravity until a new contact (the front wheel) closed the kinematic gap.
The rubric rewards drawing both FBDs, recognising that the lock is a
partial constraint (not a built-in support), and identifying the
allowed degree of freedom (rotation/slide within lock clearance) as what
let the bike move.
\end{exercise}

\subsection*{Failure analysis}\pathtag{core}

\begin{exercise}
Read the executive summary of the NBS report on the Hyatt Regency
walkway collapse \cite{acc:nbs-hyatt-1982}. Draw two FBDs of the
upper box-beam connection: one for the original single-rod design,
one for the modified two-rod design. State the connection load in
each case and the factor by which the modification increased the
load.
\paragraph{Model-answer sketch.} The two FBDs are drawn in
figure~\ref{fig:vol03ch01:hyatt-loadpath}. In the original single-rod
design, the upper box-beam connection carries only the fourth-floor
walkway load $W_4$: the second-floor walkway hangs from the rod below
the upper nut, so its load $W_2$ bypasses the upper connection and
reaches the ceiling through the continuous rod. In the modified two-rod
design, the lower rod hangs the second-floor walkway from the
fourth-floor box-beam, so the upper connection now carries $W_4 + W_2$.
With the two walkways of comparable weight, the connection load roughly
\emph{doubles}, a factor of about $2$. A strong answer states both load
paths explicitly, draws the load through the nut bearing on the box-beam
in each case, and reports that the NBS analysis found the as-built
connection capacity at roughly $0.6$ of the design-code requirement,
that is overloaded before any crowd was added. The rubric rewards the
correct load path in each FBD and the factor-of-two result.
\end{exercise}

\begin{exercise}
The NBS report identifies the modification as a fabrication-time
change to a shop drawing. Identify, in one paragraph, the step in
the FBD discipline of section 1.3 that, if applied, would have
caught the modified load path. State whether the FBD discipline alone
suffices or whether a procedural change at the engineer-fabricator
boundary is also required.
\paragraph{Model-answer sketch.} The decisive step is step 3, replace
each contact with its forces, applied to the \emph{as-built} geometry
rather than the original. Redrawing the upper connection's FBD for the
two-rod arrangement shows the lower rod bearing on the fourth-floor
box-beam, which puts $W_4 + W_2$ through the upper nut instead of $W_4$
alone. Equivalently, the principle that the FBD is a recalculation
triggered by every change in geometry or connection (section 1.3's
closing claim) is the one that was skipped. A complete answer then makes
the second point: the FBD discipline alone is necessary but not
sufficient, because the change crossed the engineer-fabricator boundary
and the engineer never redrew the diagram. The procedural fix is a
verification gate at that boundary, a requirement that any change to a
stamped connection return to the engineer for recalculation before
fabrication proceeds. The rubric rewards naming step 3 (or the
recalculation-on-change principle) and recognising that a discipline is
only as good as the process that triggers it.
\end{exercise}

\begin{exercise}
A footbridge collapses under a tour group of $30$ people on a static
crossing. The investigation finds that one of the bridge's two cable
anchorages had been replaced during a recent renovation with a
different bolt pattern. Argue, in $200$ to $300$ words, what kind
of FBD-discipline failure most likely occurred, and what verification
the renovation should have triggered. The argument does not require
the specific structural facts; it requires that the reader name the
discipline gap.
\paragraph{Model-answer sketch.} The most likely failure is the Hyatt
pattern at a different scale: a change to a load-bearing connection that
was not accompanied by a redrawn FBD and a recomputed capacity. A new
bolt pattern changes the anchorage's load path: the number of bolts, the
spacing, the edge distances, and the failure mode (bolt shear, plate
bearing, tear-out) all shift, and the anchorage's capacity can fall well
below the original even when the new pattern "looks adequate." A strong
answer argues that the renovation should have triggered a fresh FBD of
the anchorage under the design crowd load, with the cable tension
resolved into the new bolt group and each bolt's share checked against
its capacity, plus a check of the plate and edge-distance failure modes.
It should note that the bridge held under light traffic but failed under
a concentrated tour group, the signature of a capacity that was reduced
but not to zero, consistent with an under-designed replacement rather
than an outright omission. The discipline gap is the same one the
chapter names: the FBD that answers the original design's question does
not answer the renovated anchorage's question, and the change crossed a
boundary (original designer to renovation contractor) at which the
recalculation was not enforced. The verification owed was a stamped
recheck of the changed connection before reopening.
\end{exercise}

\subsection*{Judgment}\pathtag{standard}

\begin{exercise}
A senior colleague draws an FBD of a beam without including the
beam's weight, on the grounds that the load on the beam is much
larger than the beam's weight. Write a one-paragraph reply that
takes the colleague seriously, agrees where agreement is honest, and
disagrees where disagreement is honest.
\paragraph{Model-answer sketch.} A strong reply concedes the honest
core: when the applied load is one or two orders of magnitude larger
than the self-weight, neglecting self-weight changes the answer by a
percent or two, below the model's other uncertainties, and the
simplification is defensible for a quick check. It then disagrees where
disagreement is honest. First, the threshold is quantitative, not a
habit: the right move is to estimate the self-weight and confirm it is
small, not to omit it by reflex, because long-span or large-section
beams can have self-weights comparable to their service loads. Second,
self-weight is distributed while a point load is concentrated, so even a
small self-weight changes the shape of the bending-moment diagram and
can matter for deflection or for a member that is otherwise lightly
loaded. The disciplined position: include self-weight by default,
neglect it only after an explicit estimate shows it is negligible, and
record that the estimate was made. The rubric rewards genuine agreement
on the small-load case, a quantitative threshold rather than a rule of
thumb, and the distributed-versus-point distinction.
\end{exercise}

\begin{exercise}
A junior engineer asks whether the FBD of a rope passing over an
``ideal pulley'' can treat the tension on both sides as equal even
when the rope is being actively pulled by a person. Write a
one-paragraph reply that states the conditions under which the
assumption is acceptable and the conditions under which it is not.
\paragraph{Model-answer sketch.} The equal-tension assumption follows
from two idealisations, not from the pulley being "ideal" in name: the
pulley must be massless (or rotating at negligible angular acceleration)
and its bearing must be frictionless, and the rope must be inextensible
and massless. Under static or quasi-static conditions, a person pulling
steadily, those hold well and the two sides carry the same tension; the
active pull does not break the assumption, because at constant speed the
pulley's angular acceleration is zero and no net torque is needed. The
assumption fails when any idealisation is violated: a heavy or rapidly
accelerating pulley needs a net torque, so the tensions differ by
$I\alpha/r$; a pulley with real bearing friction or rope stiffness shows
a tension ratio set by a capstan-type relation (the side being pulled
toward is higher); and a rope with significant mass or stretch carries
different tension along its length. The disciplined reply: the equal-
tension model is acceptable for a light, free-running pulley under
steady or slow motion, and must be replaced by a torque balance on the
pulley when mass, acceleration, or friction is appreciable. The rubric
rewards naming the specific idealisations, recognising that steady
active pulling is still quasi-static, and identifying the failure cases.
\end{exercise}

\begin{exercise}
Identify a mechanical artifact in your immediate environment whose
load path you have never explicitly traced. State, in one paragraph,
the FBD of the artifact's chosen body, the constraints, and the
single force whose magnitude most surprises you on calculation.
\paragraph{Model-answer sketch.} The exercise is open by design; a strong
answer chooses a real object and carries the FBD through to a number. A
good submission has three marks: it isolates one body and lists every
external force on it (weight, contacts, fasteners), it names each
constraint by type (pin, roller, fixed, cable) and the reactions it
implies, and it computes at least one force and reports a genuine
surprise. Typical surprises that the calculation reliably produces: the
tension in a nearly horizontal cable or wire (much larger than the
weight it carries, the divergence of section 1.6.3); the bolt tension in
a cantilevered bracket (amplified by the load-arm-to-bracket-height
ratio); the muscle or hinge force in a body held at a long lever arm
(amplified by the short internal moment arm). The rubric rewards a real
object, a complete force list, correctly typed constraints, and a
computed force with a stated reason the magnitude is surprising. An
answer that asserts a surprise without a number does not meet the bar.
\end{exercise}

\begin{exercise}
Argue, in $200$ to $300$ words, whether the FBD discipline of section
1.3 should be required of every change to a stamped design before the
stamp is reapplied, or whether trained engineering judgement is
sufficient. Take a position. The Hyatt Regency case is admissible as
evidence on either side.
\paragraph{Model-answer sketch.} The exercise asks for a defended
position, not a single correct answer; the rubric grades the argument's
structure. A strong answer takes a side and engages the strongest
counter. The case for a mandatory redrawn FBD on every change: the Hyatt
collapse shows that a change a competent engineer judged local was not
local, and that judgement under fabrication-schedule pressure failed
exactly where it was most needed; a cheap, bounded procedure (redraw the
connection's FBD, recompute its capacity) would have caught it, and the
cost of the procedure is trivial against the consequence. The case for
trained judgement: not every change touches a load path, a blanket
requirement to redraw every FBD invites checkbox compliance and dilutes
attention across trivial changes, and the real failure at Hyatt was not
the absence of a rule but the engineer not recognising that the change
was load-bearing, which a rule does not fix. The defensible synthesis,
which a top answer reaches: require the recalculation specifically for
changes to load-bearing connections and load paths (a triaged rule, not
a blanket one), and treat the FBD redraw as the trigger that forces the
engineer to confront the changed load path rather than judge it from the
drawing. The rubric rewards a clear position, honest use of Hyatt on
both sides, and recognition that "rule versus judgement" is a false
binary resolved by a triaged trigger.
\end{exercise}
